\documentclass[11pt,oneside]{book}
\usepackage[margin=1.2in]{geometry}
\usepackage[toc,page]{appendix}
\usepackage{graphicx}
\usepackage{natbib}
\usepackage{lipsum}
\usepackage{caption}
\usepackage[T1]{fontenc}
\usepackage{titlesec, blindtext, color}
\usepackage{xcolor,tikz}
\usetikzlibrary{patterns}
\usetikzlibrary{decorations.markings}
\usetikzlibrary{decorations.pathmorphing}
\usepackage{amsmath,amssymb,amsthm,mathrsfs,amsfonts,xfrac,pifont,bbold,physics}
\usepackage[utf8]{inputenc}
\usepackage{amsthm}
\usepackage[breakable, theorems, skins]{tcolorbox}
\usepackage[colorlinks = true,
            linkcolor = red,
            urlcolor  = blue,
            citecolor = red,
            anchorcolor = red]{hyperref}
\usepackage{cleveref}

\usepackage{soul}
\usepackage{frcursive}
\usepackage{booktabs}
\usepackage{enumitem}
\usepackage{slashed}

% Below for fancy path integral D
\usepackage[mathscr]{euscript}
\DeclareSymbolFont{rsfs}{U}{rsfs}{m}{n}
\DeclareSymbolFontAlphabet{\mathscrsfs}{rsfs}

% -------------------------------------------------------------------
% Theorem Styles
% -------------------------------------------------------------------

\theoremstyle{definition} % Define theorem styles here based on the definition style (used for definitions and examples)
\newtheorem*{definition}{Definition}

\theoremstyle{plain} % Define theorem styles here based on the plain style (used for theorems, lemmas, propositions)
\newtheorem{theorem}{Theorem}[section]
\newtheorem{axiom}{Axiom}
\newtheorem{corollary}[theorem]{Corollary}
\newtheorem{lemma}[theorem]{Lemma}
\newtheorem{proposition}[theorem]{Proposition}

\theoremstyle{remark} % Define theorem styles here based on the remark style (used for remarks and notes)
\newtheorem*{notation}{Notation}
\newtheorem*{solution}{Solution}

\newtheoremstyle{underline}% name
{}        % Space above, empty = `usual value'
{}              % Space below
{}              % Body font
{}    % Indent amount (empty = no indent, \parindent = para indent)
{}              % Thm head font
{.}             % Punctuation after thm head
{1.5mm}         % Space after thm head: \newline = linebreak
{{\underline{\textit{\thmname{#1}\thmnumber{ #2}}~\thmnote{(#3)}\unskip}}}  % Thm head spec

\theoremstyle{underline}

\newtheorem{remark}[theorem]{Remark}
\newtheorem{example}[theorem]{Example}
\newtheorem{claim}[theorem]{Claim}
% -------------------------------------------------------------------
% Chapter Headings
% -------------------------------------------------------------------

%\setcounter{chapter}{-1}

\makeatletter
%\renewcommand{\@chapapp}{Lecture}
\makeatother
\definecolor{lightergray}{rgb}{0.9,0.9,0.9}

\usepackage{titlesec}
\titleformat{\section}{\large\bfseries\raggedright}{}{0em}{\colorsection}[\titlerule]
\titleformat{name=\section,numberless}{\large\scshape\bfseries\raggedright}{}{0em}{\colorsectionnonumber}[\titlerule]

\titleformat{\subsection}{\bfseries\raggedright}{}{0em}{\colorsubsection}
\titleformat{name=\subsection,numberless}{\bfseries\raggedright}{}{0em}{\colorsubsectionnonumber}

\newcommand{\colorsection}[1]{%
    \colorbox{lightergray}{\parbox{\dimexpr\textwidth-2\fboxsep}{\thesection\ \ #1}}}
\newcommand{\colorsectionnonumber}[1]{%
    \colorbox{lightergray}{\parbox{\dimexpr\textwidth-2\fboxsep}{#1}}}
    
\newcommand{\colorsubsection}[1]{%
    \colorbox{lightergray}{\parbox{\dimexpr\textwidth-2\fboxsep}{\thesubsection\ #1}}}
\newcommand{\colorsubsectionnonumber}[1]{%
    \colorbox{lightergray}{\parbox{\dimexpr\textwidth-2\fboxsep}{#1}}}
    
\definecolor{gray75}{gray}{0.75}
\newcommand{\hsp}{\hspace{20pt}}
\titleformat{\chapter}[hang]{\Huge\bfseries}{\thechapter\hsp\textcolor{gray75}{|}\hsp}{0pt}{\Huge\bfseries}

% Begin/end Shortcuts
\def\bc{\begin{corollary}}
\def\ec{\end{corollary}}
\def\bcl{\begin{claim}}
\def\ecl{\end{claim}}
\def\bd{\begin{definition}}
\def\ed{\end{definition}}
\def\ben{\begin{enumerate}}
\def\een{\end{enumerate}}
\def\be{\begin{equation}}
\def\ee{\end{equation}}
\def\bse{\begin{equation*}}
\def\ese{\end{equation*}}
\def\bex{\begin{example}}
\def\eex{\end{example}}
\def\bit{\begin{itemize}}
\def\eit{\end{itemize}}
\def\bl{\begin{lemma}}
\def\el{\end{lemma}}
\def\bnn{\begin{notation}}
\def\enn{\end{notation}}
\def\bn{\begin{note}}
\def\en{\end{note}}
\def\bp{\begin{proposition}}
\def\ep{\end{proposition}}
\def\bq{\begin{proof}}
\def\eq{\end{proof}}
\def\br{\begin{remark}}
\def\er{\end{remark}}
\def\bs{\begin{solution}}
\def\es{\end{solution}}
\def\btab{\begin{table}}
\def\etab{\end{table}}
\def\btb{\begin{tabular}}
\def\etb{\end{tabular}}
\def\bt{\begin{theorem}}
\def\et{\end{theorem}}
\def\btik{\begin{tikzpicture}}
\def\etik{\end{tikzpicture}}

% Letter Shortcuts
% Greek letters
\def\a{\alpha}
\def\b{\beta}
\def\del{\delta}
\def\g{\gamma}
\def\l{\lambda}
\def\n{\nabla}
\def\sig{\sigma}

% Mathbb
\def\b1{\mathbb{1}} % identity
\def\C{\mathbb{C}} % complex numbers
\def\F{\mathbb{F}}
\def\N{\mathbb{N}} % Natural
\def\Q{\mathbb{Q}} % Rational
\def\R{\mathbb{R}} % Reals
\def\Z{\mathbb{Z}} % Integers

% Mathcal
\def\cA{\mathcal{A}}
\def\cB{\mathcal{B}}
\def\cD{\mathcal{D}}
\def\cF{\mathcal{F}}
\def\cH{\mathcal{H}}
\def\cJ{\mathcal{J}}
\def\cK{\mathcal{K}}
\def\cL{\mathcal{L}}
\def\cM{\mathcal{M}}
\def\cN{\mathcal{N}}
\def\cO{\mathcal{O}}
\def\cP{\mathcal{P}}
\def\cT{\mathcal{T}}

% Mathfrak 
\DeclareMathOperator{\GL}{GL}
\def\gl{\mathfrak{gl}}
\DeclareMathOperator{\Ort}{O}
\def\ort{\mathfrak{o}}
\DeclareMathOperator{\SL}{SL}
\def\sl{\mathfrak{sl}}
\DeclareMathOperator{\SO}{SO}
\def\so{\mathfrak{so}}
\DeclareMathOperator{\SU}{SU}
\def\su{\mathfrak{su}}
\def\fF{\mathfrak{F}}

% Misc. 
\def\diag{\text{diag}}
\def\p{\partial}
\def\ss{\subset}
\def\se{\subseteq}
\def\sm{\setminus}
\def\ve{\varepsilon}
\def\vn{\varnothing}
\def\cl{\colon}
\def\la{\langle}
\def\ra{\rangle}
\def\lra{\overset{\leftrightarrow}}
\def\lar{\overset{\leftarrow}}
\def\pD{\mathscrsfs{D}}

% boxes 
\DeclareRobustCommand{\mybox}[2][colframe=blue!10!black,before skip=10pt]{%
\begin{tcolorbox}[
        breakable,
        width=\dimexpr\textwidth\relax, 
        boxsep=5pt,
        arc=5pt,outer arc=5pt,
        ]
        #2
\end{tcolorbox}
}
\def\bbox{\begin{tcolorbox}[colback=white!5,colframe=blue!75!black,title=Exercise]}
\def\ebox{\end{tcolorbox}}

% Tikz Stuff 
\def\midarrow{\draw[thick, decoration={markings, mark=at position 0.5 with {\arrow{>}}}, postaction={decorate}]}
\def\midarrowblue{\draw[thick, blue, decoration={markings, mark=at position 0.5 with {\arrow{>}}}, postaction={decorate}]}
\def\midarrowred{\draw[thick, red, decoration={markings, mark=at position 0.5 with {\arrow{>}}}, postaction={decorate}]}
\def\aftermidarrow{\draw[thick, decoration={markings, mark=at position 0.75 with {\arrow{>}}}, postaction={decorate}]}
\def\aftermidarrowblue{\draw[thick, blue, decoration={markings, mark=at position 0.75 with {\arrow{>}}}, postaction={decorate}]}
\def\aftermidarrowred{\draw[thick, red, decoration={markings, mark=at position 0.75 with {\arrow{>}}}, postaction={decorate}]}
\def\beforemidarrow{\draw[thick, decoration={markings, mark=at position 0.25 with {\arrow{>}}}, postaction={decorate}]}
\def\wavey{\draw[thick, decorate, decoration={snake, segment length=2.5mm, amplitude=0.75mm}]}
\newcommand{\AxisRotatorL}[1][rotate=0]{%
    \tikz [x=0.25cm,y=0.60cm,line width=.2ex,-stealth,#1] \draw (0,0) arc (-150:150:1 and 1);%
}
\newcommand{\AxisRotatorR}[1][rotate=0]{%
    \tikz [x=0.25cm,y=0.60cm,line width=.2ex,-stealth,#1] \draw (0,0) arc (150:-150:1 and 1);%
}
\def\gluon{\draw[thick,decoration={aspect=0.5, segment length=1.5mm, amplitude=1.5mm,coil},decorate]}

\begin{document}

\captionsetup[figure]{margin=1.5cm,font=small,labelfont={bf},name={Figure},labelsep=colon,textfont={it}}
\captionsetup[table]{margin=1.5cm,font=small,labelfont={bf},name={Table},labelsep=colon,textfont={it}}
\setlipsumdefault{1}

\frontmatter

\begin{titlepage}
	\centering
	    \scshape % Use small caps for all text on the title page
        \vspace*{\baselineskip} % White space at the top of the page
        
	    \rule{\textwidth}{1.6pt}\vspace*{-\baselineskip}\vspace*{2pt} % Thick horizontal rule
	    \rule{\textwidth}{0.4pt} % Thin horizontal rule
	    
	    \vspace{0.75\baselineskip} % Whitespace above the title
	    
	    {\LARGE QFT II} % Title
	    
	    \vspace{0.75\baselineskip} % Whitespace below the title
	    
	    \rule{\textwidth}{0.4pt}\vspace*{-\baselineskip}\vspace{3.2pt} % Thin horizontal rule
	    \rule{\textwidth}{1.6pt} % Thick horizontal rule
	    
        \vspace{5\baselineskip} % Whitespace after the title block

	    Course delivered in 2019 by 
	
	    \vspace{0.5\baselineskip} % Whitespace before the editors
	
	    {\scshape\Large Dr. Nabil Iqbal} % Lecturer Name
	
	    \vspace{0.5\baselineskip} % Whitespace below the editor list
	
	    \textit{Durham University} % Lecturer Institution
	    
	    \vspace{2\baselineskip} % Whitespace after the title block
	    
	    \includegraphics[width=8cm]{images/DurhamLogo.png}\\[1cm] % Logo 

	    Notes taken by 
	
	    \vspace{0.5\baselineskip} % Whitespace before my name
	
	    {\scshape\Large Richie Dadhley} % Lecturer Name
	   
	    \vspace{0.5\baselineskip} % Whitespace below my name
	    \textit{richie@dadhley.com} % Email
	
	    \vfill % Whitespace between editor names and publisher logo
\end{titlepage}



% -------------------------------------------------------------------
% Acknowledgements
% -------------------------------------------------------------------

\newpage
\section*{Acknowledgements}

These are my notes on the 2019 lecture course "QFT II" taught by Dr. Nabil Iqbal at Durham University as part of the Particles, Strings and Cosmology Msc. For reference, the course was meant to last 16 hours and be taught over 4 weeks. \\

\mybox{There was a strike during the course and 6 hours of the course were not taught. However, Dr. Iqbal gave us notes for the bits missed, so I have also included those here. This is the content of chapters 5, 6 \& 8. For this reason I have grouped the material according to Dr. Iqbal's chapters rather then the lecture numbers (as I normally do).}

I have tried to correct any typos and/or mistakes I think I have noticed over the course. I have also tried to include additional information that I think supports the taught material well, which sometimes has resulted in modifying the order the material was taught. Obviously, any mistakes made because of either of these points are entirely mine and should not reflect on the taught material in any way. \\

% For any other notable sources used:
% One of the main sources I have used to obtain the additional information is Dr. David Tong’s String Theory notes, available via the following link  

%\begin{center}
%    \textcolor{red}{Link}
%\end{center}

I would like to extend a message of thanks to Dr. Iqbal for his brilliant teaching of this course. \\

If you have any comments and/or questions please feel free to contact me via the email provided on the title page. \\

These notes are now done. A list of other notes/works I have available, visit my blog site

\begin{center}
    \href{https://richie291.wixsite.com/theoreticalphysics}{https://richie291.wixsite.com/theoreticalphysics}
\end{center}

These notes are not endorsed by Dr. Iqbal or Durham University.

\vspace{1cm}

\begin{flushright}
    \Huge{{\cursive\setul{0.1ex}{}\ul{~~Richie Dadhley~~}}}
\end{flushright}

% -------------------------------------------------------------------
% Contents
% -------------------------------------------------------------------

\tableofcontents

% -------------------------------------------------------------------
% Main sections (as required)
% -------------------------------------------------------------------

\mainmatter

\chapter{Conventions}

This course is a continuation of the QFT I course taught by Dr. Douglas Smith, and as such it is expected that the reader has taken this course/one similar. There's a couple tiny convention changes between these two courses, so here we just highlight these. 

We still work with the field theorist's metric convention of $\eta^{\mu\nu} = \text{diag}(+1,-1,-1,-1)$, and still use Greek letters to denote spacetime indices (i.e. $\mu=0,...,(d-1)$) and Latin letters to denote just spatial indices (e.g. $a=1,...,(d-1)$). Of course we will often use $4$-dimensional spacetime, however we will try to avoid doing this too often as a lot of the results will hold for higher dimensions also. 

However in contrast to the QFT I course, here we will work in \textit{natural units} and set $\hbar=c=1$. The other different convention is the definition of the source term in the partition function. In QFT I, we had that the partition function for QM was\footnote{I've slightly changed notation to use a fancy $\pD$ for the path integral measure, in contrast to the regular boring $D$ I used in QFT I. I have also included square brackets around the variable, in agreement with Dr. Iqbal's notation. I have done this purely because I think its a nice way to keep track of what the paths are. I might go back to QFT I and change this later, we'll see.}
\bse 
    Z[J] = \cN \int [\pD q] \exp\bigg(\frac{i}{\hbar} S[q] + \int dt \, q(t) J(t)\bigg),
\ese 
where the $i/\hbar$ only appears with the action term. In this course we shall redefine $J(t) \to -i\hbar J(t)$, so that (with $\hbar=1$) the partition function for QM is 
\be 
\label{eqn:QMPathIntegral}
    Z[J] = \cN \int [\pD q] \exp\bigg(iS[q] + i\int dt \, q(t) J(t)\bigg),
\ee

From this definition it follows that we also have to alter our functional derivatives relating the partition function to the $N$-point Green's functions to include the $-i$ factors. That is
\bse 
    G(t_1,...,t_N) = \bra{0}\cT[q(t_1)...q(t_N)]\ket{0} =  \frac{(-i)^N}{Z_0} \frac{\del Z[J]}{\del J(t_1) ... \del J(t_N)}\bigg|_{J=0},
\ese 
where $\cT[...]$ is the time ordering function, and $Z_0 = Z[0]$ is the normalisation factor.

\br 
    One quick remark before we start. The aim of this course is essentially to represent QFT in terms of path integrals. For this reason we will derive \textit{a lot} of relationships we have seen before from the canonical treatment (e.g. in IFT). However as much as possible it is advised that we try forget that we know the results before hand, so that we can appreciate the self contained nature of the path integral approach. I will include small remarks here and there to remind us of this.  
\er 
\chapter{Free Theories}

First let's just remind ourselves why we even bother considering QFT from the path integral language, given that we already have the canonical approach. There are two main reasons 
\ben[label=(\roman*)]
    \item The symmetries of the problem are \textit{manifest}, as the action is the main object here. 
    \item Path integrals allow us to ease over subtle points such as the question "what is the structure of the Hilbert space of our system exactly?" This is very convenient when discussing things like gauge field theories, where subtle constraints on the Hilbert space arise.
\een 

\section{Path Integral Of Fields}

Before we can actually study free QFTs in terms of the path integral, we need to adapt our QM path integral, \Cref{eqn:QMPathIntegral}, to work for fields. That is, we need to change our path integral so that we integrate over \textit{fields} $\phi$ not just position variables $q$. 

So how do we do this? Well recall that the path integral was derived using the intuitive notion of "all the different paths connecting two points". We imagined having a bunch of screens, which represented our time slices, with the points on the screen labelling the position at that time. This position label is exactly our $q(t)$. 

Now it is reasonably easy to see how we could extend this concept to two dimensions by giving $q(t)$ an index $q_a(t)$, with $a=1,2$. We would then get \textit{sheets} for our equal time slices and the two $q_a(t)$s would tell us the `coordinate' positions on the sheets. We then just let $a$ run from $1$ to some large number $N$. We then make the argument that we let $a$ run over all the different points in spacetime and we formally get the definition of a field $\phi(x)=\phi(t,\Vec{x})$. So all we have to do is replace our path integral over positions with a path integral over fields:
\bse
    Z := \int [\pD\phi] \exp\big(iS[\phi]\big),
\ese
where $S[\phi]$ is the action for our field theory. 

The rest of the definitions extend similarly, for example our $N$-point Green's functions become functions of spacetime variables, e.g. the $2$-point Green's function is 
\bse
    G(x,y) = \bra{0}\cT[\phi(x)\phi(y)]\ket{0} = \frac{1}{Z_0} \int [\pD\phi] \phi(x)\phi(y) \exp\big(S[\phi]\big).
\ese 

\section{Real Scalar Field}

Ok so let's look at a QFT then. As always we go to the simplest case, which is the free real scalar field, which has the action\footnote{Here we use the common notation of $(\p\phi)^2 := \p_{\mu}\phi\p^{\mu}\phi$, as well as $\p^2 = \p_{\mu}\p^{\mu}$.}
\be 
\label{eqn:RealScalarAction}
    \begin{split}
        S[\phi] & = \frac{1}{2}\int d^4 x \big( (\p\phi)^2 - m^2\phi\big) \\
        & = -\frac{1}{2}\int d^4x \big(\phi (\p^2+m^2)\phi\big),
    \end{split}
\ee 
where the second line is a clever rewriting\footnote{Note the minus sign out the front and the plus sign between $\p^2$ and $m^2$.} that will come in use shortly. The equations of motion\footnote{Recall these just come from the Euler-Lagrange equations. See IFT, or any similar course.}
\be 
\label{eqn:KleinGordanEquation}
    (\p^2+m^2)\phi = 0,
\ee 
known as the \textit{Klein-Gordan equation}.

\subsection{Feynman Propagator: Green's Function Of KG Equation}

\bd[Feynman Propagator]
    We define the \textit{Feynman propagator} to be the time-ordered $2$-point function
    \be 
    \label{eqn:FeynmanPropagatorDefinition}
        D_F(x,y) := \bra{0} \cT [\phi(x)\phi(y)]\ket{0}.
    \ee 
\ed 

The Feynman propagator is a very important object in QFT, and we can obtain it by introducing a source term $J(x)$ into our path integral, just as we did in QFT I:
\mybox{
    \be 
    \label{eqn:QFTPathIntegral}
        Z := \cN \int [\pD\phi] \exp\bigg(iS[\phi] + i \int dx J(x)\phi(x)\bigg).
    \ee 
}
\noindent where $\cN$ is some normalisation. From here, we can rederive the result from QFT I that related $N$-point Green's functions to functional derivatives of the partition function, e.g. the two point function is given by
\be 
\label{eqn:2PointFunctionDefinition}
    G(x,y) = \frac{(-i)^2}{Z_0} \frac{\del^2 Z[J]}{\del J(x) \del J(y)}\bigg|_{J=0}. 
\ee 
This tells us that the Feynman propagator is given by
\bse 
    D_F(x,y) = \frac{(-i)^2}{Z_0} \frac{\del^2 Z[J]}{\del J(x) \del J(y)}\bigg|_{J=0},
\ese
where the partition function is given by
\be 
\label{eqn:PartitionFunctionRealScalar}
    Z[J] = \cN \int [\pD\phi] \exp\bigg( -\frac{i}{2}\int d^4x\, \Big(\phi(\p^2+m^2)\phi\Big) + i \int dx\, J(x)\phi(x)\bigg).
\ee 
In order to get a nice result for the Feynman propagator, we introduce the following proposition. 

\bp 
    Let $\mathbf{A}$ be a real, symmetric $N\times N$ matrix with positive eigenvalues $\{a_1,...,a_N\}$, and let $\mathbf{x}$ and $\mathbf{J}$ be $N$-component column matrices. Then the following integral identity holds\footnote{Here the notation is such that an upper index tells us the row, while lower indices tell us the column.}
    \be 
    \label{eqn:ExtendedGaussianIntegral}
        \int dx_1 ... dx_N \exp\bigg( -\frac{1}{2}x_a{A^a}_bx^b + iJ_ax^a\bigg) = \sqrt{\frac{(2\pi)^N}{\det\mathbf{A}}} \exp\bigg(-\frac{1}{2} J_a {\big(A^{-1}\big)^a}_bJ^b\bigg),
    \ee 
    where $\det\mathbf{A}$ means the product of the eigenvalues of $\mathbf{A}$. 
\ep 

\bq 
    First we prove a slightly simpler result. Consider just
    \bse 
        I_1 = \int dx_1...dx_N \exp\bigg( -\frac{1}{2}x_a{A^a}_bx^b \bigg).
    \ese 
    Now because $\mathbf{A}$ is a real, symmetric $N\times N$ matrix, we can write it in terms of a diagonal matrix $\mathbf{D}$ by using some orthogonal matrix $\mathbf{R}\in O(N)$:
    \bse 
        \mathbf{A} = \mathbf{R}^T \mathbf{D}\mathbf{R},
    \ese 
    which in terms of the components reads 
    \bse 
        {A^a}_b = {R^a}_c {D^c}_d {R^d}_b.
    \ese
    Now we use the change of variables 
    \bse 
        \mathbf{y} = \mathbf{R}\mathbf{x} \qquad \implies \qquad y^a = {R^a}_b x^b
    \ese 
    which has Jacobean $|\det\mathbf{R}|$, to obtain 
    \bse
        \begin{split}
            I_1 & = \int dx_1 ... dx_N \exp\bigg(-\frac{1}{2} x_a {R^a}_c {D^c}_d {R^d}_b x^b\bigg) \\
            & = \int dy_1...dy_N |\det\mathbf{R}| \exp\bigg(-\frac{1}{2} y_c {D^c}_d y^d\bigg) \\
            & = \int dy_1 ... dy_N \exp\bigg(-\frac{1}{2} \sum_{i=1}^N a_i y_i^2\bigg) \\
            & = \prod_{i=1}^n \Bigg[\int dy_i \exp\bigg(-\frac{1}{2}a_i y_i^2\bigg)\Bigg],
        \end{split}
    \ese
    where we have used $\det\mathbf{R} = \pm 1$. Now we use the standard Gaussian integral result 
    \bse 
        \int dx \exp\big(-ax^2\big) = \sqrt{\frac{\pi}{a}},
    \ese 
    to obtain 
    \bse 
        I_1 = \prod_{i=1}^N \sqrt{\frac{2\pi}{a_i}} = \sqrt{\frac{(2\pi)^N}{\det\mathbf{A}}}.
    \ese 
    Now we just have to manipulate the integral appearing in \Cref{eqn:ExtendedGaussianIntegral} so that we can use this result. We obtain this by completing the square: 
    \bse 
        -\frac{1}{2}x_a {A^a}_b x^b + i J_ax^a = - \frac{1}{2} \Big(x_a - i{\big(A^{-1}\big)^c}_aJ_c \Big) {A^a}_b \bigg(x^b - i{\big(A^{-1}\big)^b}_dJ^d \Big) - \frac{1}{2}J_c {\big(A^{-1}\big)^c}_a J^a,
    \ese 
    so if we define 
    \bse 
        \widetilde{x}_a := \Big(x_a - i{\big(A^{-1}\big)^c}_aJ_c \Big),
    \ese 
    we get 
    \bse 
        -\frac{1}{2}x_a {A^a}_b x^b + i J_ax^a = -\frac{1}{2}\widetilde{x}_a {A^a}_b \widetilde{x}^b - \frac{1}{2}J_c {\big(A^{-1}\big)^c}_a J^a. 
    \ese 
    Now the last term is independent of the $x_i$s, so we can take it out the integral, and then the result \Cref{eqn:ExtendedGaussianIntegral} follows immediately. 
\eq 

\bbox 
    Convince yourself that the completing the square above is valid and that it does indeed give us \Cref{eqn:ExtendedGaussianIntegral}. 
\ebox 

Ok so we can use the result of this proposition to rewrite our partition function, \Cref{eqn:PartitionFunctionRealScalar}, with $A=i(\p^2+m^2)$. This gives us 
\be 
\label{eqn:PartitionDetKG}
    Z[J] = \det(\frac{\p^2+m^2}{2\pi i})^{-1/2} \cN \exp\bigg(-\frac{1}{2}\int d^4x d^4y J(x) \Delta_F(x,y) J(y), \bigg) 
\ee 
where $\Delta_F(x,y)$ is the inverse of our operator, that is 
\bse 
    i(\p^2+m^2)\Delta_F(x,y) = \del^{(4)}(x-y),
\ese 
in other words it is a Green's function of the Klein-Gordan equation, \Cref{eqn:KleinGordanEquation}. 

The above result allows us to take the functional derivatives easily, giving us\footnote{Note we have used $Z_0:=Z[0]=1$ here.}
\bse 
    D_F(x,y) := \bra{0}\cT[\phi(x)\phi(y)]\ket{0} = \Delta_F(x,y). 
\ese 
So what we have done is \textit{derived} that the Feynman propagator is the Green's function of the Klein-Gordan equation, i.e. 
\mybox{
    \be 
    \label{eqn:FeynmanPropagatorKleinGordanEquation}
        i(\p^2+m^2)D_F(x,y) = \del^{(4)}(x-y).
    \ee 
}
\br 
    Note the notation above alone obviously does not tell us that \Cref{eqn:FeynmanPropagatorKleinGordanEquation} holds, by which I mean we cannot use prior information from our IFT module to say "ah well we know that the Feynman propagator is a Green's function of the Klein-Gordan equation", but we need to explicitly show it. Indeed here I have chosen to use $D_F(x,y)$ to define the Feynman propagator, while in the IFT notes I used $\Delta_F(x,y)$ to be the defining symbol. I did this to hopefully stray us away from trying to use prior knowledge to jump to the answers. 
\er 

\br 
    Note that the determinant appearing in \Cref{eqn:PartitionDetKG} is a \textit{functional determinant}. It is the determinant of an infinite number of eigenvalues. This might sound scary and ill-defined but is actually well defined. If it scares you, fear not as we will not probe it in much detail in this course. In fact we're just going to absorb it into the normalisation factor $\cN$ and essentially forget about it. 
\er 

\br 
    As we have just seen, the Feynman propagatoris given by the inverse of the operator appearing as 
    \bse 
        -\frac{1}{2}\phi A \phi.
    \ese
    This is not something that is specific to the case considered here but is a general result for the path integral approach to QFT. That is, given the action for the QFT you can find the Feynman propagator by finding the "inverse" of the operator appearing between the quadratic terms. By inverse here we mean the function that is mapped to the identity of the target. For example if we were considering an action where the fields had index structure we would be looking for the inverse of 
    \bse 
        -\frac{1}{2} \psi^{\mu} A_{\mu\nu} \psi^{\nu}.
    \ese 
    The inverse $D_F^{\mu\nu}(x,y)$ would then need to satisfy\footnote{The $-i$ on the right-hand side just comes from the $i$ on the left-hand side of the expressions above.}
    \bse 
        A_{\mu\nu} D_F^{\nu\rho}(x,y) = -i\del^{\rho}_{\mu} \del^{(4)}(x-y).
    \ese 
    We will see an explicit calculation of this (and the problems that can arise) when looking at the quantisation of QED later.
\er 
\subsection{Calculating The Feynman Propagator}

We have seen that the Feynman propagator is a Green's function of the Klein-Gordan equation, but now we actually want to obtain an explicit expression for it. The way we do this is by considering the Fourier transform 
\be 
\label{eqn:FeynmanPropagatorFourierTransform}
    D_F(x,y) = \int \frac{d^4 p}{(2\pi)^4} e^{-ip\cdot(x-y)} \widetilde{D}_F(p),
\ee 
so that \Cref{eqn:FeynmanPropagatorKleinGordanEquation} becomes 
\bse 
    \int \frac{d^4 p}{(2\pi)^4} (-p^2+m^2) e^{-ip\cdot(x-y)} \widetilde{D}_F(p) = -i\del^{(4)}(x-y). 
\ese 
We now recall the result 
\bse 
    \int \frac{d^4p}{(2\pi)^4} e^{ip\cdot(x-y)} = \del^{(4)}(x-y). 
\ese 
which tells us 
\be 
\label{eqn:FeynmanPropagatorMomentumSpace}
    \widetilde{D}_F(p) = \frac{i}{p^2-m^2}. 
\ee 
If we then substitute \Cref{eqn:FeynmanPropagatorMomentumSpace} into \Cref{eqn:FeynmanPropagatorFourierTransform}, we get\footnote{If you can't see how to get to the second line, see the proof of Claim 5.2.1 of my IFT notes. This is on page 41.}
\bse
    \begin{split}
        D_F(x,y) & = \int \frac{dp_0}{2\pi} \int \frac{d^3\Vec{p}}{(2\pi)^3} \frac{i}{p_0^2 - \Vec{p}^2 - m^2} e^{-ip_0(x^0-y^0)+i\Vec{p}\cdot(\Vec{x}-\Vec{y})} \\
        & = i\int \frac{dp_0}{2\pi} \int \frac{d^3\Vec{p}}{(2\pi)^3} \frac{1}{p_0 + \omega_{\Vec{p}}} \frac{1}{p_0 - \omega_{\Vec{p}}} e^{-ip_0(x^0-y^0)+i\Vec{p}\cdot(\Vec{x}-\Vec{y})}
    \end{split}
\ese 
where we have defined 
\bse 
    \omega_{\Vec{p}} := \sqrt{\Vec{p}^2 +m^2}.
\ese 
So if we consider the complexification of our integral (i.e. we let $p_0\in\C$), then we see we get two poles at $p_0 = \pm \omega_{\Vec{p}}$. We then have to ask "how do we go around these poles?" 

This is a point where the path integral perhaps gives a nicer answer then the canonical approach. In the latter we just claimed that if we wanted to get the Feynman propagator then we had to take the contour so that we only picked up one pole. That is we took the following integral 
\begin{center}
    \btik 
        \draw[thick, ->] (-3.5,0) -- (3.5,0);
        \node at (4.2,0) {\large{Re($p_0$)}};
        \draw[thick, ->] (0,-1.5) -- (0,1.5);
        \node at (-0.7,1.5) {\large{Im($p_0$)}};
        \draw[fill=black] (-2,0) circle [radius=0.07cm];
        \node at (-2,0.5) {\large{$-\omega_{\vec{p}}$}};
        \draw[fill=black] (2,0) circle [radius=0.07cm];
        \node at (2,-0.5) {\large{$+\omega_{\vec{p}}$}};
        \draw[ultra thick, blue, ->] (-0.5,0) -- (-0.49,0);
        \draw[ultra thick, blue, ->] (0.5,0) -- (0.51,0);
        \draw[ultra thick, blue] (-3,0) -- (-2.5,0) .. controls (-2.25,-0.5) and (-1.75,-0.5) .. (-1.5,0)-- (1.5,0) .. controls (1.75,0.5) and (2.25,0.5) .. (2.5,0) -- (3,0);
    \etik 
\end{center}
This way we pick up one pole for each case $x^0>y^0$ and $x^0<y^0$. It was perhaps not entirely obvious why doing this would give us the Feynman propagator over, say, taking a contour which picked up both poles for $x^0>y^0$ and neither for $x^0<y^0$. 

In the path integral approach we can see exactly why. We recall from QFT I that the path integral is only really well defined when we take the Wick rotation.\footnote{We need this to check our integrals are convergent, see section 1.5 of those notes.} If we send $t\to -i\tau$, then we get the Euclidean momentum: $p_0 = ip_0^E$. So our integral $dp_0$ is to be taken along the $\Im(p_0)$ axis, and we have to take the orientation running up the the axis, i.e. we take 
\begin{center}
    \btik 
        \draw[thick, ->] (-3.5,0) -- (3.5,0);
        \node at (4.2,0) {\large{Re($p_0$)}};
        \draw[thick, ->] (0,-1.5) -- (0,1.5);
        \node at (-0.7,1.5) {\large{Im($p_0$)}};
        \draw[fill=black] (-2,0) circle [radius=0.07cm];
        \node at (-2,0.5) {\large{$-\omega_{\vec{p}}$}};
        \draw[fill=black] (2,0) circle [radius=0.07cm];
        \node at (2,-0.5) {\large{$+\omega_{\vec{p}}$}};
        \draw[ultra thick, blue] (0,-1.5) -- (0,1.5);
        \draw[ultra thick, blue, ->] (0,-.8) -- (0,-0.75);
        \draw[ultra thick, blue, ->] (0,0.7) -- (0,0.75);
    \etik 
\end{center}
This comes from the fact that the Wick rotation is defined by 
\bse 
    t \to e^{-i\theta} t, \qquad \theta \in [0,\pi/2],
\ese 
so we just rotate the contour from the real axis anticlockwise by $\pi/2$. 

Now, from $p_0 = ip_0^E$, we have
\bse 
    D_F(x,y) \sim \int dp_0 \, e^{-p_0^E(x^0-y^0)}
\ese 
and so for $x^0>y^0$ we pick up the $+\omega_{\vec{p}}$ pole, as otherwise the integral diverges at the infinite boundary. That is we have 
\bse 
    D_F(x,y) = \frac{(-2\pi i)i}{2\pi} \int \frac{d^3\vec{p}}{(2\pi)^3} \frac{1}{2\omega_{\vec{p}}} e^{-i\omega_{\vec{p}}\,(x^0-y^0) + i \vec{p}\cdot (\vec{x}-\vec{y})} \qquad \text{for} \qquad  x^0>y^0,
\ese 
where we get the minus sign in $(-2\pi i)$ because we close the contour clockwise.\footnote{That is the convention says we get $+2\pi i$ for anticlockwise contours.} Similarly we close the contour anticlockwise and pick up the $-\omega_{\vec{p}}$ for $x^0<y^0$, giving us 
\bse 
    D_F(x,y) = \frac{(2\pi i)i}{2\pi} \int \frac{d^3\vec{p}}{(2\pi)^3} \frac{1}{-2\omega_{\vec{p}}} e^{+i\omega_{\vec{p}}\,(x^0-y^0) + i \vec{p}\cdot (\vec{x}-\vec{y})} \qquad \text{for} \qquad  x^0<y^0.
\ese 
We then note that we can take the minus sign appearing in the denominator here and put it with the $(2\pi i)$ prefactor, so the only difference between the two cases is the sign appearing in the $\omega_{\vec{p}}\,(x^0-y^0)$ exponential. 

It is for this reason that we often present the momentum space Feynman propagator by shifting the two poles slightly off the real axis. That is we rewrite \Cref{eqn:FeynmanPropagatorMomentumSpace} as 
\mybox{
\be 
\label{eq:FeynmanPropagatorMomentumSpaceEpsilon}
    \widetilde{D}_F(p) = \frac{i}{p^2-m^2 + i\epsilon}
\ee 
}

\br 
    This result should look very familiar from IFT: it is the momentum space Feynman rule for internal $\phi$ line.
\er 

\section{Higher Point Functions}

So far what we have computed is the Feynman propagator which corresponds to the two point function $\bra{0}\cT[\phi(x)\phi(y)]\ket{0}$, however obviously we can look at higher point functions. The nice thing\footnote{Or perhaps boring thing...} about free theories is that all the higher point functions can be expressed in terms of the Feynman propagator. Let's see this explicitly. 

\bex 
    Let's find the 4-point function 
    \bse 
        G(x_1,x_2,x_3,x_4) = \bra{0}\cT[\phi(x_1)\phi(x_2)\phi(x_3)\phi(x_4)]\ket{0},
    \ese 
    where the subscripts just label the different variables (i.e. they're not spacetime indices). We have already seen that this is given by taking functional derivatives, i.e. 
    \bse 
        G(x_1,x_2,x_3,x_4) = (-i)^4 \frac{1}{Z_0}\frac{\del^4 Z[J]}{\del J(x_1) \del J(x_2) \del J(x_3) \del J(x_4)}\Bigg|_{J=0}.
    \ese 
    Then using \Cref{eqn:PartitionDetKG} with $\Delta_F(x,y)=D_F(x,y)$ we get 
    \bse 
        G(x_1,x_2,x_3,x_4) = (-i)^4 \Big[ D_F(x_1,x_2)D_F(x_3,x_4) + D_F(x_1,x_3)D_F(x_2,x_4) + D_F(x_1,x_4)D_F(x_2,x_4)\Big].
    \ese 
    We can obtain this pictorially by plotting four points and then connecting them in all the possible ways, i.e. 
    \begin{center}
        \btik 
            \begin{scope}[xshift=-4cm]
                \draw[fill=black] (-1,1) circle [radius=0.07cm] node [left] {$1$}; 
                \draw[fill=black] (-1,-1) circle [radius=0.07cm] node [left] {$2$};
                \draw[fill=black] (1,1) circle [radius=0.07cm] node [right] {$3$};
                \draw[fill=black] (1,-1) circle [radius=0.07cm] node [right] {$4$};
                \draw[thick] (-1,-1) -- (-1,1);
                \draw[thick] (1,1) -- (1,-1);
                \node at (2,0) {\Large{$+$}};
            \end{scope}
            \begin{scope}
                \draw[fill=black] (-1,1) circle [radius=0.07cm] node [left] {$1$}; 
                \draw[fill=black] (-1,-1) circle [radius=0.07cm] node [left] {$2$};
                \draw[fill=black] (1,1) circle [radius=0.07cm] node [right] {$3$};
                \draw[fill=black] (1,-1) circle [radius=0.07cm] node [right] {$4$};
                \draw[thick] (-1,-1) -- (1,-1);
                \draw[thick] (1,1) -- (-1,1);
                \node at (2,0) {\Large{$+$}};
            \end{scope}
            \begin{scope}[xshift=4cm]
                \draw[fill=black] (-1,1) circle [radius=0.07cm] node [left] {$1$}; 
                \draw[fill=black] (-1,-1) circle [radius=0.07cm] node [left] {$2$};
                \draw[fill=black] (1,1) circle [radius=0.07cm] node [right] {$3$};
                \draw[fill=black] (1,-1) circle [radius=0.07cm] node [right] {$4$};
                \draw[thick] (-1,-1) -- (-0.1,-0.1);
                \draw[thick] (0.1,0.1) -- (1,1);
                \draw[thick] (-1,1) -- (1,-1);
            \end{scope}
        \etik 
    \end{center}
    This is the content of \textit{Wick's theorem} in the path integral approach. 
\eex 
\chapter{Interacting Theories \& LSZ Theorem}

So far we have discussed free theories, but obviously physics is more interesting than this and we want to look at interactions. These interacting field theories will contain higher powers in the fields, e.g. $\phi^4$, and it is these terms that correspond to the scattering. As we are familiar from our previous courses on field theory, we cannot solve interacting field theories exactly and the usual approach is to study the theory perturbatively. This will, of course, lead to our Feynman diagrams. 

However before diving in and computing something, let's just give an outline of the kinds of things we can expect for an interacting QFT. As we saw above, for the free theory essentially everything boils down to the $2$-point function $\bra{0}\cT[\phi(x)\phi(y)]\ket{0}$, and we saw that this was just the Feynman propagator $D_F(x,y)$. The question we want to ask is "what happens to this result for an interacting theory?" That is what is the two point function equal to in the interacting theory? 

Well first we have to replace the free vacuum $\ket{0}$ with what is known as the \textit{interacting vacuum} and is denoted by $\ket{\Omega}$. This is essentially just defined the same way as for the free vacuum (i.e. it has the minimum energy) but now it is relative to the full interacting Hamiltonian, not just the free part of it. So our question is 
\bse 
    \bra{\Omega}\cT[\phi(x)\phi(y)]\ket{\Omega} = ?
\ese
The answer is what is known as \textit{K\"{a}ll\'{e}n-Lehmann spectral representation}:
\be
\label{eqn:KallenLehmann}
    \bra{\Omega} \cT[\phi(x)\phi(y)]\ket{\Omega} = \int_0^{\infty} \frac{dM^2}{2\pi} \rho(M^2) D_F(x,y;M^2),
\ee 
where $D_F(x,y;M^2)$ is the Feynman propagator for a \textit{free} field of mass $M$, and $\rho(M^2)$ is a positive \textit{spectral density} function. 

So what is this mysterious spectral density function? Well, for a free theory we see that we require
\bse 
    \rho(M^2) = 2\pi \del(M^2-m^2),
\ese 
which tells us that we just get single particle states with mass $m=M$. However it turns out for interacting theories we have\footnote{See section 7.1 of Peskin and Schroeder for more details.} 
\bse 
    \rho(M^2) = 2\pi Z \del(M^2-m^2) + \big(\text{other stuff when } M^2 \geq (2m)^2\big).
\ese 
The second term corresponds to producing two particle states. We note that we have some new factor $Z$ with our one particle states. This is called the \textit{field-strength renormalisation} and corresponds to the fact that we no longer have unit probability to produce one-particle states. It is clear that we require $0\leq Z \leq 1$, and whenever $Z>0$ we can still produce single particle states while when $Z=0$ our states \textit{must} contain at least two particles. An example of when the latter occurs is in QCD, where confinement basically says that it is impossible to produce a single quark state. We can depict this spectral density as the following plot:
\begin{center}
    \btik 
        \begin{scope}[xshift=-7cm]
            \draw[->] (-0.5,0) -- (4,0);
            \node at (4,-0.3) {$M^2$};
            \draw[->] (0,-0.5) -- (0,2);
            \node at (-0.7,2) {$\rho(M^2)$};
            \draw[ultra thick] (1,0) -- (1,2);
            \node at (1,-0.3) {$m^2$};
            \node at (2,-1) {Free Theory};
        \end{scope}
        \begin{scope}
            \draw[->] (-0.5,0) -- (4,0);
            \node at (4,-0.3) {$M^2$};
            \draw[->] (0,-0.5) -- (0,2);
            \node at (-0.7,2) {$\rho(M^2)$};
            \draw[ultra thick] (1,0) -- (1,1.5);
            \node at (1,-0.3) {$m^2$};
            \draw[ultra thick] (2,0) .. controls (2,1.5) and (3,0.55) .. (4,0.5);
            \node at (2.1,-0.3) {$(2m)^2$};
            \begin{scope}
                \clip (2,0) .. controls (2,1.5) and (3,0.55) .. (4,0.5) -- (4,0) -- (2,0);
                \draw[fill = gray!40, opacity = 0.8] (2,0) .. controls (2,1.5) and (3,0.55) .. (4,0.5) -- (4.2,0.5) -- (4.2,0) -- (2,0);
            \end{scope}
            \node at (2,-1) {Interacting Theory};
        \end{scope}
    \etik 
\end{center}

\br 
    Note that the two-particle states have a continuous spectral density (i.e. the shaded area). We can see this by considering the Fourier transform for the two-point function. Using \Cref{eqn:KallenLehmann}, this is 
    \bse 
        \int d^4 x \, e^{ip\cdot x} \bra{\Omega}\cT[\phi(x)\phi(y)]\ket{\Omega} = \frac{iZ}{p^2-m^2+i\epsilon} + \int_{\sim 4m^2}^{\infty} \frac{dM^2}{2\pi} \rho(M^2) \frac{i}{p^2-M^2+i\epsilon},
    \ese 
    where the first term is our one particle state, and the second term corresponds to the two-particle state. The former just has a single pole at $p^2=m^2$, which corresponds to the spike at $M^2=m^2$, while the second term corresponds to an entire \textit{branch cut} starting at $M^2=4m^2$ and going to infinity. It is all the points in this branch cut that give the continuous shaded area in the above plot. 
\er 

\br 
    In this course we will assume $Z>0$ at all times.
\er 

\bnn 
    In order to save time writing, in this course we shall sometimes adopt the short hand notation 
    \be 
    \label{eqn:VevNotation}
        \la \cT[...]\ra := \bra{\Omega} \cT[...]\ket{\Omega}.
    \ee 
    This will always mean the interacting vacuum expectation value (vev). That is, if we have external states we will always use the full bra-ket notation. 
\enn 

\br 
    Despite the above national definition, it is likely that I will default to writing $\ket{0}$ for the ground state. I shall try not to do this, but this is just a warning in case I do. 
\er 

\section{$\phi^4$ Theory}

The proxy we will use to study interacting field theories is the usual $\phi^4$ theory which has the following action
\be 
\label{eqn:Phi4Action}
    S[\phi] = \int d^4 x \, \bigg[ \frac{1}{2}(\p\phi)^2 - \frac{m^2}{2}\phi^2 - \frac{\l}{4!}\phi^4\bigg],
\ee
which physically says that now our particles can scatter off each other with $\l$ roughly giving the probability to scatter. Our perturbation theory expansion comes from us taking $\l$ to be a small parameter, i.e. the fields are weakly interacting, and so we categories the contributions to the full scattering process by their power of $\l$. This should be a familiar idea from other field theory courses, but if not it will become clear going forward.

\section{LSZ Reduction Formula}

Fundamentally in QFT the main thing we want to find are scattering cross sections. These answer questions like: if you smash together two protons, what is the probability we get one Higgs boson, given we know the momentum of the protons? In order to be able to answer such questions, we need some method that relates inner products of the form
\bse 
    \bra{\text{Higgs}} \text{Stuff} \ket{2\text{ protons}}
\ese 
to things we know, namely time-ordered correlation functions $\la \cT[...]\ra$. This relation is exactly the content of what is known as the \textit{LSZ theorem}, and the aim of this section is to derive this. We shall work it out in scalar field theory, but other field theories follow simply from here. We will consider $2\to 2$ scattering, again generalisations follow easily.\footnote{I have actually typed the general proof in my IFT notes, so see those for that.}

First, we need to prepare our initial and final state particles. As we are talking about states, it is instructive to temporarily use the canonical formalism, where we recall that we can write our real scalar fields in terms of the creation and annihilation operators:
\bse 
    \phi(\Vec{x},t) = \int \frac{d^3p}{(2\pi)^3} \frac{1}{\sqrt{2\omega_{\Vec{p}}}} \Big[ a_{\vec{p}} \, e^{-ip\cdot x} + a_{\vec{p}}^{\dagger}\, e^{ip\cdot x}\Big],
\ese 
with $p_{\mu}=(\omega_p,\Vec{p})$ and $\omega_{\Vec{p}}=\sqrt{\vec{p}^2+m^2}$. In the free theory the creation/annihilation operators are time independent, however in interacting theories they are time dependant, and we can write them in the form 
\bse 
    a_{\vec{p}}(t) = \frac{1}{\sqrt{2\omega_{\vec{p}}}}\int d^3 x e^{ip\cdot x} \Big( i\p_0 \phi(x) + \omega_{\vec{p}}\,  \phi(x) \Big).
\ese 
We can rewrite this as 
\be 
\label{eqn:AnnihilationOperatorLeftRightArrow}
    a_{\vec{p}}(t) = \frac{i}{\sqrt{2\omega_{\vec{p}}}} \int d^3 x e^{ip\cdot x} \lra{\p}_0 \phi(x).
\ee 
where we have introduced the definition 
\be 
\label{eqn:LeftRightArrow}
    f \lra{\p} g := f \p g - (\p f) g.
\ee 
Similarly we have
\be 
\label{eqn:CreationOperatorLeftRightArrow} 
    a_{\vec{p}}^{\dagger}(t) = -\frac{i}{\sqrt{2\omega_{\vec{p}}}} \int d^3 x e^{-ip\cdot x} \lra{\p}_0 \phi(x).
\ee

Now, for technical reasons which will become clear, we will scatter wavepackets \textit{not} momentum eigenstates. We create the envelope function 
\bse 
    f_1(\vec{k}) = \int \exp\bigg(\frac{(\vec{k}-\vec{k}_1)^2}{4\sig^2}\bigg),
\ese 
i.e. a Gaussian centred around $\vec{k}_1$. Now construct the momentum-smeared creation operator
\bse 
    a_1^{\dagger} := \int d^3 k f_1(\vec{k}) a_{\vec{k}}^{\dagger}.
\ese 

Recall that in the free theory we make initial states by acting with $a^{\dagger}_{\vec{p}}$, i.e. 
\bse 
    \ket{\vec{p}} = \sqrt{2\omega_{\vec{p}}} \, a^{\dagger}_{\vec{p}}\ket{0},
\ese
where the $\sqrt{2\omega_{\vec{p}}}$ is a relativistic normalisation factor.\footnote{See a canonical QFT course, e.g. section 4.2 of my IFT notes.} We are going to assume in the interacting theory, we make the initial state by acting with $a^{\dagger}$ in the distant past. This is our definition of the initial state: 
\bse 
    \ket{\text{initial packet}; \vec{k}_1} = \frac{\sqrt{2\omega_1}}{\sqrt{Z}} \lim_{t\to-\infty} a_1^{\dagger}(t)\ket{\Omega},
\ese 
where the $Z$ here is the same field-strength renormalisation we introduced in the spectral density above.\footnote{It's fine if it isn't clear why this is the case. To be honest I don't know why... I should really find out why!} Similarly we define \bse 
    \ket{\text{final packet}; \vec{k}_1} = \frac{\sqrt{2\omega_1}}{\sqrt{Z}} \lim_{t\to+\infty} a_1^{\dagger}(t)\ket{\Omega}.
\ese

So let's look at our $2\to2$ scattering. We shall assume the initial state fields have momenta $k_1$ and $k_2$ and the final state fields have momenta $k_1'$ and $k_2'$. Our initial state is therefore 
\bse 
    \ket{i} = \frac{\sqrt{2\omega_1} \sqrt{2\omega_2}}{Z} \lim_{t\to-\infty} a_1^{\dagger}a_2^{\dagger}\ket{\Omega}.
\ese 
Similarly the final state is 
\bse 
    \ket{i} = \frac{\sqrt{2\omega_{1'}} \sqrt{2\omega_{2'}}}{Z} \lim_{t\to+\infty} a_{1'}^{\dagger}a_{2'}^{\dagger}\ket{\Omega}.
\ese 

We want to find the inner product
\bse 
    \braket{f}{i} = \frac{\sqrt{2\omega_1} \sqrt{2\omega_2} \sqrt{2\omega_{1'}} \sqrt{2\omega_{2'}} }{Z^2} \bra{\Omega} a_{1'}(+\infty) a_{2'}(+\infty)a_1^{\dagger}(-\infty)a_2^{\dagger}(-\infty)\ket{\Omega}.
\ese 
in terms of our time-ordered correlation functions. So what do we do? Well firstly we notice that everything is already time ordered so we can simply rewrite this as
\bse 
    \braket{f}{i} = \frac{\sqrt{2\omega_1} \sqrt{2\omega_2} \sqrt{2\omega_{1'}} \sqrt{\omega_{2'}}}{Z^2} \bra{\Omega} \cT[a_{1'}(+\infty) a_{2'}(+\infty)a_1^{\dagger}(-\infty)a_2^{\dagger}(-\infty)]\ket{\Omega}.
\ese 

Ok now what do we do? Well in the free theory we know that if the annihilation operator $a$ acts on a vacuum ket vector the result vanishes, and similarly if a creation operator acts to the left on a bra vector. We want to do a similar thing here in the interacting theory, but we have to remember that now our creation/annihilation operators are time dependent. So if we want our annihilation operator to annihilate our initial state vacuum, we need to `trade' the $+\infty$ for a $-\infty$ in its argument. If we do this, the time ordering will take care of the rest, as the annihilation operator is automatically pushed all the way to the right.\footnote{You will actually pick up a bunch of delta functions when doing this. These correspond to non-fully connected diagrams and so we don't consider them here. For more information on this, see my IFT notes.} So how do we do this trading? Well, first consider their difference:
\bse 
    \begin{split}
        a_1^{\dagger}(+\infty) - a_1^{\dagger}(-\infty) & = \int_{-\infty}^{\infty} dt \, \p_t a^{\dagger}_1(t) \\
        & = \int d^3 \vec{k} \, f_1(\vec{k}) \frac{1}{\sqrt{2\omega_{\vec{k}}}} \int d^4 x \,  \p_t \Big[ e^{-ik\cdot x} \big(-i\p_t\phi(x) + \omega_{\vec{k}}\, \phi(x)\big)\Big] \\
        & = -i \int d^3\vec{k} \, f_1(\vec{k}) \frac{1}{\sqrt{2\omega_{\vec{k}}}} \int d^4 x \, e^{-ik\cdot x} \big[ \p_t^2 + \omega_{\vec{k}}^2\big] \phi(x)  \\
        & = - i\int d^3 \vec{k} \, f_1(\vec{k}) \frac{1}{\sqrt{2\omega_{\vec{k}}}} \int d^4 x \, e^{-ik\cdot x} \big( \p_t^2 - \lar{\nabla}^2 + m^2\big) \phi(x)
    \end{split}
\ese 
where the arrow above $\nabla$ means acts to the left. 
We now do integration by parts in space and use the fact that the wavepackets vanish at the boundary (this is why we used them). This gives us 
\bse 
    a_1^{\dagger}(-\infty) = a_1^{\dagger}(+\infty)  +  i\int d^3 \vec{k} \, f_1(\vec{k}) \frac{1}{\sqrt{2\omega_{\vec{k}}}} \int d^4x \, e^{-ik\cdot x} \big(\p^2+m^2\big)\phi(x).
\ese
We get a similar expression for the annihilation operator\footnote{We can get it by simply taking the dagger of the above expression. Note that we have switched the $\pm \infty$s to counteract the change in sign from the $i$.}
\bse 
    a_1(+\infty) = a_1(-\infty)  +  i\int d^3 \vec{k} \, f_1(\vec{k}) \frac{1}{\sqrt{2\omega_{\vec{k}}}} \int d^4x \, e^{ik\cdot x} \big(\p^2+m^2\big)\phi(x).
\ese

\br 
    Note for a free field the second term on the right hand sides vanish --- as they contain precisely the Klein-Gordan equation, \Cref{eqn:KleinGordanEquation} --- which we obviously want. However for an interacting field we get this additional term. 
\er 

We can now substitute these expressions into our expression for $\braket{f}{i}$. As we explained above, if we do this the $a_{1'}^{\dagger}(+\infty)$ gets replaced with a $a_{1'}(-\infty)$ + stuff and then the time ordering sends it all the way to the right and the annihilation operator gives $0$ acting on the vacuum. The same idea applies for creation operators. Finally, now that we're done with the wavepacket stuff let's set $f_1(\vec{k}) = \del(\vec{k}-\vec{k}_1)$, when we do this all our $(2\omega_{\vec{k}})^{-1/2}$ terms will become exactly the terms needed to cancel the factors in the numerator of $\braket{f}{i}$.

Now there was nothing special about $2\to2$ scattering in the above, so for general $n\to n'$ scattering, we get 
\mybox{
    \be 
    \label{eqn:LSZPosition}
        \begin{split}
            \braket{f}{i} & = i^{n+n'} \int d^4 x_1 Z^{-1/2} e^{-ik_1\cdot x_1} (\p_1^2+m^2) ... \int d^4 x_1' Z^{-1/2} e^{-ik_1'\cdot x_1'} (\p_{1'}^2+m^2) \\
            & \qquad \times \bra{\Omega}\cT[\phi(x_1)...\phi(x_1')...]\ket{\Omega}.
        \end{split}
    \ee
}
\noindent This is the LSZ reduction formula. It tells you how to relate an inner product of Hilbert space states to a time ordered product of correlation functions of fields. We can make this look nicer by writing everything in momentum space instead. 
\mybox{
    \be 
    \label{eqn:LSZMomentum}
        \prod_{i=1}^n \frac{i\sqrt{Z}}{k_i^2-m^2}\prod_{i'=1}^{n'} \frac{i\sqrt{Z}}{(k_i')^2-m^2} \braket{f}{i} = \bra{\Omega} \phi(k_1) ... \phi(k_1') ... \ket{\Omega}.
    \ee 
}
Notice that we have poles on the left-hand side. This tells us that if we want to find the transition amplitude (i.e. the inner product), compute the Fourier transformed time ordered correlation function and then put the external momenta near the on-shell values $k_i^2 \sim m^2$ and $(k_i')^2\sim m^2$. In this case the correlation function develops a pole, and the residue of the pole is the matrix element that you want. If you don't pick up the poles then it tells you perturbation theory has broken down. 

\br 
    Note that we get factors of $i$ in both \Cref{eqn:LSZPosition,eqn:LSZMomentum} even though they appear on the opposite sides of the equation. This at first might seem like a mistake,\footnote{By which I mean using $1/i = -i$.} however note that the derivatives appearing on the right-hand side of \Cref{eqn:LSZPosition} will give us $-k_i^2+m^2$, but the denominators on the left-hand side of \Cref{eqn:LSZMomentum} are $k_i^2-m^2$. This is where the minus signs are taken care of. 
\er 

\br 
    Note the pole condition above corresponds exactly to considering amputated diagrams.\footnote{If you don't know what amputated diagrams are, this remark basically explains it. If you want a pictorial version, see, e.g., section 8.2 of my IFT notes.} Why? Well because amputated diagrams are simply those where the incoming particles are taken to be on-shell. This is a very neat result of the LSZ theorem and it allows us to \textit{drastically} reduce the complexity of our problem. 
\er 

\section{Perturbation Theory \& Feynman Diagrams}

Ok we can now relate our goal, measuring transition amplitudes, to stuff we know, $N$-point functions, however there is a small problem in the interacting theory. To explain this problem, we will use $\phi^4$ theory as a proxy, but the argument trivially carries over to other interacting theories. The action is
\bse 
    S[\phi] = \int d^4 x \, \bigg[ \frac{1}{2}(\p\phi)^2 - \frac{m^2}{2}\phi^2 - \frac{\l}{4!}\phi^4\bigg].
\ese
In order to get our $N$-point functions, we need to compute the partition function
\bse 
    Z[J] = \int [\pD\phi] \exp\bigg( iS[\phi] + i\int d^4 x J(x) \phi(x)\bigg),
\ese
and then we can take functional derivatives. We saw before that we could do this nicely when our action was Gaussian-looking (i.e. quadratic in $\phi$). However now we have a $\phi^4$ term, and there is no hope for us to take the functional derivatives for general $\l$. The way we get around this is by considering $\l$ to be small (i.e. we have weak coupling) and then we Taylor expand in $\l$. To do this, let's write the action as 
\bse 
    S[\phi] = S_0[\phi] + \int d^4 x \cL_I\big(\phi(x)\big),
\ese 
with 
\bse 
    S_0[\phi] = \int d^4 x \, \bigg[ \frac{1}{2}(\p\phi)^2 - \frac{m^2}{2}\phi^2\bigg], \qand \cL_I = -\frac{\l}{4!}\phi^4.
\ese 

How does this help us? Well we recall from QFT I the argument made leading up to the Schwinger Dyson equation: that we can essentially replace $\phi$ factors in the path integral with functional derivatives of $J$. That is, we can replace 
\bse 
    \cL_I(\phi) \to \cL \bigg(-i\frac{\del}{\del J(x)}\bigg)
\ese 
in the path integral. This is now explicitly independent\footnote{Of course it is still \textit{implicitly} dependent} of $\phi$ and so we can take it outside our path integral, leaving us with a path integral over our free theory, $S_0[\phi]$. 

Explicitly what we're saying above is the following:
\bse 
    \begin{split}
        Z[J] & = \int [\pD\phi] \exp\bigg( iS_0[\phi] + i\int d^4x \big[\cL_I(\phi) + J(x) \phi(x)\big]\bigg) \\
        & =  \exp\bigg( i\int d^4 x \cL\bigg(\frac{\del}{\del J(x)}\bigg)\bigg) \int [\pD\phi] \exp\bigg( iS_0[\phi] + i\int d^4x J(x) \phi(x)\big]\bigg) \\
        & = \exp\bigg( i\int d^4 x \cL\bigg(\frac{\del}{\del J(x)}\bigg)\bigg)  Z_0[J],
    \end{split}
\ese 
where $Z_0[J]$ is given by \Cref{eqn:PartitionDetKG}. Plugging this result in\footnote{All of the normalisation $\cN$ etc is contained in $Z_0[0]$.}
\bse 
    Z[J] = Z_0[0] \exp\bigg( i\int d^4 x \cL\bigg(-i\frac{\del}{\del J(x)}\bigg)\bigg) \exp\bigg( -\frac{1}{2}\int d^4x d^4y J(x) D_F(x,y) J(y)\bigg)
\ese 

This is a formal expression which works for any theory with any interactions. Let us now specialise to $\phi^4$ theory. Here we have 
\bse 
    Z[J] \propto \exp\Bigg[ -\frac{i\l}{4!}\int d^4 x \bigg(-i\frac{\del}{\del J(x)}\bigg)^4\Bigg] \exp\bigg( -\frac{1}{2}\int d^4x d^4y J(x) D_F(x,y) J(y)\bigg)
\ese 
We expand both sets of exponentials, and the result is exactly the Feynman diagram expansion. We get 
\bse 
    \begin{split}
        \exp\Bigg[-\frac{i\l}{4!}\int d^4 x \bigg(-i\frac{\del}{\del J(x)}\bigg)^4\Bigg] = 1 - \frac{i\l}{4!}\int d^4 x  \bigg(-i\frac{\del}{\del J(x)}\bigg)^4 + \frac{1}{2} \Bigg[ -\frac{i\l}{4!}\int d^4 x \bigg(-i\frac{\del}{\del J(x)}\bigg)^4\Bigg]^2 + ...
    \end{split}
\ese 
and 
\bse 
    \exp\bigg(-\frac{1}{2}J\cdot D_F\cdot J\bigg) = 1 - \frac{1}{2}J\cdot D_F\cdot J + \frac{1}{2}\bigg(- \frac{1}{2}J\cdot D_F\cdot J \bigg)^2 + ...,
\ese 
where we have introduced the notation\footnote{See section 4.1 of my QFT I notes for more examples using this notation.} 
\bse 
    J\cdot D_F \cdot J := \int d^4x d^4 y J(x) D_F(x,y) J(y).
\ese 

We then simply put this expansion together and look at the terms. The first non-trivial\footnote{I.e. not just $1$ but we actually have functional derivatives acting.} term is 
\bse 
    -\frac{i\l}{4!}\int d^4 x \bigg(-i\frac{\del}{\del J(x)}\bigg)^4 \frac{1}{2}\bigg(-\frac{1}{2}J\cdot D_F\cdot J\bigg) \bigg(- \frac{1}{2}J\cdot D_F\cdot J \bigg) = -\frac{N}{4!} i\l \int d^4x D_F(x,x) D_F(x,x),
\ese 
where the $N$ on the right-hand side is a number that will be explained shortly. We can depict this result as a Feynman diagram of the form
\begin{center}
    \btik 
        \draw[thick] (-0.75,0) circle [radius=0.75cm];
        \draw[thick] (0.75,0) circle [radius=0.75cm];
        \draw[fill=black] (0,0) circle [radius=0.07cm];
    \etik 
\end{center}
We will explain in just a moment how to obtain the Feynman diagrams, but let's give another example first. The term 
\bse 
    \frac{1}{2} \bigg(\frac{i\l}{4!}\bigg)^2 \int d^4 x d^4 y  \bigg(-i\frac{\del}{\del J(x)}\bigg)^4  \bigg(-i\frac{\del}{\del J(x)}\bigg)^4  \frac{1}{2} \bigg(-\frac{1}{2}J\cdot D\cdot J\bigg)^4.
\ese 
corresponds to a diagram with $4$ propagators and $2$ vertices. There are three such diagrams, which we draw below:
\begin{center}
    \btik 
        \begin{scope}
            \draw (-3.5,-1) -- (-3.5,1) -- (3.5,1) -- (3.5,-1) -- (-3.5,-1);
            \draw[thick] (-2.5,0) circle [radius=0.75cm];
            \draw[thick] (-1,0) circle [radius=0.75cm];
            \draw[fill=black] (-1.75,0) circle [radius=0.07cm];
            %
            \draw[thick] (1,0) circle [radius=0.75cm];
            \draw[thick] (2.5,0) circle [radius=0.75cm];
            \draw[fill=black] (1.75,0) circle [radius=0.07cm];
        \end{scope}
        \begin{scope}[xshift=6.5cm]
            \draw (-2.5,-1) -- (-2.5,1) -- (2.5,1) -- (2.5,-1) -- (-2.5,-1);
            \draw[thick] (-1.5,0) circle [radius=0.75cm];
            \draw[thick] (0,0) circle [radius=0.75cm];
            \draw[thick] (1.5,0) circle [radius=0.75cm];
            \draw[fill=black] (-0.75,0) circle [radius=0.07cm];
            \draw[fill=black] (0.75,0) circle [radius=0.07cm];
        \end{scope}
        \begin{scope}[xshift=10.5cm]
            \draw (-1,-1) -- (-1,1) -- (1,1) -- (1,-1) -- (-1,-1);
            \draw[thick] (0,0) circle [radius=0.75cm];
            \draw[thick] (-0.75,0) .. controls (-0.25,0.5) and (0.25,0.5) .. (0.75,0);
            \draw[thick] (-0.75,0) .. controls (-0.25,-0.5) and (0.25,-0.5) .. (0.75,0);
            \draw[fill=black] (-0.75,0) circle [radius=0.07cm];
            \draw[fill=black] (0.75,0) circle [radius=0.07cm];
        \end{scope}
    \etik 
\end{center}
where the boxes around the diagrams mean nothing but are just put there to distinguish the three separate diagrams.

\br 
\label{rem:DiagramsWithJ}
    The diagrams we have drawn above correspond to terms in the expansion which have no $J$s left over, i.e. there are the same number of $J$s as there are functional derivatives. Of course we also get terms where these two don't agree and we have $J$s left over after we've done the functional derivatives. For example the term 
    \bse 
        -\frac{i\l}{4!} \int d^4x \bigg(-i\frac{\del}{\del J(x)}\bigg)^4 \frac{1}{2}\bigg(-\frac{1}{2}J\cdot D_F \cdot J\bigg)^4
    \ese 
    will have $4$ $J$s left over. This corresponds to two diagrams:
    \begin{center}
        \btik 
            \begin{scope}[xshift=-3cm]
                \draw[thick] (1,1) -- (-1,-1);
                \draw[thick] (1,-1) -- (-1,1);
                \draw[fill=black] (0,0) circle [radius=0.07cm];
                \node at (-1.2,1.2) {$J$};
                \node at (-1.2,-1.2) {$J$};
                \node at (1.2,1.2) {$J$};
                \node at (1.2,-1.2) {$J$};
            \end{scope}
            \node at (0,0) {and};
            \begin{scope}[xshift=3cm]
                \draw[thick] (-0.75,0) circle [radius=0.75cm];
                \draw[thick] (0.75,0) circle [radius=0.75cm];
                \draw[fill=black] (0,0) circle [radius=0.07cm];
                \draw[thick] (2.5,1) -- (3.5,1);
                \draw[thick] (2.5,-1) -- (3.5,-1);
                \node at (2.3,1) {$J$};
                \node at (2.3,-1) {$J$};
                \node at (3.7,1) {$J$};
                \node at (3.7,-1) {$J$};
            \end{scope}
        \etik  
    \end{center}
    However because ultimately what we're interested in is the $N$-point Green's function, and this contains a restriction to $J=0$ at the end (see \Cref{eqn:2PointFunctionDefinition}) terms like this will vanish, so we can essentially forget about them. It is important to realise that terms like this do exist though, its only that they vanish when we consider $N$-point functions with $N$ less then the number of $J$s appearing in the expansion.
\er 

Ok so how did we draw the above diagrams? Well we used the position space Feynman rules, which we now list: 

\mybox{
    \begin{enumerate}
        \item Each \textit{vertex}, 
        \begin{center}
            \btik 
                \draw[thick] (-0.5,-0.5) -- (0.5,0.5);
                \draw[thick] (-0.5,0.5) -- (0.5,-0.5);
                \draw[fill=black] (0,0) circle [radius=0.07cm];
            \etik 
        \end{center}
        comes with a factor of
        \bse 
            -i\l \int d^4 x 
        \ese 
        \item Each propagator
        \begin{center}
            \btik 
                \draw[thick] (-1,0) -- (1,0);
                \node at (-1,-0.3) {$x$};
                \node at (1,-0.3) {$y$};
            \etik 
        \end{center}
        comes with a $D_F(x,y)$ factor.
        \item Each external $J$ 
        \begin{center}
            \btik 
                \draw[thick] (-1,0) -- (1,0);
                \draw[fill=black] (1,0) circle [radius=0.07cm];
                \node at (1,-0.3) {$J$};
            \etik 
        \end{center}
        comes with $\int d^4 x J(x)$ factor.
    \end{enumerate}
}
Before moving on to the next subsection, let's just make few comments.
\ben[label=(\roman*)]
    \item Note for the vertex we loose the $1/4!$ that appears with the functional derivatives. This is essentially because there are $4!$ different ways to arrange the $4$ functional derivatives. This means that this diagram actually contains $4!$ different terms, corresponding to the combinatorics. This will hopefully be more clear in the following subsection.
    \item Both the vertex and an external $J$ come with integrals $\int d^4x$. These two integrals should not be confused: the vertex integral corresponds to the overall integral that appears all the way to the left before the functional derivatives; the external $J$ integral corresponds to the integrals appearing in $J\cdot G \cdot J$.
    \item Note that we label the ends of the propagators with $x/y$. This allows us to see what is contracted with what, in the sense that we integrate them and we need to know which integration variable goes where. 
    \item Relating the two above comments, the integration variable $x$ appearing in the vertex term corresponds to labelling the vertex point (filled in circle) with a variable $x$. 
    \item As we said above, these are the \textit{position space} Feynman rules. We will shortly present them in momentum space where they are nicer. 
\een 

\bbox 
    Convince yourself that the two diagrams in \Cref{rem:DiagramsWithJ} are indeed the two diagrams you get. \textit{Hint: Take the derivatives appearing in the expression, and note that you are left with four factors of $D_F$, four factors of $J$ and one factor of $-i\l$. This corresponds to diagrams that have one vertex, four propagators (i.e. solid lines) and four $J$ insertions (i.e. end points to lines).}
\ebox 

\subsection{Combinatorics And Symmetry Factors}

As we mentioned in (i) above, each diagram contains potential combinatoric factors, and we obviously need some way in order to see these. Let's look at these more carefully now.

A general term takes the form 
\be 
\label{eqn:GeneralDiagram}
    \sum_{V,P=0}^{\infty} \frac{1}{V!} \Bigg[ -i\int d^4 x \frac{\l}{4!}\bigg(-i\frac{\del}{\del J(x)}\bigg)^4\Bigg]^V \frac{1}{P!} \bigg(- \frac{1}{2}J\cdot D\cdot J \bigg)^P
\ee 
where $V$ is the number of vertices and $P$ the number of propagators. When we draw the diagram really we are capturing multiple different terms from this expression, corresponding to the different ways to get the same result by contracting with the functional derivatives in different ways. The question is "how many different contraction patterns does it correspond to?" 

Well na\"{i}vely, we have $4!$ different ways to rearrange the functional derivatives at each vertex, $V!$ different ways to rearrange the vertices amongst themselves, similarly $P!$ different ways to arrange $P$ propagators amongst themselves, and finally we have $2$ ways to pick which end of the propagator goes where. So we expect to get a factor of
\bse 
    (4!)^V V! P! 2^P,
\ese 
which has the exact form to cancel all the denominators in \Cref{eqn:GeneralDiagram}. 

Now recall before we introduced a factor of $N$ in our first non-trivial diagram (the \textit{two kissing-bubble} diagram) and said we would return to it. This $N$ corresponds exactly\footnote{Well we have to strip off a factor of $1/4!$ because that appeared explicitly in our expression.} to these combinatoric calculations, and so we have just argued  that $N$ is trivial. 

Well, turns out we lied; this isn't true, as sometimes we over count. When we considered the rearrangements we said that each one was unique, but this isn't true, as some `cancel' each other. To see what we mean let's consider our two kissing bubble example: 
\begin{center}
    \btik 
        \draw[thick] (-0.75,0) circle [radius=0.75cm];
        \draw[thick] (0.75,0) circle [radius=0.75cm];
        \draw[fill=black] (0,0) circle [radius=0.07cm];
        %
        \draw[ultra thick,->] (2,0) -- (3.75,0);
        %
        \begin{scope}
            \clip (5,-1) -- (5,1) -- (4,1) -- (4,-1) -- (5,-1);
            \draw[thick] (5,0) circle [radius=0.75cm];
        \end{scope}
        \draw[fill=red] (5,0.75) circle [radius=0.07cm];
        \draw[fill=blue] (5,-0.75) circle [radius=0.07cm];
        %
        \draw[thick] (5.5,0.75) -- (7,-0.75);
        \draw[thick] (5.5,-0.75) -- (7,0.75);
        \draw[fill=black] (6.25,0) circle [radius=0.07cm];
        \draw[fill=red] (5.5,0.75) circle [radius=0.07cm];
        \draw[fill=blue] (5.5,-0.75) circle [radius=0.07cm];
        \draw[fill=green] (7,0.75) circle [radius=0.07cm];
        \draw[fill=yellow] (7,-0.75) circle [radius=0.07cm];
        %
        \begin{scope}
            \clip (7.5,-1) -- (7.5,1) -- (8.5,1) -- (8.5,-1) -- (7.5,-1);
            \draw[thick] (7.5,0) circle [radius=0.75cm];
        \end{scope}
        \draw[fill=green] (7.5,0.75) circle [radius=0.07cm];
        \draw[fill=yellow] (7.5,-0.75) circle [radius=0.07cm];
    \etik 
\end{center}
where the diagram on the right is meant to represent the `breaking apart' of the diagram into one vertex and two propagators. The coloured circles tell us how to `glue it back together' to get the diagram we started with. What our above na\"{i}ve calculation said was "there are $(4!)^1 1! 2! 2^2$ ways to glue this back together with a different colour gluing process" which clearly isn't the case. 

Now that sentence was probably not super clear, so let's give an example of what \textit{doesn't} give a new "colour gluing process" to help clarify what we mean. Imagine we swapped the red tipped and blue tipped lines coming out the vertex and \textit{also} swapped the red and blue ends of the left detached propagator. Then we get 
\begin{center}
    \btik 
        \begin{scope}
            \clip (5,-1) -- (5,1) -- (4,1) -- (4,-1) -- (5,-1);
            \draw[thick] (5,0) circle [radius=0.75cm];
        \end{scope}
        \draw[fill=blue] (5,0.75) circle [radius=0.07cm];
        \draw[fill=red] (5,-0.75) circle [radius=0.07cm];
        %
        \draw[thick] (5.5,0.75) -- (7,-0.75);
        \draw[thick] (5.5,-0.75) -- (7,0.75);
        \draw[fill=black] (6.25,0) circle [radius=0.07cm];
        \draw[fill=blue] (5.5,0.75) circle [radius=0.07cm];
        \draw[fill=red] (5.5,-0.75) circle [radius=0.07cm];
        \draw[fill=green] (7,0.75) circle [radius=0.07cm];
        \draw[fill=yellow] (7,-0.75) circle [radius=0.07cm];
        %
        \begin{scope}
            \clip (7.5,-1) -- (7.5,1) -- (8.5,1) -- (8.5,-1) -- (7.5,-1);
            \draw[thick] (7.5,0) circle [radius=0.75cm];
        \end{scope}
        \draw[fill=green] (7.5,0.75) circle [radius=0.07cm];
        \draw[fill=yellow] (7.5,-0.75) circle [radius=0.07cm];
    \etik 
\end{center}
but clearly if we just glue this back together we get the same thing, i.e. red with red and blue with blue, so we shouldn't count it twice in our combinatorics. This process corresponds to us switching two of the functional derivatives round while simultaneously swapping the two $J$s that appear in the $J\cdot D_F \cdot J$ that these derivatives act on. 

Contrast this with the "colour gluing process" of \textit{only} swapping the vertex lines (i.e. just swapping the functional derivatives), which would correspond to 
\begin{center}
    \btik 
        \begin{scope}
            \clip (5,-1) -- (5,1) -- (4,1) -- (4,-1) -- (5,-1);
            \draw[thick] (5,0) circle [radius=0.75cm];
        \end{scope}
        \draw[fill=red] (5,0.75) circle [radius=0.07cm];
        \draw[fill=blue] (5,-0.75) circle [radius=0.07cm];
        %
        \draw[thick] (5.5,0.75) -- (7,-0.75);
        \draw[thick] (5.5,-0.75) -- (7,0.75);
        \draw[fill=black] (6.25,0) circle [radius=0.07cm];
        \draw[fill=blue] (5.5,0.75) circle [radius=0.07cm];
        \draw[fill=red] (5.5,-0.75) circle [radius=0.07cm];
        \draw[fill=green] (7,0.75) circle [radius=0.07cm];
        \draw[fill=yellow] (7,-0.75) circle [radius=0.07cm];
        %
        \begin{scope}
            \clip (7.5,-1) -- (7.5,1) -- (8.5,1) -- (8.5,-1) -- (7.5,-1);
            \draw[thick] (7.5,0) circle [radius=0.75cm];
        \end{scope}
        \draw[fill=green] (7.5,0.75) circle [radius=0.07cm];
        \draw[fill=yellow] (7.5,-0.75) circle [radius=0.07cm];
    \etik 
\end{center}
but now we have red with blue and blue with red, which \textit{is} different. 

\bbox
    Prove that our na\"{i}ve argument for the two kissing bubble diagram over counts by $2\cdot 2\cdot 2 = 8$. 
    \textit{Hint: The first factor of $2$ is the one just explained, the second factor should be trivial given this, and the third factor shouldn't be too hard to work out.}
\ebox 

The $8$ appearing in the above exercise is the \textit{symmetry factor}, $S$, of this diagram.\footnote{Note this is clearly the same symmetry factor as the one we encountered in the IFT course, and the argument used to obtain it is essentially the same, just there we talked about permutations of Wick contractions and here we're talking about different ways to take functional derivatives.} It is the number of "colour gluing" equivalent ways to draw the same diagram. Therefore if we want to convert our diagram back into the maths we have to divide by this factor to get the correct answer, that is our two kissing bubbles is given by the expression
\bse 
    \frac{i\l}{8} \int d^4 x D_F(x,x) D_F(x,x). 
\ese 

\bbox 
    Verify that this result is true explicitly. That is, show that 
    \bse 
        -\frac{i\l}{4!}\int d^4 x \bigg(-i\frac{\del}{\del J(x)}\bigg)^4 \frac{1}{2}\bigg(-\frac{1}{2}J\cdot D_F\cdot J\bigg) \bigg(- \frac{1}{2}J\cdot D_F\cdot J \bigg)
    \ese 
    contains exactly this term in it. \textit{Hint: If you took QFT I at the same time as me,\footnote{Equally if Dr. Smith hasn't changed the questions.} this result should look somewhat familiar from the homework sheet...}
\ebox  

\bbox 
    Show that the symmetry factor for the following diagram is $4! \cdot 2$. 
    \begin{center}
        \btik 
            \draw[thick] (0,0) circle [radius=0.75cm];
            \draw[thick] (-0.75,0) .. controls (-0.25,0.5) and (0.25,0.5) .. (0.75,0);
            \draw[thick] (-0.75,0) .. controls (-0.25,-0.5) and (0.25,-0.5) .. (0.75,0);
            \draw[fill=black] (-0.75,0) circle [radius=0.07cm];
            \draw[fill=black] (0.75,0) circle [radius=0.07cm];
        \etik 
    \end{center}
    \textit{Hint: Break it up into two vertices connected by $4$ propagators, and then use the result to see where the symmetry factors lie.}
\ebox 

\section{Connected Diagrams}

Ok cool, we have seen how to represent a particular $N$-point function in terms of diagrams. Now recall from QFT I that the partition function is essentially a sum over all orders of $N$-point functions. As we have seen for a given $N$-point function, we get multiple different diagrams and some of these look like two copies of the same thing. For example
\begin{center}
    \btik 
        \begin{scope}[xshift=-1.75cm]
            \draw[thick] (-0.75,0) circle [radius=0.75cm];
            \draw[thick] (0.75,0) circle [radius=0.75cm];
            \draw[fill=black] (0,0) circle [radius=0.07cm];
        \end{scope}
        \begin{scope}[xshift=1.75cm]
            \draw[thick] (-0.75,0) circle [radius=0.75cm];
            \draw[thick] (0.75,0) circle [radius=0.75cm];
            \draw[fill=black] (0,0) circle [radius=0.07cm];
        \end{scope}
    \etik 
\end{center}
looks like two copies of our two kissing bubbles. We call diagrams of this form \textit{disconnected}. If the diagram is not disconnected it is \textit{connected}.\footnote{Note Prof. Spannowsky made a further distinction in IFT between "connected" and "fully connected". In this course, connected means fully connected, and disconnected means both connected and disconnected in the IFT sense.} 

The most generic contribution to $Z[J]$ will be from disconnected diagrams, and so there is a natural question to ask: "can we just consider the connected diagrams and still get the full $Z[J]$?" The answer is yes, and we shall now describe how. 

The idea is label each connected diagram by $I$ and denote the contribution of this diagram $C_I$, then the general disconnected diagram gives contribution
\bse 
    D = \frac{1}{S_D} \prod_{I} \big(C_I\big)^{n_I}
\ese 
where $I$ runs over the set of diagrams, $n_I$ is the number of times the connected diagram appears, and $S_D$ 
\bse 
    S_D := \prod_I n_I!
\ese 
is a new symmetry factor which takes care of permuting the connected subdiagrams (e.g. for the one above, it's $2$ as we can swap the two kissing bubbles over and get the same thing). 

Putting this all together we have 
\bse 
    \begin{split}
        Z[J] & = \sum_{\{n_I\}} \prod_I \frac{1}{n_I!} \big(C_I\big)^{n_I} \\
        & = \prod_I \sum_{n_I=0}^{\infty} \frac{1}{n_I!} \big(C_I\big)^{n_I} \\
        & = \prod_I \exp\big(C_I\big) \\
        & = \exp\Big(\sum_IC_I\Big)
    \end{split}
\ese 
This tells us that the \textit{full} partition function is just given by an exponential of all the connected diagrams. This is a super nice result and massively reduces the number of different diagrams we have to compute. This result is actually so nice it is convenient to introduce a new generating functional $W[J]$ in the following manner:
\mybox{
    \be
    \label{eqn:W[J]}
        iW[J] := \sum_I C_I \qquad \implies \qquad Z[J] = \exp\big(iW[J]\big).
    \ee 
}

\section{Scattering Amplitudes}

Right so we have spent a long time developing/massaging the mathematics for calculating cross sections, let's now actually find one. We will use our $\phi^4$ theory, and consider the $2\to 2$ scattering we talked about while deriving the LSZ theorem. That is we want to find
\bse 
    \braket{f}{i} = \braket{k_1',k_2';\text{out}}{k_1,k_2;\text{in}}.
\ese 
The LSZ formula givs us
\bse 
    \braket{f}{i} = \bigg(\frac{i}{\sqrt{Z}}\bigg)^{2+2} \int d^4x_1 e^{-ik_1\cdot x_1} ... d^4x_2' e^{-ik_2'\cdot x_2'}  \bra{0}T\psi(x_1)\phi(x_2)\phi(x_1')\phi(x_2') \ket{0} (k_1^2 - m^2) ... (k_2'^2 - m^2).
\ese 
The notation is already getting quite heavy and long, so let's introduce some short hand notations. We define
\bse 
    G^{(4)}(\{x_i\}) := G\big(\phi(x_1),\phi(x_2),\phi(x_1')\phi(x_2')\big) =  \bra{0}T\psi(x_1)\phi(x_2)\phi(x_1')\phi(x_2') \ket{0},
\ese 
and
\bse 
    \del_{i} := \frac{-i\del}{\del J(x_i)}
\ese
so that 
\bse 
    G^{(4)}(\{x_i\}) = \frac{1}{Z[0]} \del_1\del_2\del_1'\del_2' Z[J]\Big|_{J=0}.
\ese 

As it is the four point function we are considering, it is only diagram with four external legs that will survive. This is just the statement that we have four functional derivatives and so, because we set $J=0$ at the end, it is only the terms in $Z[J]$ with exactly four $J$s that survive. Let's do this order by order in $\l$. 
\ben
    \item Order $\l^0$: 
    \begin{center}
        \btik
            \begin{scope}[xshift=-3cm]
                \draw[thick] (-0.5,0.5) -- (-0.5,-0.5);
                \draw[thick] (0.5,0.5) -- (0.5,-0.5);
                \node at (-0.7,0.7) {$1$};
                \node at (-0.7,-0.7) {$2$};
                \node at (0.7,0.7) {$1'$};
                \node at (0.7,-0.7) {$2'$};
            \end{scope}
            %
            \begin{scope}
                \draw[thick] (-0.5,0.5) -- (0.5,0.5);
                \draw[thick] (-0.5,-0.5) -- (0.5,-0.5);
                \node at (-0.7,0.7) {$1$};
                \node at (-0.7,-0.7) {$2$};
                \node at (0.7,0.7) {$1'$};
                \node at (0.7,-0.7) {$2'$};
            \end{scope}
            \begin{scope}[xshift=3cm]
                \draw[thick] (-0.5,0.5) -- (0.5,-0.5);
                \draw[thick] (-0.5,-0.5) -- (-0.1,-0.1);
                \draw[thick] (0.1,0.1) -- (0.5,0.5);
                \node at (-0.7,0.7) {$1$};
                \node at (-0.7,-0.7) {$2$};
                \node at (0.7,0.7) {$1'$};
                \node at (0.7,-0.7) {$2'$};
            \end{scope}
        \etik 
    \end{center}
    The claim is these claim is these don't contribute to our scattering. Why? Well let's look at the Fourier transform expressions, the second term is 
    \bse 
        \int dx_1 e^{-ik_1\cdot x_1} ... dx_1' e^{ik_1'\cdot x_1'} D_F(x_1,x_1') D_F(x_2,x_2').
    \ese 
    However the Feynman propagator is transnational invariant, meaning it only depends on the \textit{difference} of its variables. We therefore have  
    \bse 
        \int dx_1 e^{-ik_1\cdot x_1} ... dx_1' e^{ik_1'\cdot x_1'} D_F(x_1-x_1') D_F(x_2-x_2').
    \ese 
    Now it would be nice if we could use some kind of delta function relation like $\int dx e^{-ix(a-b)} = \del(a-b)$. Well we note that Feynman propagators only depend on the difference and not the sum, i.e. $x_1+x_1'$ doesn't appear in them. We therefore split the measure 
    \bse 
        dx_1dx_1' = d(x_1+x_1')d(x_1-x_1')
    \ese
    and terms appearing the exponential
    \bse 
        k_1'\cdot x_1' - k_1\cdot x_1 = \frac{1}{2}\big(k_1-k_1'\big)\big(x_1+x_1'\big) + \frac{1}{2}\big(k_1+k_1'\big)\big(x_1-x_1'\big)
    \ese 
    so that we can do the integral over $x_1+x_1'$ and get a delta function. Our scattering amplitude is therefore proportional to
    \bse
        \del^{(4)}(k_1-k_1').
    \ese 
    We can do a similar thing to obtain a $\del^{(4)}(k_2-k_2')$. 
    
    What this result is telling us is that the ingoing momentum is \textit{the same} as the outgoing momentum, that is there is \textit{no scattering}. This is clear from the picture: the two initial state particles just become the final state ones without interacting at all. The same kind of argument can be made for the other two diagrams.
    
    So we see that scattering only comes from connected diagrams, so from now on we will only consider connected ones. This means that in the LSZ formula really we should just be considering\footnote{Note $\log Z[J] =iW[J]$, so we don't need the $1/Z[0]$ factor.}
    \bse 
        G^{(4)}(\{x_i\})_C = \del_1\del_2\del_1'\del_2' (iW[J])\big|_{J=0},
    \ese
    This makes the above calculation obvious because W[J] is trivial to $\l^0$ (i.e. it contains no $J$s at this order).
    \item Go to first non-trivial order, $\l^1$. 
    \begin{center}
        \btik 
            \draw[thick] (1,1) -- (-1,-1);
            \draw[thick] (-1,1) -- (1,-1);
            \draw[fill=black] (0,0) circle [radius=0.07cm];
            \node at (-1.2,1.2) {$1$};
            \node at (-1.2,-1.2) {$2$};
            \node at (1.2,1.2) {$1'$};
            \node at (1.2,-1.2) {$2'$};
        \etik  
    \end{center}
    The position space expression for this is\footnote{Note that each Feynman propagator contains a $y$. This is because we have labelled the intersection point by $y$, hopefully this clears up any confusion regarding the comments made after the position space Feynman rules.} 
    \bse 
        G^{(4)} = -i\l \int d^4y D_F(x_1-y)D_F(x_2-y)D_F(x_1'-y)D_F(x_2'-y).
    \ese 
    As we know from IFT, it is often very useful to go to momentum space. The Feynman propagator is given by 
    \bse 
        D_F(x-y) \int \frac{d^4p}{(2\pi)^4} e^{\pm ip(x-y)} \frac{i}{p^2-m^2}.
    \ese 
    Note we are free to pick the sign in the exponential, and we pick the one that agrees with the inverse Fourier transform appearing in the LSZ formula. This corresponds to picking the positive sign for in states and negative sign for the out states. We therefore have 
    \bse
        \begin{split}
            G^{(4)}_C(\{x_i\}) = -i\l \int d^4y \prod_{i=1}^4 \bigg[\frac{d^4p_i}{(2\pi)^4}\bigg]e^{ip_1(x_1-y) + ip_2(x_2-y) - ip_1'(x_1'-y) -ip_2'(x_2'-y)} \frac{i}{p_1^2-m^2} ... \frac{i}{p_2'^2-m^2}
        \end{split}
    \ese 
    If we do the $y$ integral, we will get a factor of $(2\pi)^4\del^{(4)}(p_1+p_2-p_1'-p_2')$. We now note that this delta function is stemming from the fact that our four lines meet at the vertex, which we labelled by $y$. Similarly we note that each propagator comes with a factor of $(p^2-m^2)^{-1}$. These comments will be useful in just a moment. 
    
    If we insert this expression into our LSZ formula, \Cref{eqn:LSZPosition}, we see that the Klein-Gordan operators, $\p^2+m^2$, will give us exactly the $p^2+m^2$ factors needed to cancel the propagator terms. Next the integrals over the $x/x'$s in \Cref{eqn:LSZPosition} will act on our exponentials and give us more delta functions $\del^{(4)}(p_i-k_i)$, and similarly for the primed momenta. What this does is tie the momenta on the legs, $p_i$, to the physical momenta $k_i$.
    
    Putting this all together we see that all the remains in the scattering amplitude is
    \bse 
        \braket{k_1',k_2';\text{out}}{k_1,k_2;\text{in}} = -i\l (2\pi)^4 \del^{(4)}(k_1+k_2-k_1'-k_2'),
    \ese
    where we have set $Z=1$. It turns out that corrections to this may arise at $\cO(\l^2)$, we will return to this shortly.
\een 

Very good, we have found out first cross section. It was a lot of work but it has given us quite a nice insight into what's at play here. We note that the delta function in our answer above is simply the statement that the total $4$-momentum for the process is conserved. This is simply because we start with total $4$-momentum $k_1+k_2$ and we finish with $k_1'+k_2'$. We obviously expect this to appear in \textit{every} physical scatting process, and so we can strip it off and simply remember to include it right at the end. It is because of this we introduce what is known as the \textit{invariant matrix element}, $i\cM$. For a general scattering it is given by
\bse 
    \braket{k_1',k_2';\text{out}}{k_1,k_2;\text{in}} = i\cM(k) (2\pi)^4 \del^{(4)}(k_{\text{in}}-k_{\text{out}}'),
\ese
where $k_{\text{in}}$ means the total incoming momenta and similarly for $k_{\text{out}}$. In our order $\l$ calculation above, we see that $i\cM = -i\l$. 

We are now in a position to state the momentum space Feynman rules, which are hopefully obvious given the comments etc. made along the way. 
\mybox{
    Momentum space Feynman rules: 
    \ben 
        \item For each amplitude you want, draw the external legs. Label the momentum flow on the legs with an arrow.
        \item Draw all possible vertices to connect the lines.
        \item Each vertex gives factor $-i\l$. 
        \item Put a momentum on each line. If the line is external, this is the momentum of the ingoing/outgoing particle. 
        \item For each internal line, include a factor 
        \bse 
            \frac{i}{p^2-m^2}
        \ese 
        \item For each external line, we have a factor $i/\sqrt{Z}$. 
        \item Make sure momentum is conserved at each vertex. 
        \item Integrate over any unfixed momenta, with measure $d^4k/(2\pi)^4$. 
    \een
}

Perhaps some comments on each rule for clarity: 
\ben
    \item This should be reasonably straight forward.
    \item This corresponds to considering all the different diagrams. 
    \item This carries straight over from the position space rules. 
    \item Again should be reasonably straight forward. For internal lines (i.e. ones that aren't the ones we drew in step 1) the momentum is not known before hand. It is an integration variable and will be fixed by momentum constraints, unless we have loops. This is what steps $6$ and $7$ take care of. 
    \item This comes from the fact that we said every Feynman propagator comes with this factor. We saw that for \textit{external} states the LSZ formula will cancel these, however for internal states nothing will cancel them so we need to include them. 
    \item This comes straight from the LSZ formula. 
    \item This says the momentum is conserved locally at the vertex. This allows us to remove a bunch of our momentum integrals mention in step 4. 
    \item If we have loops (see next chapter), then we won't have enough delta functions to cancel all our integrals, so we need to include these remaining integrals. 
\een 
\chapter{Loops and Renormalisation}

So far all that we have considered is \textit{tree-level} diagrams. These diagrams are a bit boring because they always work out nicely. This is a repercussion of the fact we have the same number of integrals as we do delta functions. We know from our canonical studies that loops mess this equality up and we end up with an undetermined momenta.\footnote{If it is not clear why this is the case, just draw a loop diagram and check. It basically comes from the fact that at least one of the internal momenta will appear in two delta functions, so integrating over it kills both!}

Now loops might seem a bit of a pain because of this, however it's actually where a lot of the interesting results lie. Indeed it's a property that only appears in the quantum theory and so the presence of loops can be thought of as an important differentiation of the classical and quantum theory. Loops also give rise to divergent terms and we need to renormalise them. So it's important we study them properly. 

\section{Loops in $\phi^4$}

We shall again use our $\phi^4$ theory as a proxy to study loops. Let's continue to consider our $2\to 2$ scattering, but now look at the $\l^2$ diagrams. There are three such diagrams:
\begin{center}
    \btik 
        \begin{scope}[scale=1.5, xshift=-3.5cm]
            \midarrow (1,1) -- (0,0.5);
            \node at (1.1,1.1) {$2$};
            \node at (0.4,0.9) {$p_2$};
            \midarrow (-1,1) -- (0,0.5);
            \node at (-1.1,1.1) {$1$};
            \node at (-0.4,0.9) {$p_1$};
            \midarrow (0,-0.5) -- (1,-1);
            \node at (1.1,-1.1) {$2'$};
            \midarrow (0,-0.5) -- (-1,-1);
            \node at (-1.1,-1.1) {$1'$};
            \draw[thick, decoration={markings, mark=at position 0.5 with {\arrow{<}}}, postaction={decorate}] (0,0) circle [radius=0.5cm];
            \draw[thick, decoration={markings, mark=at position 0.5 with {\arrow{<}}}, postaction={decorate}] (0,0) circle [radius=-0.5cm];
            \draw[fill=black] (0,0.5) circle [radius=0.047cm];
            \draw[fill=black] (0,-0.5) circle [radius=0.047cm];
            \node at (1.2,0) {$k+p_1+p_2$};
            \node at (-0.75,0) {$k$};
        \end{scope}
        \begin{scope}[scale=1.5]
            \midarrow (1,1) -- (0.5,0);
            \node at (1.1,1.1) {$2$};
            \midarrow (-1,1) -- (-0.5,0);
            \node at (-1.1,1.1) {$1$};
            \node at (-1,0.5) {$p_1$};
            \midarrow (0.5,0) -- (1,-1);
            \node at (1.1,-1.1) {$2'$};
            \midarrow (-0.5,0) -- (-1,-1);
            \node at (-1.1,-1.1) {$1'$};
            \node at (-1,-0.5) {$p_1'$};
            \beforemidarrow (0,0) circle [radius=0.5cm];
            \beforemidarrow (0,0) circle [radius=-0.5cm];
            \draw[fill=black] (0.5,0) circle [radius=0.047cm];
            \draw[fill=black] (-0.5,0) circle [radius=0.047cm];
            \node at (0,0.7) {$k$};
            \node at (0,-0.7) {$k+p_1-p_1'$};
        \end{scope}
        \begin{scope}[scale=1.5, xshift=3.5cm]
            \midarrow (1,1) -- (0.5,0);
            \node at (1.1,1.1) {$2$};
            \midarrow (-1,1) -- (-0.5,0);
            \node at (-1.1,1.1) {$1$};
            \node at (-1,0.5) {$p_1$};
            \midarrow (0.35,-0.425) -- (1.5,-1);
            \draw[thick] (-0.5,0) -- (0.25,-0.375);
            \node at (1.35,-1.1) {$2'$};
            \draw[thick] (0.5,0) -- (0.05,-0.225);
            \draw[thick] (-0.05,-0.275) -- (-0.25,-0.375);
            \midarrow (-0.35,-0.425) -- (-1.5,-1);
            \node at (-1.35,-1.1) {$1'$};
            \node at (1,-0.5) {$p_2'$};
            \beforemidarrow (0,0) circle [radius=0.5cm];
            \beforemidarrow (0,0) circle [radius=-0.5cm];
            \draw[fill=black] (0.5,0) circle [radius=0.047cm];
            \draw[fill=black] (-0.5,0) circle [radius=0.047cm];
            \node at (0,0.7) {$k$};
            \node at (0,-0.7) {$k+p_1-p_2'$};
        \end{scope}
    \etik 
\end{center}

We do not calculate these diagrams in full detail in these notes, but just look at the form they take. Recall the Mandelstam variables 
\be 
\label{eqn:Mandelstam}
    \begin{split}
        s = (p_1+p_2)^2 \\
        t = (p_1'-p_1)^2 \\
        u = (p_2'-p_1)^2
    \end{split}
\ee
We can use the Feynman rules to obtain an integral equation for the diagrams. Let's look at the left-most diagram. The loop momentum $k$ is undetermined\footnote{It flows into both vertices so it appears in two delta functions. This is what we meant in the previous footnote.} and so we must integrate over it. We therefore get
\bse 
    D_{\text{one-loop}} = \frac{(-i\l)^2}{2} \int \frac{d^4k}{(2\pi)^4} \frac{i}{k^2-m^2} \frac{i}{(k+p_1+p_2)^2-m^2}
\ese 
Now note that this is a Lorentz scalar (i.e. it has no Lorentz index), and because we integrate over all $k$, the only thing it can depend on $p_1+p_2$ through the scalar part of it, i.e. $s= (p_1+p_2)^2$.\footnote{For a few more examples/explanation of this see my QED notes.} Therefore 
\be 
\label{eqn:D1Vs}
    D_{\text{one-loop}} = (-i\l)^2 V(s).
\ee
This is nice, but what can we conclude about $V(s)$? Does our integral even converge? Well, at large $k$ it is controlled by 
\bse 
    I = \int \frac{d^4k}{(2\pi)^4} \frac{1}{k^4},
\ese 
where again we integrate over \textit{all} momenta. If we denote the maximum $k$ by $\Lambda$, we see that 
\bse 
    I \sim \log \Lambda \qquad \implies \qquad D_{\text{one-loop}} \sim \l^2 \log\Lambda.
\ese
This is bad: we want to take $\Lambda\to\infty$, and so our integral will diverge! At first we might panic and think "we've broken our QFT!" However, as we are hopefully familiar with from the canonical study, infinities are actually at the very heart of QFT and give us most of the interesting behaviour. Let's look into what this infinity is and where it's coming from. We will make one simplifying assumption first: Assume that we are in a kinematic regime where $s,t,u >> m$. Therefore we can essentially set $m=0$. We do this so that we can focus just on the high momenta divergence and not worry about other complicating subtleties.

Ok so let's carefully calculate the divergent part of $V(s)$. We have just seen this comes directly from $I$. Luckily there is a standard procedure of how to deal with terms like this: 
\ben
    \item \textbf{Wick rotate}: Note that our integral is over \textit{Lorentzian} $4$-momenta. If we wick rotate to \textit{Euclidean} momenta
    \bse 
        k_0 = ik_0^E, \qquad (k^2)_L = (-(k^E)^2)_E,
    \ese 
    with
    \bse 
        d^4 k = i d^4k^E
    \ese 
    then our integral becomes 
    \bse 
        I = i \int \frac{d^4k^E}{(2\pi)^4} \frac{1}{k_E^4}.
    \ese 
    We now note that this is spherically symmetric under SO(4). We can therefore use the $4$D analogue of spherical coordinates:
    \bse 
        \int d^4k^E = \text{Vol}\big(S^3\big) \int_0^{\infty} d|k_E| |k_E|^3 = 2\pi^2 \int_0^{\infty} d|k_E| |k_E|^3,
    \ese 
    where we have used Vol$(S^3) = 2\pi^2$.\footnote{This is easily shown by taking the $4$D coordinates. See, e.g., page 193 of Peskin \& Schroeder.} This turns our integral into 
    \bse 
        I = \frac{i}{8\pi^2}\int_0^{\infty} d|k_E| \frac{1}{|k_E|}.
    \ese 
    \item \textbf{Regulate the integral}: As we just saw, it is the limit $\Lambda\to\infty$ that caused our integral to diverge. We therefore cut off integral at high, but finite, momenta. There are many ways to do this, we will use what is known as a \textit{hard cut off}, i.e. take the max Euclidean momentum $|k_E^{\text{max}}|=\Lambda$ for some finite $\Lambda$.\footnote{This type of regularisation is \textit{not} gauge invariant, and so we could run into problems for gauge theories. See my QED notes for more info on other regularisation schemes.} So the divergent part is
    \be 
    \label{eqn:IDiv}
        I_{\text{div}} = \frac{i}{8\pi^2}\int_0^{\Lambda} d|k_E| \frac{1}{|k_E|} = \frac{i}{8\pi^2} \log \Lambda.
    \ee
\een 

Ok, let's return to thinking about $V(s)$. We know that this is dimensionless, but we have also said that it depends on $s$, which is dimensionful. The only other dimensionful thing ($m=0$, remember) entering $V(s)$ is $\Lambda$. We therefore require
\bse 
    V(s) = f\bigg(\frac{\Lambda^2}{s}\bigg).
\ese
So our problem is now to find the function $f(\Lambda^2/s)$. Well we just found the $\Lambda$ dependence, \Cref{eqn:IDiv}, so we also know the $s$ dependence:
\bse 
    V(s) = -\frac{i}{32\pi^2} \log \bigg(\frac{\Lambda^2}{s}\bigg) + \text{const.},
\ese 
where we have included the $i^2/2$ factor from $D_{\text{one-loop}}$ (see the definition \Cref{eqn:D1Vs}) and an additional $1/2$ in the prefactor to get $\Lambda^2$ in the log. We have also included some constant term which must be $s$-independent. However this term will become negligible in our large $s$ limit, so we will forget about it from now on. 

Note that there was nothing special about the fact that we were considering the $s$ diagram, and indeed the $t$ and $u$ diagrams follow trivially from here. The full set of one loop diagrams, then, give 
\mybox{
    \be 
    \label{eqn:DOneLoop}
        D_{\text{one-loop}} = \frac{i\l^2}{32\pi^2} \bigg[ \log \bigg(\frac{\Lambda^2}{s}\bigg) + \log \bigg(\frac{\Lambda^2}{t}\bigg) + \log \bigg(\frac{\Lambda^2}{u}\bigg)\bigg]. 
    \ee 
}

We have now found the $2\to 2$ scattering to second order, we simply add the above result to our tree-level result:
\be 
\label{eqn:iMOneLoop}
    i\cM(s,t,u; \Lambda) = -i\l + \frac{i\l^2}{32\pi^2} \bigg[ \log \bigg(\frac{\Lambda^2}{s}\bigg) + \log \bigg(\frac{\Lambda^2}{t}\bigg) + \log \bigg(\frac{\Lambda^2}{u}\bigg)\bigg]
\ee 
where we note that the matrix element is obviously dependent on $s,t$ and $u$. This is actually what we expect: the matrix element tells us about scattering probabilities and these are obviously dependent on the momenta in question which are exactly $s,t$ and $u$.

\section{Coming To Terms With Divergences}

We just said that our matrix element depends nicely depends on $s,t$ and $u$ but we also seem to have $\Lambda$ dependence. This might make us panic! Why? Well if we want to be solving our original integral, we are supposed to keep in mind that we want to take the limit $\Lambda\to\infty$. This would give us an infinite probability for our $2\to 2$ scattering to happen! There is no need to panic though, let's calm down a bit and see what's going on here. 

Let's make a few comments of what we've done leading up to here
\ben 
    \item QFT has told us that in order to scatter two $\phi$ particles we need integrate to arbitrarily high momentum, which are arbitrarily small distances. This seems like a weird thing to require: we don't expect to need to use some full quantum gravity theory just to study a simple process like this. Indeed we expect our theory to break down in this limit and so wouldn't necessarily even trust our results even if they didn't diverge. 
    \item In normal physics, we relate one observable quantity to another. For example $pV=NkT$ tells us how to relate the pressure of a gas to the temperature given we know the volume. However our formula above is different. It relates the observable $i\cM$\footnote{Really we mean the probablity which is related to $|i\cM|^2$, but the argument follows through.} to something that is not observable, namely $\l$! It is a parameter appearing in the Lagrangian and these are never directly detectable. We shall refer to this $\l$ as the \textit{bare coupling}. 
\een     
    
So what do we do? Well, as we just said, the above infinity arises when we try relate something unmeasurable to something physical. This infinity is therefore not physical itself and so it looks like everything is nice again provided that whatever the physical thing is doesn't give us some other infinity. Let's try to give $\l$ some kind of physical meaning, then. 

How do we do this? We go to our experimentalist friends, who have done the real life scattering at some energy scale $s_0,t_0,u_0$ and got some number $i\cM(s_0,t_0,u_0)$. We then \textit{define} our physical $\l$, denoted $\l_p$, using this result as
\be 
\label{eqn:LambdaPhysical}
    -i\l_p := i\cM(s_0,t_0,u_0). 
\ee 
This kind of definition is called a \textit{renormalisation condition}.

How does this relate to our bare coupling, $\l$? Well we can simply use our theoretical result, \Cref{eqn:iMOneLoop}:
\bse 
    -i\l_p = -i\l + \frac{i\l^2}{32\pi^2} \bigg[ \log \bigg(\frac{\Lambda^2}{s_0}\bigg) + \log \bigg(\frac{\Lambda^2}{t_0}\bigg) + \log \bigg(\frac{\Lambda^2}{u_0}\bigg)\bigg].
\ese 
We want to solve this for $\l(\l_p)$. Now, as will perhaps be clearer in a moment,\footnote{Wait for \Cref{rem:LambdaLambdaPOrder} and \Cref{eqn:DelLambdaFirstOrder}.} in perturbation theory we treat $\l$ and $\l_p$ to be of the same order, and so if we only work to leading order we are free to replace the $\l^2$ term with $\l_p^2$. We therefore have
\bse 
    \l = \l_p+ \frac{i\l_p^2}{32\pi^2} \bigg[ \log \bigg(\frac{s_0}{\Lambda^2}\bigg) + \log \bigg(\frac{t_0}{\Lambda^2}\bigg) + \log \bigg(\frac{u_0}{\Lambda^2}\bigg)\bigg].
\ese 
We can then plug this into $i\cM(s,t,u)$ to get an expression that relates an observable, $\l_p$, to another observable, $i\cM(t,s,u)$:
\mybox{
    \be 
    \label{eqn:iMOneLoopRenormalised}
        i\cM(s,t,u) = -i\l_p + \frac{i\l_p^2}{32\pi^2} \bigg[ \log \bigg(\frac{s_0}{s}\bigg) + \log \bigg(\frac{t_0}{t}\bigg) + \log \bigg(\frac{u_0}{u}\bigg)\bigg].
    \ee
}
\noindent We note that this is no longer divergent as our $\Lambda$ `dependence' has dropped out. This is why we separated the $\Lambda$ in the argument of \Cref{eqn:iMOneLoop} with a semi-colon. Note that we can study the theory to any momenta (i.e. the arguments on the left-hand side are proper variables not fixed numbers) by comparing it to our reference value $(s_0,t_0,u_0)$.

What we have just done is the infamous \textit{renormalisation}; that is renormalisation is simply the `common sense' idea of setting your theory up relative to some reference point. It is important to note that we did \textit{not} remove our divergence by `subtracting it away' or anything, it was simply never observable so we just insisted we look only at observable things in the first place. Where did our divergence go, then? Well it sits in the relationship between $\l$ and $\l_p$, so essentially it sits in our, physically unmeasurable,\footnote{Are you getting the hint yet?} Lagrangian parameter $\l$. 

As we have just mentioned, there is a small cost to our renormalisation process: we can only do physics relative to our measured reference point $(s_0,t_0,u_0)$. This is essentially a loss of a tiny amount of information. However this tiny loss in information has come at a \textit{huge} benefit: we have obtained the momentum dependence of our scattering amplitude! 

Finally it's worth noting that we are going to have to do this loss of information procedure for every separate infinity we get from the Lagrangian. However if we have the same number of parameters in our Lagrangian as we do infinities of this kind, we will be able to absorb them all into these bare parameters and move on. We obviously require that the infinities that arise at higher order loops corrections (i.e. two loops, three loops, etc.) are essentially related to the original ones. If they are not and they give rise to something completely new, we will have to introduce some new term into our Lagrangian to absorb our infinity. This process can then get wildly out of hand and we can end up having to include new parameters for \textit{every} higher order loop correction. In the full perturbative expansion, then, we need to include an \textit{infinite} number of new terms all of which need to be fixed by experimental measurement. Doing this will obviously \textit{completely} constrain our equation to only work for these specific values and so becomes completely useless for predicting other behaviour. A theory of this kind is, creatively, called \textit{non-renormalisable}. 

\section{Renormalised Perturbation Theory}

What we have done above is give an idea of the concept of renormalisation, but now we actually need to see that it actually \textit{works}, and if it does work if we can always do it. We will answer these two questions in turn. 

\subsection{Counterterms \& Renormalisation Conditions}

First a comment on notation. We will drop the $p$ appearing in the subscript $\l_p$, and simply use $\l$ to denote the physical coupling.\footnote{If you do get confused, I recommend looking at the QED course as there we were explicit about what was the bare parameter and what was the renormalised one.} This hopefully won't cause any confusion as we will only be working with the physical one from now on. We will also use $m^2$ to denote the physical mass, and we shall assume $Z=1$. This collection of conditions is our \textit{renormalisation conditions}, and the latter two conditions correspond to taking the "on-mass-shell" scheme.\footnote{Again see QED for more info on what this means.}

So what do we do? Well essentially what we're saying above is take our original Lagrangian and plug in the substitutions
\bse 
    \l \to \l + \del_{\l}, \qquad m^2 \to m^2 + \del_m, \qand \phi \to \phi\sqrt{1+ \del_{Z}}.
\ese 
This corresponds exactly to what we were doing above when we had $\l = \l_p + (...)$. This is why we no longer need to worry about the subscript $p$. 

\br 
\label{rem:LambdaLambdaPOrder}
    Note that our substitution above corresponds exactly to what we said before about taking $\l$ and $\l_p$ to be the same order in perturbation theory. This will become explicit in \Cref{eqn:DelLambdaFirstOrder}.
\er 

Ok so our $\l\phi^4$ theory becomes
\bse 
    \cL = \frac{1}{2} (\p\phi)^2 -\frac{1}{2}m^2 \phi^2 - \frac{\l}{4!}\phi^4 + \frac{1}{2}\del_Z (\p\phi)^2 - \frac{1}{2}\del_m\phi^2 - \frac{\del_{\l}}{4!} \phi^4
\ese
This might \textit{look} like a new theory, but of course it is exactly the same one just expressed in terms of the physical parameters and some additional terms. These additional terms are called \textit{counter terms} and they are what we will have to adjust order by order in perturbation theory if we want to maintain our renormalisation conditions, that is keep on on-mass-shell conditions.

Now recall that our Feynman diagrams come from the terms in our Lagrangian, and so now we have new Feynman rules corresponding to our counter terms. For this theory we get two new terms, one corresponding to the propagator the other to the vertex. The standard way to denote these is with a $\otimes$ symbol, i.e. our new Feynman rules look like:
\begin{center}
    \btik 
        \begin{scope}[xshift=-3.5cm]
            \draw[thick] (-1,0) -- (-0.25,0);
            \draw[thick] (0,0) circle [radius=0.25];
            \draw[thick, rotate around={45:(0,0)}] (-0.25,0) -- (0.25,0);
            \draw[thick, rotate around={-45:(0,0)}] (-0.25,0) -- (0.25,0);
            \draw[thick] (0.25,0) -- (1,0);
            \node at (2.5,0) {$= i (p^2\del_Z-\del_m)$};
        \end{scope}
        \begin{scope}[xshift=3.5cm]
            \draw[thick,rotate around={45:(0,0)}] (-1,0) -- (-0.25,0);
            \draw[thick, rotate around={-45:(0,0)}] (-1,0) -- (-0.25,0);
            \draw[thick,rotate around={45:(0,0)}] (0.25,0) -- (1,0);
            \draw[thick, rotate around={-45:(0,0)}] (0.25,0) -- (1,0);
            \draw[thick] (0,0) circle [radius=0.25];
            \draw[thick, rotate around={45:(0,0)}] (-0.25,0) -- (0.25,0);
            \draw[thick, rotate around={-45:(0,0)}] (-0.25,0) -- (0.25,0);
            \node at (1.75,0) {$=-i\del_{\l}$};
        \end{scope}
    \etik 
\end{center}

\br 
\label{rem:HigherOrderCorrections}
    Note that the expressions we have obtained for these correction diagrams look like the original expressions. That is the propagator term looks like a propagator "$p^2-m^2$", and the vertex correction looks like a vertex $-i\l$. This is an important result, and it essentially tells us that when we consider higher order corrections we do not need to introduce new diagrams, but we simply alter our correction term values, $\del_Z, \del_m$ and $\del_{\l}$.\footnote{For a slightly more detailed discussion of this, see section 3.3.4 of my QED notes.}
\er 

Ok so we have our new diagrams, we now need to ensure our renormalisation conditions are satisfied. As we've mentioned countless times, we do this by taking some experimental result and using this result as a reference point. How do we translate this into conditions on our propagator and coupling? Well we remember where they come from.

\begin{itemize}
    \item \textbf{Propagator}: The propagator comes from the two point function $\bra{\Omega}\cT[\phi(x)\phi(y)]\ket{\Omega}$, and so our on-mass-shell condition tells us to make the following demand\footnote{Note we have just set $y=0$ on the second term to simplify things a bit. We can do this because $N$-point functions only depend on the difference $(x-y)$.}
    \be
    \label{eqn:PropagatorRenormCondition}
        \int d^4x \, e^{ip\cdot x} \bra{\Omega}\cT[\phi(x)\phi(0)]\ket{\Omega} \overset{!}{=} \frac{i}{p^2-m^2} + (\text{regular at } p^2=m^2).
    \ee 
    What this says is that the pole of our propagator, $p^2=m^2$, gives us the physical mass and that the residue is $1$. These are exactly the two conditions we wanted. If we compare this to our LSZ theorem, \Cref{eqn:LSZMomentum}, we see that the residue = 1 condition is equivalent to requiring $Z=1$. This should be viewed as a normalisation on the field $\phi$, and so $Z$ is often given the name \textit{wave-function normalisation}. 
    \item \textbf{Coupling}: The coupling $\l$ comes from the $4$-point function. We can relate this to the first order amplitude, as there are no other factors contributing to the diagram, which gives us our familiar condition 
    \be
    \label{eqn:CouplingRenormCondition}
        i\cM(s_0,t_0,u_0) \overset{!}{=} -i\l
    \ee 
    for some fixed energy scale $(s_0,t_0,u_0)$. This condition is clearly dependent on \textit{which} $(s_0,t_0,u_0)$ we choose,\footnote{Note, though, that it doesn't matter which one we pick as this is only a \textit{reference} point.} and it is standard to use $(m^2,0,0)$. However remember that in these notes we are working with high energies, and so we can't use this condition. Instead we shall just leave the values simply as $(s_0,t_0,u_0)$ without specifying their values. This will actually be nice because it will allow us to see how these factors carry through in our calculations. 
\end{itemize}

\subsection{Determining The Counterterms}

Now that we have our renormalisation conditions, we can use them to work out what our counter terms are. That is we can use \Cref{eqn:PropagatorRenormCondition,eqn:CouplingRenormCondition} to determine $\del_{\l}, \del_Z$ and $\del_m$. At first we might think this is impossible as we only have 2 equations but three things to find, but we have to remember \Cref{eqn:PropagatorRenormCondition} contains two pieces of information: the position of the pole \textit{and} the residue value. 

\subsubsection{Coupling}

Let's do the coupling first as it is more straight forward. To first order the only two Feynman diagrams we get are 
\begin{center}
    \btik 
        \draw[thick] (-0.5,0.5) -- (0.5,-0.5);
        \draw[thick] (-0.5,-0.5) -- (0.5,0.5);
        \draw[fill=black] (0,0) circle [radius=0.07cm];
        %
        \node at (1.25,0) {$+$};
        % 
        \draw[thick] (2,0.5) -- (3,-0.5);
        \draw[thick] (2,-0.5) -- (3,0.5);
        \draw[thick] (2.5,0) circle [radius=0.25cm];
        %
        \node at (4.5,0) {$= -i\l -i\del_{\l}$};
    \etik 
\end{center}
which, from \Cref{eqn:CouplingRenormCondition}, tells us that we require 
\be
\label{eqn:DelLambdaFirstOrder}
    \del_{\l} = 0 + \cO(\l^2).
\ee 
This is the condition we used before to set $\l^2=\l_p^2$. 

Now let's look at second order. Before we jump in and start drawing diagrams, its important that we note that our expansion is in $\l$ \textit{not} $\del_{\l}$. What we mean by this is we consider diagrams with two $\l$ vertices but we still only have the single $\del_{\l}$ diagram drawn above. This obviously then means we have to alter our result \Cref{eqn:DelLambdaFirstOrder} to compensate for all the diagrams, not just the first order "cross" diagram. The fact that we can do this (i.e. not needing to consider diagrams with two $\del_{\l}$ vertices) is related to the fact that our correction term looks like a vertex again, see \Cref{rem:HigherOrderCorrections}. 

So at second order we have the diagrams 
\begin{center}
    \btik 
        \begin{scope}[xshift=-3.5cm]
            \draw[thick] (1,1) -- (0,0.5);
            \draw[thick] (-1,1) -- (0,0.5);
            \draw[thick] (0,-0.5) -- (1,-1);
            \draw[thick] (0,-0.5) -- (-1,-1);
            \draw[thick] (0,0) circle [radius=0.5cm];
            \draw[fill=black] (0,0.5) circle [radius=0.07cm];
            \draw[fill=black] (0,-0.5) circle [radius=0.07cm];
            \node at (1.75,0) {$+$};
        \end{scope}
        \begin{scope}
            \draw[thick] (1,1) -- (0.5,0);
            \draw[thick] (-1,1) -- (-0.5,0);
            \draw[thick] (0.5,0) -- (1,-1);
            \draw[thick] (-0.5,0) -- (-1,-1);
            \draw[thick] (0,0) circle [radius=-0.5cm];
            \draw[fill=black] (0.5,0) circle [radius=0.07cm];
            \draw[fill=black] (-0.5,0) circle [radius=0.07cm];
            \node at (1.75,0) {$+$};
        \end{scope}
        \begin{scope}[xshift=3.5cm]
            \draw[thick] (1,1) -- (0.5,0);
            \draw[thick] (-1,1) -- (-0.5,0);
            \draw[thick] (0.35,-0.425) -- (1.5,-1);
            \draw[thick] (-0.5,0) -- (0.25,-0.375);
            \draw[thick] (0.5,0) -- (0.05,-0.225);
            \draw[thick] (-0.05,-0.275) -- (-0.25,-0.375);
            \draw[thick] (-0.35,-0.425) -- (-1.5,-1);
            \draw[thick] (0,0) circle [radius=-0.5cm];
            \draw[fill=black] (0.5,0) circle [radius=0.07cm];
            \draw[fill=black] (-0.5,0) circle [radius=0.07cm];
            \node at (1.75,0) {$+$};
        \end{scope}
        \begin{scope}[xshift=7cm]
            \draw[thick] (1,1) -- (-1,-1);
            \draw[thick] (1,-1) -- (-1,1);
            \draw[thick] (0,0) circle [radius=0.25cm];
        \end{scope}
    \etik
\end{center}
which gives us 
\bse 
    i\cM(s,t,u) = -i\l + \frac{\l^2}{32\pi^2}\bigg[ \log(\frac{\Lambda^2}{s}) + \log(\frac{\Lambda^2}{t}) + \log(\frac{\Lambda^2}{u})\bigg] -i\del_{\l},
\ese 
and then \Cref{eqn:CouplingRenormCondition} gives us 
\bse 
    i\del_{\l} = i\frac{\l^2}{32\pi^2}\bigg[ \log(\frac{\Lambda^2}{s_0}) + \log(\frac{\Lambda^2}{t_0}) + \log(\frac{\Lambda^2}{u_0})\bigg] + \cO(\l^3).
\ese 
This then gives us the same result we obtained before
\bse 
    i\cM = -i\l + \frac{\l^2}{32\pi^2}\bigg[ \log(\frac{s_0}{s}) + \log(\frac{t_0}{t}) + \log(\frac{u_0}{u})\bigg].
\ese

We have bothered to obtain this result this way for two reasons: firstly it gives further support for our result and secondly it allows us to see how we proceed for higher orders. Note in particular that $\del_{\l}$ is always going to be to the same order as the truncation of the perturbation, i.e. at third order we expect $\del_{\l} \sim \l^3$ etc. We therefore see that our renormalisation scheme does indeed work to all orders, we just have to continuously keep adjusting the value of $\del_{\l}$. 

\br 
    Again the fact that this works to all orders is a direct repercussion of the fact that our correction term takes the same form as our original coupling. If it had not we would have to introduce a new term into the Lagrangian at every order and we would end up needing an infinite number of them for the full perturbation expansion, and this would kill our ability to predict anything. 
\er 

\subsubsection{Propagator}

Ok that's the coupling done, now let's do the propagator. This will take a little bit more time and involves us introducing what are known as \textit{one-particle irreducible} (1PI) diagrams. 

Recall that our on-mass-shell renormalisation condition required that the pole of our theory gives us the physical mass of the particle. We therefore need to keep track of how the pole shifts, and in order to do this will need sum an infinite set of diagrams. To gain more clarity on what we mean, let's look at the structure of the two-point function. 

At order $\l^0$ we of course just have the exact result 
\bse 
    \int d^4 x \, e^{ip\cdot x} \bra{\Omega}\cT[\phi(x)\phi(0)]\ket{\Omega} = D_F(p).
\ese 
This is what we expect as the free propagator is non-divergent. However at order $\l$ we get the following diagram 
\begin{center}
    \btik 
        \midarrow (-1,0) -- (0,0);
        \node at (-0.5,-0.3) {$p$};
        \draw[thick] (0,0) -- (1,0);
        \draw[thick] (0,0) .. controls (-1.5,1.5) and (1.5,1.5) .. (0,0);
        \draw[fill=black] (0,0) circle [radius=0.07cm];
    \etik  
\end{center}
which is divergent and needs to be counteracted with our corrected propagator
\begin{center}
    \btik 
        \midarrow (-1,0) -- (-0.25,0);
        \node at (-0.625,-0.3) {$p$};
        \draw[thick] (0,0) circle [radius=0.25];
        \draw[thick, rotate around={45:(0,0)}] (-0.25,0) -- (0.25,0);
        \draw[thick, rotate around={-45:(0,0)}] (-0.25,0) -- (0.25,0);
        \draw[thick] (0.25,0) -- (1,0);
    \etik 
\end{center}
where in both diagrams we have labelled the incoming momentum. 

What about second order in $\l$? Well there's two diagrams we can get, namely:
\begin{center}
    \btik 
        \begin{scope}[xshift=-3cm]
            \midarrow (-1,0) -- (0,0);
            \node at (-0.5,-0.3) {$p$};
            \draw[thick] (0,0) -- (1,0);
            \draw[thick] (0,0) .. controls (-1.5,1.5) and (1.5,1.5) .. (0,0);
            \draw[fill=black] (0,0) circle [radius=0.07cm];
            \draw[thick] (1,0) -- (1.5,0);
            \draw[thick] (1.5,0) -- (2.5,0);
            \draw[thick] (1.5,0) .. controls (0,1.5) and (3,1.5) .. (1.5,0);
            \draw[fill=black] (1.5,0) circle [radius=0.07cm];
        \end{scope}
        \begin{scope}[xshift=3cm]
            \midarrow (-1,0) -- (0,0);
            \node at (-0.5,-0.3) {$p$};
            \draw[thick] (0,0) -- (1,0);
            \draw[thick] (0,0) .. controls (-1.5,1.5) and (1.5,1.5) .. (0,0);
            \draw[fill=black] (0,0) circle [radius=0.07cm];
            \draw[thick] (0,1.14) .. controls (-1.5,2.64) and (1.5,2.64) .. (0,1.14);
            \draw[fill=black] (0,1.14) circle [radius=0.07cm];
        \end{scope}
    \etik 
\end{center}
Now the second one looks new (and perhaps a bit scary) but the first one just looks like two copies of our first order diagram "stuck together". Indeed it is exactly what it is, and the continuation of this argument motives the definition of our 1PIs. 

\bd[One-Particle Irreducible]
    If a diagram can be cut along a propagator line and give to separate diagrams of lower order it is \textit{not} 1PI. If we cannot do this it is 1PI. We denote the full set of 1PI diagrams as $-i\Sigma(p)$. 
\ed 

\noindent For clarity on the above definition, the right second order diagram is 1PI but the left one is not. It is also important to note that our corrected propagator (i.e. the one with the $\otimes$ on it) is counted in our 1PI definition. 

\bbox 
    There are three\footnote{Ignoring the $\otimes$ diagrams.} 1PI diagrams for $\phi^4$ theory at order $\l^3$. Draw them. 
    
    \textit{Hint: I'm just going to leave this here...}
    \begin{center}
        \includegraphics[scale=0.2]{images/MickeyMouse.png}
    \end{center}
    \textit{Also the word "sunset" might help...}
\ebox  

Our full propagator term is then given by the sum 
\begin{center}
    \btik 
        \midarrow (-0.5,0) -- (0.5,0);
        %
        \node at (1,0) {$+$};
        %
        \midarrow (1.5,0) -- (2,0);
        \draw[thick] (2.5,0) circle [radius=0.5cm];
        \node at (2.5,0) {1PI};
        \midarrow (3,0) -- (3.5,0);
        % 
        \node at (4,0) {$+$};
        % 
        \midarrow (4.5,0) -- (5,0);
        \draw[thick] (5.5,0) circle [radius=0.5cm];
        \node at (5.5,0) {1PI};
        \midarrow (6,0) -- (7,0);
        \draw[thick] (7.5,0) circle [radius=0.5cm];
        \node at (7.5,0) {1PI};
        \midarrow (8,0) -- (8.5,0);
        %
        \node at (9,0) {$+$};
        \node at (9.5,0) {$...$};
    \etik 
\end{center}
which mathematically is 
\bse 
    D_F(p) + D_F(p)\big(-i\Sigma(p)\big) D_F(p) + ... = D_F(p)\sum_{n=0}^{\infty}\Big[\big(-i\Sigma(p)\big) D_F(p)\Big]^n,
\ese
which is just a geometric sum,\footnote{Note really we should be a bit more careful because for Fermions the propagator is a matrix valued expression. This is discussed a bit more in my QED notes.} and so we get 
\bse 
    D_F(p) \frac{1}{1+ i\Sigma(p)D_F(p)} = \frac{i}{p^2-m^2-\Sigma(p)},
\ese 
where we have used the result $D_F(p) = (p^2-m^2)^{-1}$. This result is really worth stressing and so we put it again in a nice box:
\mybox{
\be 
\label{eqn:SF(p)}
    S_F(p) = \frac{i}{p^2-m^2-\Sigma(p)}
\ee 
}
\noindent where we have just introduced the common notational definition $S_F(p)$. 

Ok so what we have just shown is that $\Sigma(p)$ is exactly the shift in our pole, and so if we are to get our on-mass-shell renormalisation condition we require 
\mybox{
\be 
\label{eqn:SigmaOnShell}
    \Sigma(p^2=m^2) = 0, \qand \bigg(\frac{d \Sigma}{dp^2}\bigg)\bigg|_{p^2=m^2}=0.
\ee 
}
\noindent where the second result gives us our residue condition.\footnote{The proof of this was an exercise on the course, so I don't want to type the answer here. If you can't see it at all, feel free to email me and I'll give you some hint.}

Let's now evaluate this to order $\l$. From the diagrams above, we have\footnote{Note the factor of $1/2$. It is a symmetry factor.}
\bse 
    -i\Sigma(p^2) = -\frac{i\l}{2} \int \frac{d^4 k}{(2\pi)^4} \frac{i}{k^2-m^2} + i(p^2\del_Z -\del_m),
\ese
where $k$ is the momentum in the loop. Now we note that the loop part is completely independent of $p^2$ and so is the $\del_m$ term, so our residue condition (the derivative one) simply gives us $\del_Z=0$ and then the pole condition gives us our $\del_m$ result. We summarise these again in a nice box:
\mybox{
\be 
\label{eqn:PropagatorCorrectionFirstOrder}
    \del_Z = 0, \qand \del_m = - \frac{\l}{2} \int \frac{d^4 k}{(2\pi)^4} \frac{i}{k^2-m^2}.
\ee 
}
\noindent We could then evaluate the integral to get an explicit expression for $\del_m$ but we do not do that here, as the result depends on our cut-off value $\Lambda$. 

\br 
    Note the fact that we got $\del_Z=0$ above is purely a repercussion of the fact we are only working to order $\l$. If we worked to higher orders, our diagrams will have non-trivial dependence on $p^2$ and so we will get something non-zero for $\del_Z$.
\er 

\section{What Theories Are Renormalisable \& What Does This Mean?}

Ok so we have successfully (at least outlined how to) renormalised our $\phi^4$ theory to all orders. The obvious question is "can we always do this? If not what are the conditions we need to ensure that we can?" The answer to the first question is "no", and the answer to the second question takes a bit of work. 

Let's recap what we did above: essentially we drew the diagrams and found the ones that were divergent and then introduced counter terms for these diagrams and then constrained the result using an experiment. What we need, therefore, is some way to determine if a diagram is divergent and how this divergence will be effected as we edit the diagram by, say, adding another vertex.

We will do this by generalising our $\phi^4$ theory in $4$-dimensions to $\phi^n$ theory in $d$-dimensions. Our action is then
\bse 
    S[\phi] = \int d^d x \, \bigg(\frac{1}{2}(\p\phi)^2 - \frac{m^2}{2}\phi^2 - \frac{\l}{n!} \phi^n\bigg).
\ese 
Now we invoke some dimensional analysis arguments. Recall that in QFT we work in so-called \textit{natural units} (i.e. $\hbar=c=1$) in which case we can categorise the dimension of everything by simply stating its mass dimension. The most important one to note here is $[x]=-1$, from which it follows that $[\p]=+1$ and $[dx]=-1$.\footnote{If this notation or argument doesn't make sense to you, see your canonical field theory notes.} 

Now recall that the action appears inside an exponential in our path integrals. This tells us\footnote{As the Taylor expansion of an exponential contains powers of the argument, and if we are to add these things they must all have the same dimension. The only way we can satisfy this is if it is dimensionless.} we \textit{must} have $[S]=0$. Now our integral measure is going to give us a dimension factor of $[d^dx]=-d$, and so we must cancel all of these with the Lagrangian, i.e. $[\cL]=+d$. We can then work from here to find the dimension of our field, and from there the dimension of our coupling constant. 
\bbox 
    Convince yourself that 
    \bse 
        [\phi]= \frac{d-2}{2} \qquad \implies \qquad [\l] = d-n\bigg(\frac{d-2}{2}\bigg)
    \ese
    for our theory. 
\ebox 

Why are we talking about this? Well because it allows us to define what is known as the \textit{superficial degree of divergence} (SODD), $D$. This is essentially how badly a diagram \textit{looks} like its diverging. That is, if the SDOD is $D$ then the diagram diverges as $\Lambda^D$.

\br
\label{rem:SDOD}
    We say "looks" because it turns out that the SDOD is not always correct: it can both underestimate and overestimate the actual divergence of the diagram. For example symmetry factors can cause two diagrams to partially cancel, thereby lowering the SDOD. For a bit more information on this see my QED notes.
\er 

With the above remark in mind, we are going to pretend that the SDOD is a good measure of the actual divergence and plough on. We want to relate $D$ to things in the diagram, but how do we do this? Well recall that our divergences are coming from the integrals and in particular by the ratio of the number of dimensions we're integrating over to the momentum powers in the denominator. This is why we were get $\log$ divergences: we were integrating $1/p^4$ with measure $d^4p$. 

Next recall that each loop comes with an undetermined momentum and therefore the Feynman rules tell us to include a $d$-dimensional integral. This is going to give us a mass dimension of $+d$. Similarly every propagator comes with a $1/(p^2-m^2)$, and in the high energy limit (which is where the divergences we are talking about come from) this simply becomes $1/p^2$. Momentum has units of energy which has units of inverse length and so each propagator gives us a factor of $-2$. We therefore get 
\bse 
    D = dL - 2P,
\ese    
where $L/P$ are the number of loops/propagators in the diagram, respectively. This is hopefully an intuitively clear result: the more loops we have the worse our divergence and the more propagators we have the better. 

Ok this is nice, but is there a way for us to express this result in terms of something easier to see, say the parameters appearing in the Lagrangian? The answer, of course,\footnote{Otherwise I wouldn't be typing this.} is "yes", let's see how. 

Imagine we have a diagram with exactly $N$ external lines (i.e. both incoming \textit{and} outgoing particles) and $V$ vertices. Now \textit{imagine} that our Lagrangian contained a term of the form $\eta\phi^N$ then we would be able to draw this diagram with exactly one vertex. This vertex will give a factor of $\eta$, and so, from our previous discussion, we see that the dimension of the diagram is just 
\be
\label{eqn:DiagramDimension}
    [\text{Diag}] = [\eta] = d - N\bigg(\frac{d-2}{2}\bigg).
\ee 
We stressed the word "imagine" above as it is possible, and indeed probable, that our theory won't contain exactly this term. We just use this argument to obtain the dimensions for the diagram. If our theory might not contain this interaction vertex, why are we bothering to discuss it? Well, if we are going to sum the diagrams, which we always do to get the full amplitude, each diagram \textit{must} have the same dimension.\footnote{By the same argument made in footnote 16.} We can then use this as the \textit{definition} of the dimension of our diagram and compare it to one that we can get for our actual theory. 

In our actual theory we construct the diagrams by `gluing together' a bunch of propagators and vertices. These vertices will each give a dimensionful contribution of $[\l]$ to the dimension of the diagram. Therefore if the diagram has $V$ vertices, the most divergent part\footnote{Note any other contributions to the diagram will only make it \textit{less} divergent, so here we're considering the worst case scenario.} of the diagram has dimension
\bse 
    V[\l] + \big[\Lambda^D\big] = V\bigg[ d - n\bigg(\frac{d-2}{2}\bigg)\bigg] + D
\ese 
where the second term comes from the definition of the SDOD. We can then equate this \Cref{eqn:DiagramDimension} and rearrange for $D$ to obtain
\mybox{
\be 
\label{eqn:SDOD}
    D = d + V\bigg[n\bigg(\frac{d-2}{2}\bigg) - d\bigg] - N\bigg(\frac{d-2}{2}\bigg). 
\ee 
}

For the purposes of the next comment, let's rewrite the above as 
\bse 
    D = d - V[\l] - N\bigg(\frac{d-2}{2}\bigg).
\ese 
Now the last term is always negative,\footnote{Provided $d>1$. As Dr. Iqbal points out the QM case, namely $d=1$, seems to be different. I have no idea why.} and the first term is obviously always positive. Let's categorise this result by the values of $[\l]$, then. 
\begin{itemize}
    \item $[\l]>0$, \textbf{Super-Renormalisable}: In this case the more vertices we include the less divergent the diagram is. This tells us that we only have a finite number of divergent \textit{diagrams} and so we can obviously renormalise the theory. Indeed there becomes a point when we no longer need to add corrections to our counter terms at all. That is once we hit the $V$ value that renders $D\leq 0$, \textit{all} higher order diagrams are convergent given our counter terms. This is why we call them \textit{super}-renormalisable, in contrast to the next case. 
    \item $[\l]=0$, \textbf{Renormalisable}: In this case adding additional vertices doesn't change anything because that term vanishes. We therefore see the only way to ensure a convergent diagram is to increase the number of external states, $N$. However \textit{all} diagrams with $N$ less then this critical value will be divergent. That is if a diagram is divergent at one-loop order it will continue to be divergent at all higher order loops. We therefore get an infinite number of divergent diagrams, but we have a \textit{finite} number of divergent \textit{amplitudes}, i.e. all scattering processes with enough external states are convergent. This amounts to the fact that we only need a finite number of experiments (as there are only finitely many divergent amplitudes) but we will have to continuously alter our correction terms to higher and higher orders. This was the case for our $\phi^4$ in $d=4$ theory, which is in agreement with $[\l]=0$ there.
    \item $[\l]<0$, \textbf{Non-Renormalisable}: In this case the more vertices we include the more external states we need to keep our amplitudes from diverging. This tells us that not only do we have an infinite number of divergent diagrams but we also have an infinite number of divergent \textit{amplitudes}. This means we would need more and more experiments at higher order loops to fix our counter terms and we would end up needing an infinite number. This would completely constrain our reference point and we would not be able to predict any other physics, rendering the theory useless.
\end{itemize}

We have therefore come up with a method for just "seeing" if the theory is divergent by simply finding the dimension of the coupling constants appearing in the Lagrangian. The method we have described is known as \textit{power-counting renormalisation}, for obvious reasons.

\br 
    As we mentioned before, the SDOD only gives us what the divergence of a diagram "looks" like. We gave an example in \Cref{rem:SDOD} of how it can overestimate the divergence of a diagram, but the more worrying case would be it underestimating the divergence of the diagram. When could this happen? Well if we had a badly divergent subdiagram, as in, for example, 
    \begin{center}
        \btik
            \draw[thick] (1,1) -- (0.5,0);
            \draw[thick] (-1,1) -- (-0.5,0);
            \draw[thick] (0.5,0) -- (1,-1);
            \draw[thick] (-0.5,0) -- (-1,-1);
            \draw[thick] (0,0) circle [radius=0.5cm];
            \draw[thick] (0,0.5) .. controls (-1,1.5) and (1,1.5) .. (0,0.5);
            \draw[fill=black] (0.5,0) circle [radius=0.07cm];
            \draw[fill=black] (-0.5,0) circle [radius=0.07cm];
            \draw[fill=black] (0,0.5) circle [radius=0.07cm];
        \etik 
    \end{center}
    It turns out\footnote{See section 3.3.4 of my QED notes, for a bit more detail.} that this will not actually effect our discussion above as we can essentially `shrink' the loop in the above diagram and replace it by our first order vertex correction. That is we can replace this diagram with 
    \begin{center}
        \btik 
            \draw[thick] (1,1) -- (0.5,0);
            \draw[thick] (-1,1) -- (-0.5,0);
            \draw[thick] (0.5,0) -- (1,-1);
            \draw[thick] (-0.5,0) -- (-1,-1);
            \draw[thick] (0,0) circle [radius=0.5cm];
            \draw[fill=black] (0.5,0) circle [radius=0.07cm];
            \draw[fill=black] (-0.5,0) circle [radius=0.07cm];
            \draw[fill=white] (0,0.5) circle [radius=0.25cm];
            \draw[thick, rotate around={45:(0,0.5)}] (-0.25,0.5) -- (0.25,0.5);
            \draw[thick, rotate around={-45:(0,0.5)}] (-0.25,0.5) -- (0.25,0.5);
        \etik 
    \end{center}
    and proceed as above.
\er 

\section{A Few Non-Renormalisable Theories}

Let's now give some more physical significance to our somewhat abstract results above about what it means for a theory to be non-renormalisable. We shall explain the idea here and then give real world examples. 

Our result above was based around the dimension of the coupling, but what does this actually mean? The answer comes from recalling that the matrix element $i\cM$ essentially tells us the probability for a particular scattering to occur;\footnote{Well technically its the modulus squared that gives us the probability, but the idea is the same. Another way you could argue this result would be the fact that, as we have seen, $i\cM$ is a complex number and so cannot possibly have physical dimensions.} it is therefore dimensionless. However we have seen explicitly that the coupling parameters appear in our matrix elements, and so if they are dimensionful we must include some other dimensionful quantity to cancel them out. 

To see why this causes us problems, lets consider $\phi^6$ theory in $d=4$ dimensions. By our arguments above this theory is non-renormalisable. The coupling constant for this term will have dimension $-2$, and so we can write the term in the Lagrangian as
\bse 
    \cL_6 = \int d^4 x \, \frac{c}{M^2} \phi^6,
\ese 
where $[M]=+1$ and $c$ is some dimensionless constant. Now imagine we have a diagram that just contains a single one of these vertices then we expect 
\bse 
    i\cM_6 \sim \frac{c}{M^2},
\ese 
and we need to include something to cancel the dimensions. Obviously we can't just stick anything in there, and with a bit of thought we see that the \textit{only} thing we can use is the physical energy of the scattering, $E$. Recalling the $[E]=+1$, we see that what we require is 
\bse 
    i\cM_6 \sim \frac{c}{M^2}E^2.
\ese
Why is this a problem? Well recall our divergences were coming from the high momentum (and therefore high energy) limit. It follows from the above, then, that the theory becomes more and more strongly coupled at higher energies, and there becomes a point at which our perturbation theory breaks down. That is, our theory is no good at predicting physics at high energies/small distances. 

So what do we do? Well we don't really have a choice, we cannot take our cut off value $\Lambda$ to infinity, and so we \textit{must} carry it through in our calculation and make sure that we are only asking questions valid in the energy limit $E << \Lambda$. We then just hope that this restriction isn't going to effect our calculations too badly. The idea we just sketched out is the basic concept of \textit{effective field theory}, and apparently we will discuss this in more detail in the Renormalisation Group course in second term.\footnote{Notes to come... Provided I don't crash and fail the exams...} For now let's just look at two examples of non-renormalisable theories.

\subsection{Pions}

Pions are a group of mesons, i.e. quark-antiquark pair. There are three of them known, denoted by $\pi^+,\pi^-$ and $\pi^0$, where the superscript indicates their electric charge. The pions are massive, and their masses are all very close, being approximately 135MeV. If we work in a limit where quarks are massless, though, we can express the action as 
\bse 
    S = \int d^4 x\, \bigg(\frac{1}{2}\p_{\mu}\Vec{\pi}\cdot \p^{\mu}\Vec{\pi} + \frac{1}{2F_{\pi}^2}(\p_{\mu}\Vec{\pi}\cdot \vec{\pi})(\p^{\mu}\Vec{\pi}\cdot\vec{\pi}) + ... \bigg)
\ese 
where $\vec{\pi}$ is a $3$ vector whose entries are the three pions and $F_{\pi}$ is the so-called \textit{pion decay constant}. The kinetic term in this action tells us that $[\vec{\pi}]=+1$ and so we it follows that $[F_{\pi}]=+1$. As it appears in the denominator, then, this theory is non-renormalisable. This theory therefore breaks down at energies $E\sim F_{\pi}$ (which has been measured to be around 93MeV), and we must replace it with something else. 

It turns out that we do in fact know what we need to replace this theory with, its QCD, which \textit{is} renormalisable. What was going wrong with our pion theory? Well at high energies it implies that the pions will `split' into their quark antiquark pairs, whereas we know from QCD that this is not possible. That is QCD tells us that quark states are \textit{not} asymptotically free, but we must \textit{always} observe quarks in at least twos. 

Let's contrast this theory with the following example. 

\subsection{Gravity}

You are probably aware that it is a big current goal to come up with a full QFT for gravity. It is often said that "gravity and quantum field theory don't get along", well, as we will see here, that is not quite the true statement. 

Recall from GR that the action for gravity is the so-called \textit{Einstein-Hilbert action}
\bse 
    S = \frac{1}{16\pi G_N} \int d^4x\, \sqrt{-g} R,
\ese 
where $R$ is the \textit{Ricci scalar}. We can boldly attempt to treat this as a QFT by expanding around flat spacetime:\footnote{If you think this seems like a strange thing to do, this is actually where the study of gravitational waves comes from. These are obviously now experimentally verified objects, and so we are not totally crazy for doing something like this.} 
\bse 
    g_{\mu\nu} = \eta_{\mu\nu} + h_{\mu\nu},
\ese 
where $\eta_{\mu\nu}$ is the Minkowski metric and $h_{\mu\nu}$ is our small perturbation parameter, which we treat as a quantum field. If we plug this bold expansion into our Einstein-Hilbert action and use the behaviour $R \sim (\p h)^2$, our action becomes 
\bse 
    S = \frac{1}{16\pi G_N}\int d^4x (\p h)^2 \big(1 + h^2 + h^4 + ...\big),
\ese 
where we have of course dropped all the index behaviour as we are only really interested in the power behaviour. This result tells us that $h$ must be dimensionless (as otherwise $h^2+h^4$ wouldn't make sense) and so it follows (note the derivatives) that $[G_N]=-2$. This defines a dimensionful scale called the \textit{Planck scale}
\bse 
    m_{Pl}^2 := \frac{1}{16\pi G_N} \sim 10^{19}GeV.
\ese 

We can make our action look more like a normal QFT (i.e. put the dimensionful stuff in the expansion) by redefining our perturbation field to be 
\bse 
    \widetilde{h} := m_{Pl}^2 h,
\ese 
giving us 
\bse 
    S = \int d^4 x\, \big(\p\widetilde{h}\big)^2 \bigg( 1 + \widetilde{h}^2 + \frac{1}{m_{Pl}^2}\widetilde{h}^4 + ... \bigg).
\ese 
We see from here straight away that this theory is non-renormalisable, and so it must break down as high energies. This is what people mean by "gravity and QFT don't get along". Note, though, that this theory \textit{does} work provided we stick to low enough energies. This therefore corrects the statement to "gravity and QFT don't get along \textit{at arbitrarily high energies}". 

The difference to the pion case, though, is that we do not as of yet know what theory we need to replace this one. The main problem in trying to find such a theory is that the energy scale we're talking about is \textit{way} beyond our current experimental capacity and so we have no way to test the theories. The most famous theory we have so far is string theory, and this is why so many people are interested in studying it. 
\chapter{Global Symmetries In The Functional Formalism}

\mybox{
    Just a reminder/disclaimer: This chapter was not part of the taught material, however I have worked from Dr. Iqbal's notes to fill it in here.
}

As we know from our other field theory courses, symmetries play a huge role in QFT and (the continuous symmetries) fall into two general categories: global and local/gauge. The former are simpler and so we consider them first here. We will return to gauge theories later and obtain some nice results. 

\section{Classical Noether's Theorem}

Recall that, classically, global continuous symmetries give rise to conserved currents, known as \textit{Noether currents}. This is the content of Noether's theorem. These conserved currents essentially trickle down to becoming the different charges our particles carry in the QFT. It is therefore important that we know how to use Noether's theorem in the path integral approach. The derivation we provide here will be slightly different to the one normally seen in a canonical field theory course, but is a standard approach in the path integral language. The reason for this change in approach will become apparent soon. 

Let's consider a system with Lagrangian $\cL$. We purposefully do not specify the exact form of the Lagrangian so that we can see the generality of the results presented. Let's denote the field content of this theory by the set $\{\phi^a\}$, where each $a$ value corresponds to a different field. That is our action is $S[\{\phi^a\}]$.

Now let's suppose that the system is invariant under some global symmetry transformation. As the global transformation is continuous, we can write it infinitesimally. That is the field transforms as 
\be
\label{eqn:GlobalFieldTransformation}
    \phi^a(x) \to \phi^a(x) + \epsilon \del \phi(x),
\ee 
where $\epsilon$ is some small \textit{constant} parameter. Our requirement that this be a symmetry of the system is then just the statement
\bse 
    S[\phi^a + \epsilon \del\phi^a] = S[\phi^a].
\ese 

Now as we are considering a global symmetry we have that $\epsilon$ is a constant. However in order to arrive at Noether's theorem we perform a trick that at first is confusing: we \textit{imagine} that $\epsilon$ was actually a function of the spacetime events, i.e. $\epsilon = \epsilon(x)$. It is very important to note that \textit{physically} we are not saying that our \textit{symmetry} is dependent on the spacetime points, as this would be a gauge symmetry. What we are simply doing is using a trick to find the form of how the action changes. 

As we just tried to clarify, we know that the action is invariant when $\epsilon$ is just a constant, but in the imagined situation there is no reason to assume this would still be true. We therefore conclude that the change in the action must be related to the \textit{derivative} of $\epsilon(x)$. The action id given by the integral of the Lagrangian (density), which is itself Lorentz invariant, and so we conclude 
\bse 
    \del_{\epsilon}S[\phi^a] = \int d^4 x \, j^{\mu}\p_{\mu}\epsilon(x),
\ese 
where $j^{\mu}$ is whatever else appears with $\p_{\mu}\epsilon(x)$ in the Lagrangian. Now we use the usual trick of integrating by parts and making the physicists typical argument of "boundary terms don't matter" we get
\be 
\label{eqn:VariationOfActionEpsilon}
    \del_{\epsilon} S[\phi^a] = - \int d^4x (\p_{\mu}j^{\mu})\epsilon(x).
\ee 

Now as we said above, there is absolutely no reason to assume that \Cref{eqn:VariationOfActionEpsilon} vanishes itself, it was only when $\epsilon$ was a constant that it vanished. So what do we do? Well we know that when the fields $\phi^a$ are \textit{on-shell}, that is they obey the equations of motion, that the action \textit{must} be minimised (as this is where the equations of motion come from) and so we have 
\be 
\label{eqn:VariationActionField}
    \del_{\phi} S[\phi^a] = 0.
\ee 
This is true for \textit{any} variation of the field, and so it must hold for the particular case of \Cref{eqn:GlobalFieldTransformation}, but this lead exactly to \Cref{eqn:VariationOfActionEpsilon}. We can therefore conclude that 
\mybox{
\be 
\label{eqn:ConservedCurrent}
    \p_{\mu}j^{\mu} = 0.
\ee 
}

\br 
    It is important to note that in order to derive \Cref{eqn:ConservedCurrent} we have used \textit{both} the fact that the action is invariant (to obtain \Cref{eqn:VariationOfActionEpsilon}) \textit{and} the fact the fields are on-shell (to use \Cref{eqn:VariationActionField}). 
\er 

As the above remark highlights, what we have shown is that classically we get a conserved current when we have a global symmetry and the fields obey the classical equations of motion. 

\subsection{Quantum Ward Identities}

As we have tried to stress above, the result \Cref{eqn:ConservedCurrent} is a \textit{classical} result. However this is a course on quantum field theory and so we obviously want to see what the quantum analogue is. 

As is hopefully natural by now, our starting point is the partition function 
\bse 
    Z = \int [\pD \phi^a] e^{iS[\phi^a]}.
\ese 
We now consider our global transformation \Cref{eqn:GlobalFieldTransformation}, with $\epsilon=\epsilon(x)$, and see what happens. Well, as far as the path integral measure is concerned, this is just a change of variables\footnote{The tilde here is not a Fourier transform, just meant to indicate the transformation.}
\bse 
    \phi^a \to \widetilde{\phi}^a = \phi^a(x) + \epsilon(x) \del \phi(x).
\ese 
We therefore have 
\bse 
    \int [\pD \phi^a] e^{iS[\phi^a]} = \int [\pD \widetilde{\phi}^a] e^{iS[\widetilde{\phi}^a]}.
\ese

Now, as the transformation we are looking at is a symmetry, the claim is that the measure itself is invariant,\footnote{Apparently this is not always true, but exceptions happen when so-called \textit{anomalies} enter the game. I do not yet know what these are, so do not wish to comment further. Dr. Iqbal gives Chapter 19 of Peskin as a reference.} namely 
\bse 
    [\pD \widetilde{\phi}^a] = [\pD \phi^a].
\ese 
Putting this together with \Cref{eqn:VariationOfActionEpsilon}, we get 
\bse 
    \int [\pD \phi^a] e^{iS[\phi^a]} = \int [\pD \phi^a] \bigg(1 - i\int d^4x (\p_{\mu}j^{\mu}) \epsilon(x) + \cO(\epsilon^2) \bigg) e^{iS[\phi^a]},
\ese 
where we have Taylor expanded the $e^{i\del_{\epsilon}S[\phi^a]}$ part. From here we can conclude that 
\mybox{
\be 
    \bra{0} \p_{\mu} j^{\mu}(x) \ket{0} \equiv \frac{1}{Z[0]}\int [\pD \phi^a] \p_{\mu}j^{\mu}(x) = 0.
\ee 
}

Note that this is \textit{not} the same thing as saying $\p_{\mu}j^{\mu}=0$. To see why this is the case, let's consider the variation of the 2-point correlation function 
\bse 
    \begin{split}
        \bra{0}\cT[ \phi^b(x_1) \phi^c(x_2) ] \ket{0} & = \frac{1}{Z[0]} \int [\pD \phi^a] \phi^b(x_1) \phi^c(x_2) e^{iS[\phi^a]} \\
        & = \frac{1}{Z[0]} \int [\pD \widetilde{\phi}^a] \widetilde{\phi}^b(x_1) \widetilde{\phi}^c(x_2) e^{iS[\widetilde{\phi}^a]}
    \end{split}
\ese 
where the second line follows from the same "change of variables" argument from above. If we then also again make the claim that the measure is invariant, we can write the last line as
\bse 
    \frac{1}{Z[0]} \int [\pD \phi^a] \big(\phi^b(x_1) + \epsilon(x_1) \del \phi^b(x_1)\big) \big(\phi^c(x_2) + \epsilon(x_2) \del \phi^c(x_2)\big) \bigg(1 -i\int d^4x (\p_{\mu}j^{\mu}) \epsilon(x) + \cO(\epsilon^2)\bigg) e^{iS[\phi^a]}.
\ese 
As before we see that the order $\epsilon$ terms must cancel and so we get
\bse 
    0 =  \frac{1}{Z[0]}\int [\pD\phi^a] \bigg(\epsilon(x_1)\del\phi^b(x_1)\phi^c(x_2) + \epsilon(x_2)\phi^b(x_1)\del\phi^c(x_2) - i \phi^b(x_1)\phi^c(x_2)\int d^4x (\p_{\mu}j^{\mu})\epsilon(x)\bigg) e^{iS[\phi^a]}.
\ese 
Now we want to write this in the form $\epsilon(x) \times (...)$, which we can achieve by writing 
\bse 
    \epsilon(x_1) \del\phi^b(x_1) \phi^c(x_2) = \int d^4x \epsilon(x) \del\phi^b(x_1) \phi^c(x_2) \del^{(4)}(x-x_1)
\ese 
and similarly for the $\del\phi^c$ term. We therefore get the result 
\be 
    i\la \phi^b(x_1)\phi^c(x_2) \p_{\mu}j^{\mu}(x) \ra = \la\epsilon(x) \del\phi^b(x_1) \phi^c(x_2)\ra \del^{(4)}(x-x_1) + \la\epsilon(x) \phi^b(x_1) \del\phi^c(x_2)\ra \del^{(4)}(x-x_2).
\ee
where we have suppressed the $\bra{0}...\ket{0}$ to just $\la ... \ra$ for notational reasons. This result is obviously specific to the case of a two point function but it is easy to convince yourself that this result extends to the following 
\mybox{
\be 
\label{eqn:WardIdentity}
    i \la \phi^{a_1}(x_1) ... \phi^{a_n}(x_n) \p_{\mu}j^{\mu}(x) \ra = \sum_{i=1}^{n} \la \phi^{a_1}(x_1) ... \del \phi^{a_i}(x_i) ... \phi^{a_n}(x_n) \ra \del^{(4)}(x-x_i).
\ee     
}
\noindent This result is known as the \textit{Ward Identity} and is a the quantum analogue of Noether's theorem, and it has very powerful uses in QFT. Essentially what it is telling us is that the current is conserved, i.e. $\p_{\mu}j^{\mu}=0$, whenever $x\neq x_i$. However when one of the delta functions on the right-hand side is triggered, we get a non-zero result. Pictorially we can think of this as a plane with the the fields $\phi(x_i)$ inserted and we "move" the $j^{\mu}$ around. When it "hits" one of our $\phi$s we trigger our delta function. Terms of this kind are commonly referred to as \textit{contact terms}.

\subsection{Comparison To Other Literature}

The derivation we have presented above for \Cref{eqn:WardIdentity} is hopefully clearly algebraically sound, however this section is included to help raise a comparison to other literature on the subject. 

As we have done it, everything is done in the \textit{Euclidean} formalism, for this reason there are no time-ordered products or anything appearing in our correlators. This therefore tell us that nothing (apart from $j^{\mu}$ obviously) on the left-hand side of \Cref{eqn:WardIdentity} depends on $x^{\mu}$, and so we are completely free to take the derivative outside the correlator
\bse 
    i \la \phi^{a_1}(x_1) ... \phi^{a_n}(x_n) \p_{\mu}j^{\mu}(x) \ra = i \p_{\mu} \la \phi^{a_1}(x_1) ... \phi^{a_n}(x_n) j^{\mu}(x) \ra.
\ese 

This is all very clear, however subtleties arise when we derive the Ward identity from the canonical picture. We do not present the canonical calculation here, apart from the relevant steps needed for the comparison.\footnote{The interested reader is directed towards Chapter 10.4 of Weinberg.} Using the 2-point function for convenience, and using $y,z$ as the variables for reasons that will be clear in a moment, the Ward identity in this approach is given by 
\be 
\label{eqn:WardIdentityCanonical}
    \p_{\mu}\la \cT[\phi^{a_1}(y)\phi^{a_2}(z) j^{\mu}(x)]\ra = \la \del \phi^{a_1}(y) \phi^{a_2}(z) \ra \del^{(4)}(x-y) + \la \phi^{a_1}(y) \del \phi^{a_2}(z) \ra \del^{(4)}(x-z), 
\ee 
where now we \textit{do} have time orderings. This makes a difference as now when we allow the derivative to act on the full correlator it will act on the Heaviside functions that order our operators and give a temporal delta function. With a bit of thought (or a quick calculation) we see that this gives us 
\bse 
    \begin{split}
        \p_{\mu}\la \cT[\phi^{a_1}(y)\phi^{a_2}(z) j^{\mu}(x)]\ra & = \del(x^0 - y^0) \la [j^0(x),\phi^{a_1}(y)] \phi^{a_2}(z)\ra + \del(x^0 - z^0) \la [j^0(x),\phi^{a_2}(z)] \phi^{a_1}(y)\ra \\
        & \qquad + \la \cT[\phi^{a_1}(y) \phi^{a_2}(z) \p_{\mu}j^{\mu}(x)]\ra 
    \end{split}
\ese 
where the commutators come from the fact that we have to consider both orderings, i.e. 
\bse 
    \cT[A(x)B(y)] = A(x)B(y) \Theta(x^0-y^0) +  B(y)A(x) \Theta(y^0-x^0).
\ese 

Now it turns out to be true that the commutator with the temporal component of $j^{\mu}$ gives the variation times a spatial delta function, 
\bse 
    [j^0(x), \phi^{a_1}(y)] = \del\phi^{a_1}(y) \del^{(3)}(\Vec{x}-\Vec{y}),
\ese 
and similarly for the other field. Putting this all together we have 
\bse 
    \begin{split}
        \p_{\mu}\la \cT[\phi^{a_1}(y)\phi^{a_2}(z) j^{\mu}(x)]\ra & = \la \del\phi^{a_1}(y) \phi^{a_2}(z)\ra \del^{(4)}(x-y) + \la \phi^{a_1}(y) \del\phi^{a_2}(z)\ra \del^{(4)}(x-z) \\
        & \qquad + \la \cT[\phi^{a_1}(y) \phi^{a_2}(z) \p_{\mu}j^{\mu}(x)]\ra. 
    \end{split}
\ese 
However if we compare this to \Cref{eqn:WardIdentityCanonical} we see that we must conclude 
\bse 
    \la \cT[\phi^{a_1}(y) \phi^{a_2}(z) \p_{\mu}j^{\mu}(x)]\ra = 0 \qquad \implies \qquad \p_{\mu}j^{\mu} = 0.
\ese
This last step makes sense as we know that in the canonical theory our current obeys such a relation. This seems harmless until we notice that this suggests that the left-hand side of \Cref{eqn:WardIdentity} vanishes! 

How do we reconcile these two things? The answer is I don't personally know, but Dr. Iqbal suggests that it could be to do with where you choose to analytically continue our Euclidean result to a Minkowski result, where the time orderings become important. Clearly we cannot do it while the derivative is still inside the correlator, as in \Cref{eqn:WardIdentity}, as then we get a vanishing result. It therefore seems that we should first remove the integral and \textit{then} do the analytic continuation on just the correlator part. The idea is that some subtleties could arise in this process that makes the two agree. This section is \textit{not} meant to explain this (again because I don't know how) but merely just to point out that the result is given in different forms depending on how it is derived. 
\chapter{Fermions}

\mybox{
    Just a reminder/disclaimer: This chapter was not part of the taught material, however I have worked from Dr. Iqbal's notes to fill it in here.
}

So far we have only studied scalar fields, however we know that the physics of QFT is much richer than this. In particular we also have Fermions. Now recall from an introduction to QFT course that Fermions are distinctly different to bosons (which the scalar fields are), with the most prominent difference being that underlying fields obey \textit{anti}commutation relations. Let's now look into Fermions and how they appear in the path integral approach. 

\section{Recap Of Canonical Approach To Fermions}

We begin by giving a recap on the canonical approach to Fermions, as this will give us some motivation for what proceeds. Recall that Fermions arise from so-called \textit{Dirac} fields, which come from the action
\bse 
    S[\psi,\overline{\psi}] = \int d^4x \, \overline{\psi}(i\slashed{\p} -m) \psi 
\ese 
where $\slashed{\p} := \g^{\mu}\p_{\mu}$, whose equation of motion is the \textit{Dirac equation}
\bse 
    (i\slashed{\p}-m)\psi = 0.
\ese 

We can show that, upon quantisation, the resulting particles are four component objects, called \textit{spinors}, and, by considering the angular momentum current, we can show they carry spin $1/2$. Equally, as we said above, we see that if you want our energy to be bounded from below\footnote{See, for example, Section 11.2.3 of my IFT notes for more details.} that we require the fields obey (equal time) anticommutation relations:
\bse 
    \{\psi_a(\Vec{x}), \psi_b^{\dagger}(\Vec{y})\} = \del^{(3)}(\vec{x} - \vec{y}) \del_{ab}
\ese 
and all others vanishing. This anticommutation behaviour translates to the creation/annihilation operators anticommuting, which in turn gives us Pauli exclusion --- two Fermions with exactly the same spin and momentum can never occupy the same state. 


\section{Grassman Variables}

The important piece of information we need from this brief recap is the fact that stuff is anticommuting and so "squares to 0". It is not clear at all how we could produce such behaviour from our path integral approach as is: all we have available to us are complex valued functions $\phi(x)$, and these obviously commute. What we want, therefore, is some kind of number that anticommutes so that we can make our fields `some-kind-of-number' valued, thereby giving us anticommuting fields. Obviously no `normal' kind of number does this and so we essentially have to make them up. Let's do that now. 

\subsection{Grassman Numbers}

As uncomfortable as it might seem, we essentially just start by saying our new fancy numbers, known as \textit{Grassman numbers}, anticommute:
\be 
\label{eqn:GrassmanNunbers}
    \theta \eta = - \eta \theta.
\ee 
From here we instantly see that they square to zero:
\be
\label{eqn:GrassmanSquared}
    \theta\theta = -\theta\theta \qquad \iff \qquad \theta^2 = 0.
\ee 
Note it follows immediately from this that \textit{any} function $f(\theta)$ can be expanded as 
\be 
\label{eqn:f(theta)}
    f(\theta) = A + B\theta,
\ee 
where $A$ and $B$ could be (but need not be) Grassman numbers themselves.  One it typically interested in commuting functions like \Cref{eqn:f(theta)}, and so we often take $A$ to be a `normal', commuting number while $B$ is a Grasmann number, this is what we shall do in this course.

\br 
   It is important to note that the above conditions on $A$ and $B$ gives us only a specific class of such functions $f(\theta)$. That is, in general $A$ and $B$ can be either Grassman or `normal'. It follows from this that $f(\theta)$ need not possess any specific commutation/anticommutation behaviour --- e.g. we could take both $A,B\in\C$ and then $f(\theta)$ neither commutes, nor anticommutes, with another general function $g(\eta)$. 
\er 

Ok great, but now what do we do? Well if we are going to use these in our path integrals we need to know how to manipulate them. More specifically we need to know how to differentiate and integrate them. There is essentially an ambiguity about how we go about doing this, and we need to chose a convention. 

\subsubsection{Differentiation}

First let's deal with differentiation. With Grassman numbers the derivative $\p_{\theta}$ is also Grassman with respect to its left/right action. That is 
\bse 
    \overset{\rightarrow}{\p}_{\theta}\theta = - \theta \lar{\p}_{\theta}.
\ese 
We therefore have the ambiguity of where to put this relative minus sign. In this course we will use the (perhaps surprising) convention that the \textit{right} action be positive:
\be
\label{eqn:GrassmanDerivativeConvention}
    \theta \lar{\p}_{\theta} = 1 \qquad \iff \qquad \p_{\theta}\theta = -1,
\ee 
where we have dropped the arrow on the left action (as we are familiar with what it means).

Therefore using our constraints on $A$ and $B$ in \Cref{eqn:f(theta)}, we have 
\bse 
    \p_{\theta}f(\theta) = -B.
\ese

Now what if we take the derivative of a product of Grassman numbers? The idea is that we don't simply use the product rule, but instead commute the numbers so that the derivative and w.r.t. number meet and then just take the derivative there. For clarity we have 
\bse 
    \p_{\theta} (\eta \theta) = -\p_{\theta}(\theta\eta) = -\eta,
\ese 
and similarly for higher products:
\bse 
    \p_{\theta}(\eta_1 \eta_2 ... \eta_n \theta \eta_{n+1} ... ) = (-1)^n \eta_1\eta_2 ... \eta_n\eta_{n+1} ...
\ese
We can therefore think of the derivative itself as a Grassman number and ask "how many other Grassman numbers do I have to anticommute it past to get to its w.r.t. number?" Obviously any `normal' numbers can simply be passed through with no problems. 

\br 
    Note the above procedure will always work and we never have to think about anything like the product rule. This is simply because we will only ever have 1 w.r.t. Grassman number in each expression. This follows simply because 
    \bse 
        \eta_1 ... \theta \eta_i ... \eta_{i+n} \theta ... \eta_m = (-1)^n \eta_1 ... \theta^2 \eta_i ... \eta_{i+n} ... \eta_m = 0,
    \ese 
    where we have used \Cref{eqn:GrassmanSquared}.
\er 

\subsubsection{Integration}

What about integrating Grassman variables, how does that work? Well the idea is to define out integral such that is obeys the kinds of properties we used for the scalar field above. The two most heavily used properties were linearity in the integrand and invariance under a shift of variable. That is we want our Gaussian integral to obey 
\bse 
    \int d\theta \big(g(\theta) + f(\theta)\big) = \int d\theta g(\theta) + \int d\theta f(\theta)
\ese 
and 
\bse 
    \int d\theta f(\theta) = \int d\theta f(\theta+\eta),
\ese 
where the second line follows from $\theta\to \theta+\eta$. 

We can put these two results together in the following manipulation
\bse 
    \begin{split}
        \int d\theta \theta & = \int d\theta (\theta + \eta) \\
        & = \int d\theta \theta + \int d\theta \eta \\
        & = \int d\theta \theta \pm \eta\int d\theta 1 
    \end{split}
\ese
where the $\pm$ is included on the last line as we don't yet know what happens when we move a Grassman $\eta$ through a Grassman $d\theta$. From which we can conclude 
\mybox{
    \be 
    \label{eqn:GrassmanIntegral1}
        \int d\theta \, 1 = 0.
    \ee
}
This seems to be a problem. Why? Well recalling that a 'normal' integral "adds variables", e.g.
\bse 
    \int dx x = x^2,
\ese 
it appears that we get 
\bse 
    \int d\theta \theta = \theta^2 = 0,
\ese 
and so the Grassman integral appears to always give zero? This is obviously wrong, but what do we do? 

Well \Cref{eqn:GrassmanIntegral1} appears to suggest the opposite to the `normal' integral, i.e. instead of "adding variables" we should "take them away". This would give us a finite result for 
\bse 
    \int d\theta \theta = C
\ese
for some constant $C$. Now we note that $C$ must be a `normal', commuting, number as the `product'\footnote{Inverted commas as this is not a real product, but we get the idea.} $d\theta \theta$ will commute (its two Grassman variables). We therefore just adopt the definition $C=1$, so that we have 
\mybox{
    \be 
    \label{eqn:GrassmanIntegralTheta}
        \int d\theta \theta = 1.
    \ee 
}

Finally we need to return to the $\pm$ in the manipulation leading up to \Cref{eqn:GrassmanIntegral1}. Well we just said that we want to treat the `product' $d\theta \theta$ as a commuting number, so it follows that if we want to move the $\eta$ past $d\theta$, we need to pick up a minus sign. We can summarise this in the following convention for doing double Grassman integrals:
\mybox{
    \be 
    \label{eqn:DoubleGrassmanIntegral}
        \int d\theta \int d\eta (\eta \theta) = + 1,
    \ee
}
\noindent which says move the inner most integral's variable to the left and do that integral first (similarly to the derivative case above). 

\br 
    Note it follows from the previous results that the integral over several Grassman numbers is only non vanishing if \textit{every} integrated number appears in the integrand. That is 
    \bse 
        \int d\theta_1  ... d\theta_n f(\{\theta_i\}) \neq 0
    \ese 
    only when (up to a prefactor)
    \bse 
        \prod_{i=1}^n\theta_i \in f(\{\theta_i\}).
    \ese
    We can write this result, along with \Cref{eqn:DoubleGrassmanIntegral}, as the following expression
    \be 
    \label{eqn:nGrassmanIntegrals}
        \int d^n\theta (\theta_{a_1}...\theta_{a_n}) := \int d\theta_{a_1} ... d\theta_{a_n} \big(\theta_{a_1} ... \theta_{a_n}\big) = \epsilon_{a_1...a_n},
    \ee 
    where we have defined what we mean by $d^n\theta$, and where $\epsilon_{a_1...a_n}$ is the $n$-dimensional Levi-Civitia symbol.
\er 

\subsubsection{Complex Conjugation}

Now recall that we have used a lot of complex conjugations in the treatment of the complex scalar field, and recall from your canonical field theory course that we treat the fields $\phi$ and $\phi^*$ as independent degrees of freedom in QFT. The question we now want to ask is "does this work for our Grassman numbers?" The answer is yes and a proof of which is the content of the next exercise. 

\bbox 
    Defining 
    \bse 
        \theta = \frac{1}{\sqrt{2}}\big(\theta_1 + i\theta_2) \qand \theta^* = \frac{1}{\sqrt{2}}\big(\theta_1 - i\theta_2),
    \ese 
    show that 
    \bse 
        \int d\theta d\theta^* (\theta^*\theta) = 1.
    \ese
    \textit{Hint: Use the linearity of $d$ to write}
    \bse 
        d\theta = \frac{1}{\sqrt{2}}\big(d\theta_1 + id\theta_2),
    \ese 
    \textit{and then use \Cref{eqn:GrassmanSquared}.}
\ebox 

It turns out to also be convenient to define the complex conjugation of the product of Grassman numbers to be like the Hermitian conjugate, namely they swap places:
\be 
\label{eqn:ComplexConjTwoGrassman}
    (\theta\eta)^* = \eta^* \theta^* = -\theta^* \eta^*.
\ee 

\subsection{Gaussian Type Integrals}

If we are going to use the methods employed on the scalar field above, we are obviously going to need to know how to do Gaussian type integrals over Grassman numbers. The following exercise will be a good hands on check that the above information has been understood.

\bbox 
    Show that 
    \be
    \label{eqn:GrassmanIntegralExp}
        \int d\theta^* d\theta e^{-\theta^* b \theta} = b,
    \ee
    where $b$ is a `normal' number. 
    
    \textit{Hint: Expand the exponential and recall  \Cref{eqn:f(theta)}.}
\ebox  

We can compare \Cref{eqn:GrassmanIntegralExp} to the `normal' Gaussian result 
\bse 
    \int dx^*dx e^{-x^* b x} = \frac{2\pi}{b}.
\ese
The factor of $2\pi$ is not important,\footnote{If it really bothered us we could redefine \Cref{eqn:GrassmanIntegralTheta} to include it.} but what we see is that the difference between the Grassman integral and `normal' one is that in the former $b$ appears in the numerator, while in the latter it appears in the denominator. We can actually extend this result by introducing the following definition. 

\bd[Pfaffian]
    The \textit{Pfaffian} of a $2n\times 2n$, antisymmetric matrix, $A$, is defined to be 
    \be 
        \text{Pfaff}(A) := \frac{1}{2^n(n!)} A^{a_1a_2} A^{a_3a_4} ... A^{a_{2n-1}a_{2n}} \epsilon_{a_1...a_{2n}}.
    \ee 
\ed 

\bcl 
    The following identity holds 
    \be 
        \big[\text{Pfaff}(A)\big]^2 = \det(A).
    \ee 
\ecl

\bq 
    The proof is provided here, but details can be found via \href{http://scipp.ucsc.edu/~haber/webpage/pfaffian2.pdf}{this link}.
\eq 

Why are we introducing this Pfaffian thing? Well now consider the $2n$-dimensional Grassman integral over the exponential of a matrix equation (which has obvious use for us given the scalar field calculations):\footnote{Note there is no complex conjugation in here, we will return to that in a moment. I just thought it's also worth seeing this result.}
\bse 
    \begin{split}
        \int d^{2N}\theta e^{-\frac{1}{2}\theta_i A^{ij} \theta_j} &= \int  d^{2N}\theta \sum_{n=0}^{N} \frac{(-1)^n}{2^n(n!)} \big(\theta_i A^{ij} \theta_j\big)^n \\
        & = \frac{(-1)^N}{2^N(N!)} A^{a_1a_2} ... A^{a_{2N-1}a_{2N}} \int d^{2N}\theta (\theta_{a_1} ... \theta_{a_{2N}}) \\
        & = (-1)^N \frac{1}{2^N(N!)} A^{a_1a_2} ... A^{a_{2N-1}a_{2N}}\epsilon_{a_1...a_{2N}} \\
        & = (-1)^N \text{Pfaff}(A) \\
        & = (-1)^N \sqrt{\det(A)},
    \end{split}
\ese 
where we have used \Cref{eqn:nGrassmanIntegrals}. 

\br 
    Note that we stopped the sum at $N$ (rather than going to infinity) as $i,j \{1,...,2N\}$ and so if we went beyond $N$ we would end up having a $\theta_i^2$ term, which would kill everything.
\er 

\br 
    Again the above result is in contrast to the `normal' case which we previous showed obeys 
    \bse 
        \int d^{2N}\phi e^{-\frac{1}{2}\phi_a A^{ab} \phi_b} = \sqrt{\frac{(2\pi)^N}{\det(A)}}.
    \ese 
\er 

Ok what about if we do the complex integrals? Well the result of \Cref{eqn:GrassmanIntegralExp} seems to suggest if we remove the $1/2$ from the exponential we will get the the same thing without the square root. Removing the half makes sense for us as it's included as a symmetry factor but when we have a complex field we drop it.\footnote{For example compare the real scalar Klein-Gordan Lagrangian to the complex case. The former has a half, while the second doesn't.} To see that we remove the square root, let's just so the manipulation again using the above as a guide:\footnote{Be careful here, we get one $\epsilon$ symbol for each kind of integration. That is one comes from all the $d\theta$ integrals and another comes from the $d\theta^*$ integrals.}
\bse 
    \begin{split}
        \int d^N \theta^* d^N\theta e^{-\theta^*_i A^{ij} \theta_j} & = \frac{(-1)^N}{N!} A^{a_1a_2} ... A^{a_{2N-1}a_{2N}} \int d^N \theta^* d^N\theta \big(\theta^*_{a_1}\theta_{a_2} ... \theta^*_{a_{2N-1}}\theta_{a_{2N}}\big) \\
        & = (-1)^N \frac{1}{N!} A^{a_1a_2} ... A^{a_{2N-1}a_{2N}} (-1)^N \epsilon_{a_1a_3...a_{2N-3}a_{2N-1}} \epsilon_{a_2a_4...a_{2N-2}a_{2N}} \\
        & = \det(A),
    \end{split} 
\ese 
where we have used the notational shorthand 
\be 
\label{eqn:ComplexGrassmanIntegralMany}
    \int d^N\theta^* d^N\theta := \int d \theta^*_{a_1} d\theta_{a_2} ... d\theta^*_{2N-1} d\theta_{2N}
\ee 
and the determinant definition 
\bse 
    \det(A) = A^{a_1a_2}... A^{a_{2N-1}a_{2N}} \epsilon_{a_1a_3...a_{2N-3}a_{2N-1}} \epsilon_{a_2a_4...a_{2N-2}a_{2N}}.
\ese 

\br 
    Note that the ordering in \Cref{eqn:ComplexGrassmanIntegralMany} matters, i.e. we can't just move the $d\theta^*$s and $d\theta$s through each other. Also note that the way we have done it means that in order to satisfy our convention, \Cref{eqn:DoubleGrassmanIntegral}, we need to anticommute $N$ Grassman numbers, hence the additional $(-1)^N$ on the penultimate line above.
\er 

Ok finally let's see one last integral result. Consider the integral 
\bse 
    I[\eta,\eta^{\dagger}] = \int d^n\theta^* d^n\theta \exp\big( -\theta^{\dagger} A  \theta + \eta^{\dagger}\theta + \theta^{\dagger} \eta\big),
\ese 
where we have adopted a more "matrix like" notation. Let's further work with the special case when $A$ is Hermitian. Now we can use the invariance of the measure under shifts (a property we insisted our integrals obeyed) to take the transformation 
\bse 
    \theta \to \theta + A^{-1}\eta \qand \theta^{\dagger} \to  \theta^{\dagger} + \eta^{\dagger}A^{-1},
\ese
where the Hermitian nature of $A$ has been used. This gives us 
\bse 
    I[\eta,\eta^{\dagger}] = \int d^n\theta^* d^n\theta \exp\big( -\theta^{\dagger}A\theta + \eta^{\dagger}A^{-1}\eta\big) 
\ese 
which we can then split and use the above result to give 
\mybox{
    \be 
    \label{eqn:Ietaeta}
        I[\eta,\eta^{\dagger}] = \det (A) \exp\big(\eta^{\dagger}A^{-1}\eta\big).
    \ee 
}

\section{Fermions From The Path Integral}

We are now armed with the structures we need to obtain Fermions from the path integral. Luckily we have essentially done all the work in the development of the scalar field, we just need to make a couple changes here and there. 

We start by recalling the Dirac action
\be 
\label{eqn:DiracAction}
    S[\psi,\overline{\psi}] = \int d^4 x\, \overline{\psi}(i\slashed{\p} -m)\psi. 
\ee
We want our fields $\psi/\overline{\psi}$ to be Grassman valued functions (so that we get our anticommutation relations). The obvious question to ask is "how do we turn $\psi(x)$ into a anticommuting thing?" Well we simply decompose it as 
\bse 
    \psi(x) = \sum_n \psi_n f_n(x)
\ese
where $f_n(x)$ are `normal' well behaved functions (e.g. plane wave solutions) and $\psi_n$ are Grassman numbers. 

Now recall that in the scalar field case we expressed the partition function $Z$ in terms of a source field $J$ via 
\bse
    Z[J] = \int [\pD\phi] \exp\bigg(iS[\phi] + i\int d^4x \, J(x) \phi(x)\bigg).
\ese 
We now want to do a similar thing for the Fermions. As we have seen we treat the complex conjugate part of a complex Grassman number independently, and so we will have two sources here: one for $\psi$ and the other for $\overline{\psi}$. We have to be a bit careful, though; Fermions are $4$ component spinors, and so we represent them as matrices. The action is meant to consist of just be a collection of numbers (i.e. no matrices) and so we need these source terms to contract the matrix-ness away. With a bit of thought it should be easy to see that what we require is 
\mybox{
    \be 
    \label{eqn:ZFermion}
        Z[\eta, \overline{\eta}] = \int [\pD\overline{\psi} \pD \psi] \exp\bigg( S[\psi,\overline{\psi}] + i\int d^4x \, \big[ \overline{\eta}(x)\psi(x) + \overline{\psi}(x)\eta(x) \big] \bigg)
    \ee 
}
\noindent so that $\eta$ is the source for $\overline{\psi}$ and $\overline{\eta}$ the source for $\psi$. As the notation suggests, $\eta/\overline{\eta}$ are anticommuting fields. 

Next recall that we got the $n$-point functions by considering functional derivatives of the of the free partition function w.r.t. the sources. We now do the same thing for the sources $\eta/\overline{\eta}$, but we have to take into account potential minus signs picked up from anticommutions. That is we have 
\bse 
    \frac{\del}{\del \eta(x)} \big(\overline{\psi}(y) \eta(y)\big) = - \overline{\psi} \frac{\del}{\del \eta(x)}\eta(y),
\ese 
and similarly for $\frac{\del}{\del\overline{\eta}}$. We then extend our derivative convention, \Cref{eqn:GrassmanDerivativeConvention}, to functional derivatives so that the left action gives a minus sign so that 
\bse 
    \begin{split}
        \frac{\del}{\del \eta(x)} \int d^4y\, \big( \overline{\eta}(y) \psi(y) + \overline{\psi}(y) \eta(y)\big) & = +\overline{\psi}(x), \\
        \frac{\del}{\del \overline{\eta}(x)} \int d^4y\, \big( \overline{\eta}(y) \psi(y) + \overline{\psi}(y) \eta(y)\big) & = -\psi(x).
    \end{split}
\ese 

In order to account for this additional minus sign we have to make a small alteration to how we obtain $n$-point Green's functions: basically put an additional minus sign in with the $\overline{\eta}$ derivatives. For example, the $2$-point function is 
\be 
\label{eqn:2PointFunctionFermion}
    \bra{0}\cT[\psi(x_1)\overline{\psi}(x_2)]\ket{0} = \frac{1}{Z_0[0]} \bigg(+i\frac{\del}{\del\overline{\eta}(x_1)}\bigg)\bigg(-i\frac{\del}{\del\eta(x_2)}\bigg) Z_0[\eta,\overline{\eta}]\Big|_{\eta,\overline{\eta}=0}.
\ee 

Recall that the two point function is related to the propagator. We can get this by manipulating \Cref{eqn:ZFermion} with the free action, \Cref{eqn:DiracAction}. It follows from \Cref{eqn:Ietaeta} that the result is 
\mybox{
    \be 
    \label{eqn:Z0Fermion}
        Z_0[\eta,\overline{\eta}] = \det\big(i\slashed{\p}-m\big) \exp\bigg(-\int d^4x d^4y \, \overline{\eta}(x) S_F(x,y) \eta(y)\bigg)
    \ee 
}
\noindent where the propagator is given by
\bse 
    (i\slashed{\p}-m) S_F(x,y) = i\del^{(4)}(x-y) \b1_{4\times 4},
\ese 
where we note that the $4\times 4$ identity matrix is included as $\slashed{\p} = \g^{\mu}\p_{\mu}$ and $\g^{\mu}$ is a $4\times 4$ matrix. As before, we find a nicer expression for the propagator by going to momentum space by replacing $\slashed{\p} \to -i\slashed{p}$ and dropping the $\del^{(4)}(x-y)$:
\mybox{
    \be 
    \label{eqn:FermionPropagator}
        (\slashed{p}-m)S_F(p) = i \qquad \iff \qquad S_F(p) = \frac{i}{\slashed{p} -m}.
    \ee 
}

\br 
    The expression above for $S_F(p)$ is a formal expression, by which we mean it doesn't really make sense to divide by $\slashed{p}$ as it is a matrix. What we mean is to take the matrix inverse. We can write the result in a more "correct" way as
    \bse 
        S_F(p) = \frac{i(\slashed{p}+m)}{p^2-m^2}.
    \ese 
    We can go between these two expressions via 
    \bse 
        \frac{(\slashed{p}-m)}{(\slashed{p}-m)}\frac{i(\slashed{p}+m)}{p^2-m^2} = \frac{i}{\slashed{p}-m}
    \ese 
    where we have used $\slashed{p}^2=p^2$.\footnote{This is easily verified using $\{\g^{\mu},\g^{\nu}\} = \eta^{\mu\nu}$. Check it if you're not familiar.} However \Cref{eqn:FermionPropagator} is less writing and is also standard notation, so we shall use that. 
\er 

\bbox 
    Using \Cref{eqn:Z0Fermion} and \Cref{eqn:2PointFunctionFermion} verify that 
    \bse 
        \bra{0}\cT[\psi(x_1)\overline{\psi}(x_2)]\ket{0} = S_F(x_1,x_2) = \int \frac{d^4p}{(2\pi)^4} \frac{i e^{ip\cdot(x_1-x_2)}}{\slashed{p}-m},
    \ese 
    where the second line just follows from the Fourier tranfsorm of \Cref{eqn:FermionPropagator}.
    
    \textit{Hint: Be careful about signs.}
\ebox 

\br 
    \textcolor{red}{There's a comment in Nabil's notes here about the fact that the determinant appears in the numerator for Grassman variables being related to the fact that the zero-point energy for Fermions being negative. He gives a reference to Chapter II.5 of Zee. Look this up later and write something here.}
\er 

\section{Feynman Rules For Fermions}

The last thing we need to do is obtain the Feynamn rules for our Fermion theory. We also want to take into consideration interactions and so we work with the common Yakawa interaction theory with action 
\bse 
    S[\psi,\overline{\psi},\phi] = \int d^4x \, \bigg[ \overline{\psi}(i\slashed{\p}-m)\psi + \frac{1}{2}(\p\phi)^2 - \frac{M^2}{2} \phi^2 + g\overline{\psi}\psi\phi\bigg].
\ese 
With a bit of thought we see that this corresponds to a theory containing Fermions of mass $m$, a real scalar field $\phi$ of mass $M$ and a $3$-point interaction. For this reason it is clear that we will get the relevant propagators, which we have already derived. We therefore just need to look at the interaction term. 

We follow everything through as above, and we expand the interaction theory around the free theory in powers of $g$:
\bse 
    Z[\eta,\overline{\eta},J] = Z_0[0] \Bigg[ ig \int d^4x \, \bigg(-i\frac{\del}{\del J}\bigg) \bigg(-i\frac{\del}{\del \eta_a}\bigg) \bigg(+i\frac{\del}{\del \overline{\eta}_a}\bigg) \Bigg] Z_0[\eta,\overline{\eta},J],
\ese
where a sum over $a$ is assumed,\footnote{This just tells us that the theory doesn't allow for $\phi$ to change the type of Fermion we are considering. This is the case, for example, in QED where the photon cannot cause intergeneration changes.} and where $\eta/\overline{\eta}/J$ are the relative sources. Our free partition function is now given by 
\bse 
    Z_0[\eta,\overline{\eta},J] \sim \exp\bigg(-\int d^4xd^4y \, \overline{\eta}(x) S_F(x,y)\eta(y)\bigg) \exp\bigg(-\frac{1}{2}\int d^4x d^4y \, J(x) D_F(x,y) J(y)\bigg).
\ese 

\bbox 
    Use the above to show that the interaction vertex term is given by 
    \bse 
        D_1 = ig \int d^4xd^4yd^4zd^4w \, \big[ \overline{\eta}(x) S_F(x,y)S_F(y,z)\eta(z)\big] D_F(y,w)J(w).
    \ese 
    This corresponds to the diagram 
    \begin{center}
        \btik 
            \midarrow[rotate around ={45:(0,0)}] (0,0) -- (-1,0);
            \node at (-0.85,-0.85) {$\overline{\eta}$};
            \midarrow[rotate around ={-45:(0,0)}]  (-1,0) -- (0,0);
            \node at (-0.85,0.85) {$\eta$};
            \draw[thick, dashed] (0,0) -- (1,0) node [right] {$J$};
            \draw[fill=black] (0,0) circle [radius=0.07cm];
        \etik 
    \end{center}
\ebox 

The result of the previous exercise basically tells us that the vertex comes with a factor of $ig$. The "trained" field theorist would "see" this by staring at the Lagrangian (i.e. basically just read off the prefactor and do the symmetry factors in your head).

\br 
    In the diagram in the above exercise we have adopted the convention that the arrows point towards the sources $\overline{\eta}$. As we are familiar with from our canonical approach, these arrows are needed whenever we have a complex field as they allow us to distinguish particles from antiparticles. 
\er 

\subsection{The Infamous "$-1$ For Closed Fermion Loops"}

We have seen all the Feynman rules for our Dirac-Yakawa theory, expect for one: include a factor of $-1$ for every closed Fermion loop. This is one that is often injected into the canonical approach\footnote{For example, see Feynman rule (v)(a) in my QED notes.} with no explanation apart from "it's true". Here we present where it comes from. 

First let's draw the closed Fermion loop diagram. With a bit of thought, we see that it is 
\begin{center}
    \btik 
        \draw[thick, dashed] (-1.5,0) -- (-0.5,0);
        \draw[thick, dashed] (0.5,0) -- (1.5,0);
        \draw[thick] (0,0) circle [radius=0.5cm];
        \draw[fill=black] (-0.5,0) circle [radius=0.07cm];
        \draw[fill=black] (0.5,0) circle [radius=0.07cm];
    \etik 
\end{center}
If we label the left-hand side vertex point with integration variable $x$ and the right-hand side one with $y$ then it's clear that this diagrams contains the following term
\be 
\label{eqn:SFLoop}
    S_F(x,y) S_F(y,x) = \del_x\overline{\del}_y \big(\overline{\eta}\cdot S_F\cdot \eta) \,  \del_y\overline{\del}_x \big(\overline{\eta}\cdot S_F\cdot \eta)
\ee
where we have used the shorthand notation 
\bse 
    \del_x := \frac{\del}{\del \eta(x)} \qand \overline{\del}_x := \frac{\del}{\del \overline{\eta}(x)}
\ese
and the same $\overline{\eta}\cdot S_F \cdot \eta$ notation that we used for the scalar field (i.e. for $J\cdot D_F \cdot J$).

Now a term of this form will come from our path integral in the form
\bse 
    \int d^4x d^4y, \del_x \overline{\del}_x \del_y \overline{\del}_y \big(\overline{\eta} \cdot S_F \cdot \eta\big)\big(\overline{\eta} \cdot S_F \cdot \eta\big),
\ese 
and so in order to get \Cref{eqn:SFLoop} we need to anticommute an odd number of functional derivatives. We therefore see that a closed Fermion loop comes with a $(-1)$ factor. 

\subsection{A Nod To SUSY}

As we have seen during this course, the difficulties in QFT come from loop diagrams. However we have also just seem that Fermion loops contribute a relative minus sign to Boson loops, so we could ask the question "is it possible these two things cancel?" This would clearly make our QFT lives a lot easier. Obviously there is absolutely no reason why this would happen a priori, and so we have to ask questions like "what constraints can we put on the system so that this happens?" One possibility would be to \textit{enforce} this result via some symmetry. This is one of the big motivations behind supersymmetry (SUSY).\footnote{Notes on SUSY to come on my website.} Unfortunately it seems as though nature doesn't want to play ball with SUSY, and so at least currently we need to continue worrying about loops. 
\chapter{Abelian Gauge Theory}

We previously worked through global symmetries in the path integral approach and derived the Ward identities. We are yet to talk about local (gauge) theories, though. The rest of this course is dedicated to exactly this discussion. We, of course, deal with the easier abelian gauge theories first and then move on to the full beast of non-abelian gauge theories. 

\section{Gauge Invariance}

In order to discuss gauge symmetry, it is often useful to first pick a specific global symmetry and then "make it local" to see how things change. Here we do this by considering a theory with a global $U(1)$ theory. We will use the Dirac action as a proxy. 

Recall that the Dirac action is given by
\bse 
    S[\psi,\overline{\psi}] = \int d^4x \, \overline{\psi} \big( i\slashed{\p} - m\big) \psi
\ese 
This admits a global $U(1)$ symmetry which acts as
\be
\label{eqn:GlobalU(1)}
    \psi(x) \to \psi'(x) = e^{i\Lambda} \psi(x), \qand \overline{\psi}(x) \to \psi'(x) = \overline{\psi}'(x) e^{i\Lambda}
\ee 
where $\Lambda$ is \textit{constant} in spacetime. Any two configurations that are related by the symmetry are both physical. That is if $\psi$ is a field we integrate over in the path integral, then we must also integrate over $\psi'$. This is an important point and something that we will return to when talking about gauge symmetries. 

Let's now "make this into a gauge symmetry", i.e. let's \textit{demand} invariance under \Cref{eqn:GlobalU(1)} with $\Lambda(x)$ being spacetime dependent. As we will see, the truth this that this is not a symmetry in the same way that the global one was. Really we should refer to it as a \textit{gauge redundancy}, however we will not use this terminology in this course much. Things that are invariant under our gauge symmetries are called \textit{gauge invariant}. 

Let's now see the repercussions of this procedure, i.e. let's ask the question "What is, and more importantly what isn't, gauge invariant in our current Lagrangian?" Obviously the mass term is fine:
\bse 
    m\overline{\psi} (x)\psi(x) \to m \overline{\psi} e^{-i\Lambda(x)} e^{i\Lambda(x)}\psi(x) = m \overline{\psi}(x)\psi(x).
\ese 
Now note that this worked out because we had the fields evaluated at the \textit{the same spacetime point}. That is if we had 
\bse 
    m\overline{\psi}(x)\psi(y) \to m\overline{\psi} e^{i(\Lambda(x)-\Lambda(y))} \psi(y) \neq m\overline{\psi}(x)\psi(y)
\ese 
when $x\neq y$. It follows from this that we expect the kinetic term with derivatives to \textit{not} be invariant. This is simply because, roughly speaking, the derivative compares neighbouring points. Indeed this turns out to be the case, and a quick calculation reveals
\bse 
    \p\psi' \to e^{i\Lambda(x)} \big(i\p_{\mu}\Lambda + \p_{\mu}\big)\psi(x). 
\ese
This tells us that our Dirac action is \textit{not} gauge invariant and so we don't have a gauge symmetry. This is not good, and we need to do something to fix it. 

This should prick up the ears of people familiar with GR: we have some derivative behaviour that doesn't quite add up. We had a similar problem in GR when taking derivatives of tensor fields. In this case we replaced the partial derivative by the covariant derivative, which was just given by the original partial derivative plus something (the connection coefficients). We draw from this previous experience and we seek to replace our partial derivative here with a so-called \textit{gauge invariant derivative}, which we define via\footnote{The $e$ that appears in here is the \textit{gauge field coupling}. If this doesn't make any sense to you, see, for example, my QED notes.}
\mybox{
    \be 
    \label{eqn:GaugeInvariantDerivative}
        D_{\mu}\psi := \big(\p_{\mu} + ieA_{\mu}\big) \psi,
    \ee 
}
\noindent where we have introduced a new object $A_{\mu}$ called the \textit{gauge field} or \textit{gauge connection}. For a abelian theory we require 
\be 
\label{eqn:Abelian}
    [A_{\mu},A_{\nu}] = 0 \qquad \forall \mu,\nu \in \{1,...,d\}. 
\ee 

We now use our gauge invariant derivative to make our action gauge invariant. How do we do this? Well we \textit{demand} that the gauge covariant derivative of $\psi$ has a \textit{nice} transformation property. What do we mean by "nice"? Well looking at the mass term calculation, we see that if we have 
\bse 
    D_{\mu}\psi(x) \to D_{\mu}' \psi'(x) \overset{!}{=} e^{-\Lambda(x)} D_{\mu}\psi(x). 
\ese 
then we expect the action to be gauge invariant. 

Ok great, but how to we make such a demand? The answer is that we allow the gauge field $A_{\mu}$ to transform under our gauge transformation, and define its trasnformation property such that our demand holds. Explicitly we have 
\bse 
    \begin{split}
        D_{\mu}' \psi'(x) & = \big(\p_{\mu} + ieA_{\mu}'\big) e^{i\Lambda(x)}\psi(x) \\
        & = e^{i\Lambda(x)}\big(i\p_{\mu} \Lambda + ieA_{\mu}' + \p_{\mu}\big) \psi(x) \\
        & \overset{!}{=} e^{i\Lambda(x)}\big(\p_{\mu} + ieA_{\mu}\big) \psi(x)
    \end{split}
\ese 
from which it follows that we require 
\mybox{
    \be
    \label{eqn:GaugeTransformationOnFields}
        A_{\mu}' = A_{\mu} - \frac{1}{e} \p_{\mu}\Lambda, \qand \psi'(x) = e^{i\Lambda(x)}\psi. 
    \ee
}

\bbox 
    Confirm that the above transformation property does indeed give us a gauge invariant action. That is show that 
    \bse 
        \overline{\psi}'(x)(i\slashed{D}'-m)\psi'(x) = \overline{\psi}(x)(i\slashed{D}-m)\psi(x).
    \ese 
\ebox 

So we have a  fully gauge-invariant Dirac action 
\bse 
    S[\psi,\overline{\psi}] = \int d^4x \, \overline{\psi}(x) \big[ i\slashed{D} - m\big] \psi(x),
\ese 
but what is $A$ good for?\footnote{Besides taking derivatives.} Well, firstly note that because $D_{\mu}\psi$ transforms nicely, therefore so does $D_{\mu}D_{\nu}\psi$.\footnote{Hopefully this is clear. Basically the idea is we defined $D_{\mu}$ to transform such that it took something that transformed as $A \to e^{i\Lambda(x)}A$ so that $D_{\mu}A \to e^{i\Lambda(x)}D_{\mu}A$. Well we have just shown that $D_{\mu}\psi$ transforms in exactly this way and so it follows immediately that $D_{\nu}D_{\mu}\psi \to e^{i\Lambda(x)}D_{\nu}D_{\mu}\psi$.} We may, then, consider the commutator of two $D$s. 
\bse 
    \begin{split}
        [D_{\mu},D_{\nu}]\psi & = [\p_{\mu},\p_{\nu}] \psi + ie \big([\p_{\mu},A_{\nu}] + [A_{\mu},\p_{\nu}]\big) \psi \\
        & = ie \big(\p_{\mu}A_{\nu} - \p_{\nu}A_{\mu}\big)\psi.
    \end{split}
\ese
where we have made use of \Cref{eqn:Abelian}. We can then define 

\mybox{
    \be     
    \label{eqn:Fmunu}
        ieF_{\mu\nu} := [D_{\mu},D_{\nu}],
    \ee 
}
known as the \textit{field strength}. This definition holds for both abelian and non-abelian theories. For the abelian case we have 
\be 
\label{eqn:FmunuAbelian}
    F_{\mu\nu}^{\text{Abelian}} = \p_{\mu}A_{\nu} - \p_{\nu}A_{\mu}
\ee 

\bnn 
    From now on we will drop the superscript "Abelian" and unless needed in order to avoid confusion. 
\enn 

\bbox 
    Show that \Cref{eqn:FmunuAbelian} is itself gauge invariant. This can be shown explicitly using \Cref{eqn:GaugeTransformationOnFields} or it can be shown more generally using
    \bse 
        [D_{\mu},D_{\nu}]\psi \to e^{i\Lambda(x)}[D_{\mu},D_{\nu}]\psi.
    \ese 
    Try using the second method. \textit{Hint: Compare this to the transformation of $\psi$ itself.}
\ebox 

Why are we bothering to define the field strength? Well because we now note that it can be included into our action in a gauge invariant and renormalisable way. That is we can extend our Dirac action to be 
\bse 
    S[\psi,\overline{\psi},A] = \int d^4 x \bigg[ -\frac{1}{4} F_{\mu\nu}F^{\mu\nu} + \overline{\psi}(x)\big(i\slashed{D} - m\big)\psi(x)\bigg].
\ese
This additional term is important as it contains a kinetic term for the gauge field. That is we get terms of the form 
\bse 
    \p A \p A,
\ese 
where the indices have been dropped. Putting this together with the fact that we do \textit{not} have a term of the form 
\bse 
    \frac{1}{2}m_A^2 A^2,
\ese 
we see that this new term makes our gauge field into a propagating \textit{massless} degree of freedom. In QED $A_{\mu}$ is exactly the photon. 

\bbox 
    Verify that the new $F_{\mu\nu}F^{\mu\nu}$ term is indeed renormalisable. Also show that a mass term is forbidden by gauge invariance. \textit{Hint: Use dimensional arguments}.
\ebox 

\br 
    As the above exercise shows, we exclude the mass term on gauge invariant grounds. It turns out that (for QED, at least) even if it wasn't excluded on these grounds it would instead be excluded because it would lead to the theory being non-renormalisable. This is a non-trivial statement to see, and comes from considering so-called \textit{box} diagrams. For more details see Section 3.3.5 of my QED notes.
\er 

The next question we can ask is "Is that it, or can we add more gauge invariant, renormalisable terms to our action?" The answer is that we can indeed add another term, giving us the most general renormalisable, gauge invariant action involving $A_{\mu}$ and $\psi$:
\bse 
    S[\psi,\overline{\psi}, A] = \int d^4 x \bigg[ -\frac{1}{4} F_{\mu\nu}F^{\mu\nu} + \overline{\psi}(x)\big(i\slashed{D} - m\big)\psi(x) + \Theta \epsilon^{\a\beta\mu\nu} F_{\mu\nu}F_{\a\beta}\bigg],
\ese 
where $\epsilon^{\a\beta\mu\nu}$ is the $4$-dimensional Levi-Civita symbol and $\Theta\in\R$ is just some constant. This last term looks scary and we would perhaps like it to not be there. Well, it turns out that 
\ben[label=(\roman*)]
    \item The $\epsilon$ symbol breaks parity and time reversal, and
    \item It is a total derivative so in the classical theory we can drop it as it doesn't change our equations of motion.
\een  
If we, therefore, just specialise to theories that preserve parity-time (a physically reasonable thing to do) we can drop it (i.e. set $\Theta=0$) and move on. 

\br 
    This last term is actually important in the quantum theory. To see why we need to use some differential geometry:\footnote{If you are not familiar with this stuff, don't worry just move past this remark.} after staring at the expression for a moment we realise that this term if just $\theta F\wedge \star F$,\footnote{$\star$ here is the Hodge operator, it maps $n$-forms to $(d-n)$-forms. So here it maps a $2$-form to a $(4-2=2)$-form.} and recalling that $F$ is a 2-form we see that this term is a top form (we are considering a $4$-dimensional spacetime). The integral over it will, therefore, just be some number, say $C$, times $\theta$. In fact it turns out that the integral over the top form is an integer times $2\pi$, and so from this (after we recall that the action appears as $e^{iS}$) it follows that $\theta \in [0,2\pi)$. This whole term is referred to as a \textit{theta term}. Classically it has no effect on the equations of motion, as per (ii) above, however when we quantise the theory this term ends up weighting different fields by different numbers\footnote{I think called a \textit{winding number}, but this could be wrong. \textcolor{red}{Come back and maybe add a bit clearer/more correct statement once you've had time to read more of Tong's Gauge theory.}} and encodes information about the topological nature of the fields. 
\er 

\section{Some Classical Aspects of Abelian Gauge Theory}

In order to get a motivation for a result we will shortly need, let us just consider the following part of the above action
\bse 
    S[A]_{\text{Maxwell}} = \int d^4 x \bigg[ -\frac{1}{4} F_{\mu\nu}F^{\mu\nu}\bigg].
\ese 
If we vary this w.r.t. $A_{\mu}$ we get
\be 
\label{eqn:Feom}
    \p_{\mu}F^{\mu\nu} = 0.
\ee 
We can use these results to show that the above is really just Maxwell dynamics. We achieve this by picking the specific basis such that  $F^{0i}=E^i$ and $F^{ij} = \epsilon^{ijk} B^k$. Then the above condition (along with some other steps) can be used to derive the Maxwell equations. 

Why are we bringing this up? Well recall Klein-Gordan equations of motion for massless scalar were 
\bse 
    \p^2 \phi = 0.
\ese 
This defines what is known as a \textit{well-posed Cauchy problem}. This basically says that our theory is predictable. Let's now show that this is indeed the case for the above equation of motion: pick a time, say $t=0$ for simplicity, and then specify $\phi(0,x^i)$ and $\p_t \phi(0,x^i)$. We can then solve the equations of motion to find out how $\phi(t,x^i)$ propagates in time. We can do this unambiguously and obtain the actual behaviour of $\phi(t,x^i)$ for all $t$. 

What about \Cref{eqn:Feom}? Well the fundamental degree of freedom is $A_{\mu}(x)$, and so this plays the role of $\phi(x)$ for the Klein-Gordan field. The question we want to ask is "do these pose a well-posed Cauchy problem?" The answer is "no", and you can obtain this by solving the problem. However we can be a bit cleverer and simply recall that $F^{\mu\nu}$ was itself gauge invariant. This tells us that if $A_{\mu}(x)$ is a solution to our problem then so is $A_{\mu}'(x) = A_{\mu}(x) -\frac{1}{e}\p_{\mu}\Lambda(x)$. $\Lambda(x)$ here can be an arbitrary function on the spacetime, and so it follows that we could choose it so that our two solutions differ at large $t$ values. That is, we could pick $\Lambda(t\approx0,x^i)=0$ but $\Lambda(t >>0, x^i) \neq 0$, and so two solutions which agreed on our our initial Cauchy surface ($t=0$) differ at a later Cauchy surface ($t >>0$). This basically tells us that we can not predict what the theory will do a later time, giving some initial data. This is essentially the statement that the theory is non-predictive. 

Hmm, this seems like a bit of a problem. How do we deal with this? Well, we declare a very important result:
\mybox{
    \begin{center}
        Things that are not gauge invariant are \textit{not} physical.
    \end{center}
}

\br 
    Note the above declaration only constrains things that aren't gauge invariant, but says nothing about gauge invariant quantities. It turns our that things that are gauge invariant will always have a well-posed Cauchy problem. 
\er 

How does this declaration help us? Well the idea is we \textit{fix a gauge} and then use that to give us a well-posed Cauchy problem. For example, we can use \textit{Lorenz}\footnote{Note there is no $t$ in his name... unfortunate for him, people don't always notice this. } gauge 
\bse 
    \p_{\mu}A^{\mu} = 0 \qquad \implies \qquad \p^2 A^{\nu} = 0.  
\ese
In Fourier space $e^{i\omega t+ ikz}$ we then get 
\bse 
    (\omega^2 - k^2) A^{\nu} = 0 \qquad \implies \qquad \omega_k=k,
\ese 
so the photon is massless. 

\br 
    This doesn't actually completely fix our gauge. That is there is a remaining redundancy, which can be used to show that there are only two polarisations. We do not discuss this further here.
\er 

\section{Quantise QED}

We want to understand the path integral\footnote{We drop the source terms for ease of notation.}
\bse 
    Z = \int [\pD A] \exp \bigg( -\frac{i}{4} \int d^4 x F_{\mu\nu}F^{\mu\nu}\bigg). 
\ese 
Of course this is only part of the full action, we have dropped the Fermion part. We are not interested in that here, but instead want to look for the propagator for our gauge field. We can get that from here as this part contains the quadratic derivative terms in $A_{\mu}$.

Recall that we get the propagator by finding the inverse of the operator appearing between two $A_{\mu}$s. Explicitly, we have
\bse 
    \begin{split}
        S[A] & = -\frac{1}{4}\int d^4x \big( \p_{\mu}A_{\nu} - \p_{\nu}A_{\mu}\big) \big( \p^{\mu}A^{\nu} - \p^{\nu}A^{\mu}\big) \\
        & = \frac{1}{2} \int d^4x \big( A_{\nu}\p^2 A^{\nu} - A_{\nu} \p^{\nu}\p_{\mu}A^{\mu} \big) \\
        & = \frac{1}{2} \int d^4 x A_{\mu} \big(\p^2 \eta^{\mu\nu} - \p^{\mu}\p^{\nu}\big)A_{\nu}. 
    \end{split}
\ese 
We then find the propagator $D^F_{\mu\nu}$, by insisting it satisfies\footnote{Note that $D_{\mu\nu}^F$ is symmetric in its indices.} 
\bse 
    \big(\p^2 \eta^{\mu\nu} - \p^{\mu}\p^{\nu}\big) D_{\nu\rho}^F(x,y) = \del^{\mu}_{\rho} \del^{(4)}(x-y).
\ese 

The problem is, there is no \textit{unique} such $D^F$. In other words, we can't invert our differential operator. To see this, let's go to Fourier space, where we have 
\bse 
    \big(p^2\eta^{\mu\nu} - p^{\mu}p^{\nu}\big)D^F_{\mu\rho} = -i\del^{\nu\rho}. 
\ese
Now consider the test function $p_{\mu} \a(p)$:
\bse 
    \big(p^2\eta^{\mu\nu} -p^{\mu}p^{\nu}\big) p_{\mu}\a(p) = \big(p^2p^{\nu} - p^2 p^{\nu}\big)\a(x) = 0,
\ese 
where $\a(p)$ is an \textit{arbitrary} function. It follows from the arbitrariness of $\a(p)$ to conclude that we don't have a unique kernel, therefore it is non-invertable.\footnote{More mathematically, the map is non-injective as it doesn't have a unique kernel.}

We can express the result more physically: any `pure gauge'
\bse 
    A_{\mu} = \frac{1}{e}\p_{\mu}\Lambda,
\ese    
has $F=0$ which implies $S[A]=0$. These are unsuppressed in the path integral, and it this gives us bad divergences. The problem is that we are still integrating over these in our path integral and that's not a clever idea as they are unphysical. We fix this by introducing something known as the \textit{Faddeev-Popov} procedure. Let's outline how this works now. 

We start by defining $G(A)$ be a function that we set to $0$ to fix our gauge, e.g. for Lorenz $G(A) := \p_{\mu}A^{\mu}$. Next recall the normal delta functions, $\del : \R \to \R$, which satisfies $\del(x)=0$ for all $x\neq 0$ and it integrates to one. 

As we have done with several things in this course already, we want to extend this to the case of functionals. That is we want to introduce the \textit{functional delta function}: it maps functions to the real numbers and\footnote{The condition $\phi(x) = 0$ here means the zero-function, that is the function $\phi : \R \to 0\in \R$.}
\bse 
    \begin{split}
        \del[\phi] & =0 \qquad \text{if} \qquad  \phi(x) \neq 0 \\
        \int [d\phi] \del[\phi] & = 1.
    \end{split}
\ese 

We now want to somehow stick $\del(G(A))$ into our path integrals. In order to do this, let's first define the gauge transformed field 
\bse 
    A^{\Lambda}_{\mu}(x) := A_{\mu} - \frac{1}{e}\p_{\mu}\Lambda.
\ese
Then we introduce a fancy way of writing 1:
\be 
\label{eqn:FunctionalDelta1}
    1 = \int [\pD \Lambda] \del\big(G(A^{\Lambda})\big) \det\bigg(\frac{\del G(A^{\Lambda})}{\del \Lambda}\bigg),
\ee
which is just the infinite dimensional version of 
\bse 
    1 = \int dy \del(y) = \int dx \del(y(x)) \frac{dy}{dx}. 
\ese

So, we take this and insert into the path integral to give us
\bse 
    Z = \int [\pD \Lambda][\pD A] \del\big(G(A^{\Lambda})\big) \det\bigg(\frac{\del G(A^{\Lambda})}{\del \Lambda}\bigg) e^{iS[A]}.
\ese 
This doesn't seem to have helped much. The thing that really doesn't look nice here is the determinant. Let us, therefore, pick a convenient $G$
\be 
\label{eqn:Gomega}
    G_{\omega}(A) = \p_{\mu}A^{\mu} - \omega, 
\ee 
so that 
\bse 
    \frac{\del G(A^{\Lambda})}{\del \Lambda} = \frac{\del}{\del\Lambda} \bigg( \p_{\mu}\bigg(A^{\mu} - \frac{1}{e} \p^{\mu}\Lambda\bigg) - \omega\bigg) = - \frac{1}{e}\p^2.
\ese 
This means our determinant can be pulled out the path integral as it doesn't depend on $A$ or $\Lambda$.\footnote{As we will see later, this is only true for the non-Abelian case.} We therefore get 
\be 
\label{eqn:Zomega}
    Z_{\omega} = \det\bigg(\frac{1}{e}\p^2\bigg) \int [\pD \Lambda] [\pD A] \del\big( G_{\omega}(A^{\Lambda})\big) e^{iS[A]}.
\ee
where the subscript is there to remind us that we're using \Cref{eqn:Gomega}. 

Ok this is a bit nicer, but we still have to worry about the path integral over our gauge parameters $[\pD\Lambda]$. Well we can fix this to by using the following claim. 

\bcl 
    The path integral measure $[\pD \Lambda]$ is gauge invariant. More specifically it is left and right action invariant:
    \be 
    \label{eqn:PathIntegralMeasureInvariant}
        [\pD (\Lambda' \Lambda)] = [\pD \Lambda] = \big[\pD (\Lambda \Lambda')\big],
    \ee 
    where $\Lambda'$ is our transformation. 
\ecl 

We then use \Cref{eqn:PathIntegralMeasureInvariant} to take the inverse transformation on the whole of \Cref{eqn:Zomega} without effecting our $[\pD \Lambda]$, i.e. we send $A_{\mu} \to A_{\mu}^{\Lambda^{-1}}$, so that $A_{\mu}^{\Lambda} \to A_{\mu}$, to give us 
\bse 
    Z_{\omega} = \det\bigg(\frac{1}{e}\p^2\bigg) \int [\pD A^{\Lambda^{-1}}] [\pD \Lambda] \del\big(G_{\omega}(A)\big) e^{iS[A^{\Lambda^{-1}}]}.
\ese 
Now we use that the measure $[\pD A]$ is gauge invariant\footnote{This can argued from the fact that a gauge transformation is just a shift in $A_{\mu}$ and so the integral measure just sees it as a change of variables} and the fact that we have a gauge invariant to obtain 
\bse 
    Z_{\omega} = \det\bigg(\frac{1}{e}\p^2\bigg) \int [\pD \Lambda] [\pD A]  \del\big(G_{\omega}(A)\big) e^{iS[A]}.
\ese
Finally we note that nothing in the integrand depends on $\Lambda$ and so we can "factor out" the path integral over $[\pD\Lambda]$, absorbing it into an overall constant
\bse 
    \cN = \det\bigg(\frac{1}{e}\p^2\bigg) \int [\pD \Lambda]. 
\ese
This is where our bad divergence was coming from; as we said before we have an unsuppressed path integral and so we get an infinite contribution. However this is going to appear for all our partition functions and so we just "forget about it" and pretend $\cN$ is finite. 

Great, so we have removed the $\Lambda$ dependence of our path integral to obtain a finite result. The last thing we need to do is evaluate our functional delta function. In order to do this we need to remove the $\omega$ dependence, the question is "how do we do that?" Well we want to average over all $\omega$ values, and so we do this with a Gaussian with weight $\xi$, giving us\footnote{$\cN'$ is clearly related to $\cN$ above, but includes other factors from our averaging.} 
\bse 
    Z = \cN' \int [\pD \omega] [\pD A] \exp\bigg(-i\int d^4x \frac{1}{2\xi} \omega^2\bigg) \del\big(\p_{\mu}A^{\mu}-\omega\big) e^{iS[A]}.
\ese
If we then do the $\omega$ integral we get 
\bse 
    Z = \cN' \int [\pD A] \exp\bigg(iS[A] - i\int d^4x \, \frac{1}{2\xi} (\p_{\mu}A^{\mu})^2\bigg),
\ese
which is the gauge-fixed form of the path integral. 

We therefore have the Faddeev-Popov action 
\mybox{
    \be 
    \label{eqn:FadPopActionQED}
        S_{FP}[A] = \int d^4 x \, \bigg( -\frac{1}{4}F_{\mu\nu}F^{\mu\nu} - \frac{1}{2\xi} (\p_{\mu}A^{\mu})^2\bigg).
    \ee
}
\noindent We can use this to find the propagator for our gauge field (the photon). The first thing is to write this in our usual form by integrating by parts
\bse 
    S_{FP}[A] = = \int d^4 x \, \frac{1}{2}\bigg( A_{\mu} \bigg[ \p^2\eta^{\mu\nu} - \bigg(1-\frac{1}{\xi}\bigg) \p^{\mu}\p^\nu\bigg)A_{\nu}\bigg),
\ese 
and so we see that, in Fourier space, we require 
\be 
\label{eqn:QEDPropagatorEquation}
    \bigg(-p^2\eta_{\mu\nu} + \bigg( 1 - \frac{1}{\xi}\bigg) p_{\mu}p_{\nu}\bigg) D_F^{\nu\rho}(p) = i \del^{\rho}_{\mu}. 
\ee 
We can guess the answer by noting that we need tensor indices and the only things that we have available to us are the metric $\eta^{\mu\nu}$ and momentum $p^{\mu}$. So we guess 
\be
\label{eqn:QEDDFGuess}
    D_F^{\nu\rho}(p) = A(p) \eta^{\nu\rho} + p^{\nu}p^{\rho} \frac{B(p)}{p^2},
\ee 
which is we solve for $A(p)$ and $B(p)$. 

\bbox 
    Plugging \Cref{eqn:QEDDFGuess} into \Cref{eqn:QEDPropagatorEquation} show that our propagator is given by
    \bse 
         D_F^{\nu\rho}(p) = -\frac{i}{p^2}\bigg( \eta^{\mu\nu} - (1-\xi)\frac{p^{\mu}p^{\nu}}{p^2}\bigg).
    \ese 
\ebox 

The result of the above exercise is very important (it's the photon propagator!), so we write it again in one of our nice boxes. 
\mybox{
    \be 
    \label{eqn:PhotonPropagator}
         D_F^{\nu\rho}(p) = -\frac{i}{p^2}\bigg( \eta^{\mu\nu} - (1-\xi)\frac{p^{\mu}p^{\nu}}{p^2}\bigg).
    \ee 
}

\noindent The natural question to ask is "what is $\xi$?" The answer is it is a gauge parameter and therefore, by gauge invariance, nothing can depend on it. We can therefore set it to anything we like! Of course a common choice is $\xi=1$ as this removes the awkward fraction term. This is known as \textit{Feynman gauge}. Other gauges exist, and more details can be found, for example, in my QED notes. 

Wow that was a lot of work, but the good news is we're done! That is we now have all the information we need to derive the full set of Feynman rules for QED and plough on with our calculations. For clarity, QED is given by the following action
\bse 
    S[\psi,\overline{\psi},A] = \int d^4 x\, \bigg[-\frac{1}{4}F^2 + \overline{\psi}(x) \big( i\slashed{D} - m\big) \psi(x)\bigg]
\ese 
with the following Feynman rules:
\ben[label=(\roman*)]
    \item All the usual stuff,
    \item The propagator is given by \Cref{eqn:PhotonPropagator},
    \item The Fermion propagator is given by \Cref{eqn:FermionPropagator}, 
    \item The vertex is simply $-ie\g^{\mu}$,\footnote{Make sure you understand why it's this.}
    \item Don't forget to include a factor of $(-1)$ for every closed Fermion loop. 
\een 
We do not go through any detailed diagram calculations here as this is done to death in the QED course, and any unfamiliar reader is directed there. 

\br 
    It is important to note that our vertex term came from the $-ie\overline{\psi}\slashed{A}\psi$ term in the Lagrangian. This term came fundamentally from our covariant gauge derivative $D_{\mu}$, and so are \textit{not} free to fix its form. This is a principle that carries over to all gauge theories: the gauge invariance fixes the kinds of interactions we can have. We will see in the next chapter that this gives some interesting results in non-Abelian theories. 
\er 
\chapter{Non-Abelian Gauge Theories}

\mybox{
    Just a reminder/disclaimer: This chapter was not part of the taught material, however I have worked from Dr. Iqbal's notes to fill it in here.
}

Finally let's discuss non-Abelian gauge theories. Being gauge theories, a lot of the tricks and tips we use here will be reminiscent (or indeed identical in some places!) of the previous chapter. However the non-Abelian nature of our gauge fields will obviously be important and so will give rise to some completely new things. It is these "additional things" that essentially result in the differences between QED (the electromagnetic force) and QCD (the strong force). 

\section{Non-Abelian Gauge Invariance}

As before, the first thing we have to do is ask the question of how things transform and how we must modify our action if we are to "promote" our global symmetry to a gauge one. So the first thing we need is our non-Abelian group to define our transformations. Of course any\footnote{Compact, simple.} one will do (if we want to stay completely general), however in order to have some relation to the standard model (which is where QCD lives), we shall keep SU($N$)\footnote{The gauge group of QCD is SU$(3)$.} in the back of our minds. 

\subsection{Recap On Group Theory}

In order to proceed (and to set notational convention) let's start by recalling some group theory. Recall that a Lie group, $G$, has an associated Lie algebra, $\mathfrak{g}$, and we can use basis vectors $T^a$, with $a=1,...,\dim\mathfrak{g}$, of the Lie algebra to \textit{generate} the Lie group. This is true on an abstract level; that is we do not need a representation in order to do this but it is a general statement about Lie groups and their algebras. To give a more concrete example, if we were working with $SU(2)$, for example, we would have $\dim\mathfrak{g}= 2^2 -1 = 3$ generators. At no point have we mentioned the Pauli matrices or any kind of representation. We shall stick to the notation of $T^a$ being the abstract generators and then introduce other notations as we go for specific representations. 

Now also recall that a Lie algebra comes equipped with a antisymmetric bilinear map, known as a Lie bracket, $[\cdot, \cdot ]$. We can use the Lie bracket in order to define the \textit{structure} constants, $f^{abc}$:\footnote{Note an $i$ is included here. This is different to the convention used on the group theory course. This is just the convention Dr. Iqbal is using and so I am sticking to it here.}
\be 
\label{eqn:StructureConstants}
    [T^a, T^b] = if^{abc}T^c,
\ee 
where a sum over $c$ is implied. It is also easy to see that $f^{abc}$ is completely antisymmetric in its indices.  
\bbox 
    Recall the Jacobi identity for the Lie bracket 
    \bse 
        \big[X,[Y,Z]\big] + \big[Z,[X,Y]\big] + \big[Y,[Z,X]\big] = 0.
    \ese 
    Use this along with \Cref{eqn:StructureConstants} to show that 
    \bse 
        f^{ade}f^{bcd} + f^{bde}f^{cad} + f^{cde}f^{abd} = 0.
    \ese 
\ebox  

Again this is all abstract, and we can get a better handle on it by giving a concrete example (using this as a chance to introduce some more notation). 

\bex 
    Let our Lie group be SU($2$). Then we can use the \textit{fundamental representation}
    \bse 
        T^a|_\text{fund.} =: t^a = \frac{1}{2}\sig^a
    \ese 
    with $a=1,2,3$ and $\sig^a$ being the Pauli matrices. We have used this to define our notational for the fundamental representation, little $t$. Our structure constants are then just the Levi-Civita tensor,\footnote{Perhaps I'm being a little bit overkill with subscripts etc here. This is just to try make sure the idea is clear.} 
    \bse 
        [t^a,t^b]_{SU(2)} = i\epsilon^{abc}t^c.
    \ese 
\eex 

\subsection{Building The Action}

Ok now let's talk about fields so we can start building our Lagrangian. As we are ultimatley leading to QCD, let's introduce a Fermion field $\psi$. We will let it transform in the fundamental representation\footnote{Note there is no reason a priori to assume this. The reason we do this is because we know that quarks carry 3 possible colours and so transform in the fundamental representation of SU(3).} of our SU($N$) group:
\bse 
    \psi(x) \mapsto \psi'(x) = U(x) \psi(x), \qand \overline{\psi}(x) \mapsto \overline{\psi}'(x) = \psi(x) U^{\dagger}(x),
\ese
where $U(x)\in SU(N)$ is our fundamental representation of the action of the gauge group. 

We can then try construct our Dirac action, 
\bse 
    S[\psi, \overline{\psi}] = \int d^4 x \, \overline{\psi}(i\slashed{\p} -m)\psi, 
\ese 
and check for gauge invariance. We already know this is not gauge invariant from our discussion above, namely because the derivative acts on $U(x)$. We already know how to fix this: we define a covariant derivative $D_{\mu}$ to counteract the bad terms so that 
\be 
\label{eqn:NonAbelianCovariantDerivativeTransformation}
    D_{\mu}\psi(x) \mapsto D_{\mu}'\psi'(x) =  U(x) D_{\mu}\psi(x).
\ee 
Now recall we did this for the Abelian case by introducing a gauge field $A_{\mu}$ and having it transform in a certain way. We want to do a similar thing, but now we need to be a bit more careful: we now have $\dim\mathfrak{g}$ different generators, and so our transformation matrix is given by 
\bse 
    U(x) = \exp\big( i \a^a(x) t^a\big),
\ese
where $\a^a\in \C$. We didn't have to worry about this before because we were dealing with a $U(1)$ gauge theory and this only has one generator (so we only had $\Lambda$ in our exponential). Now recalling that the generators come from the basis of the Lie algebra, it follows that they are linearly independent, and so we can only deal with them one at a time. We therefore have to introduce $\dim\mathfrak{g}$ different gauge fields in our covariant derivative. Put another way, when the derivative acts on $U(x)$ we will get $\dim\mathfrak{g}$ different terms by 
\bse 
    (\p_{\mu}\a^1) t^1 + (\p_{\mu}\a^2) t^2 + ... + (\p_{\mu}\a^{\dim\mathfrak{g}}) t^{\dim\mathfrak{g}},
\ese 
all of which are linearly independent. So we have to introduce a different $A_{\mu}$ for every term in order to remove all the behaviour. It is hopefully clear that the answer to this problem is to define our covariant derivative via 
\mybox{
    \be 
    \label{eqn:NonAbelianCovariantDerivative}
        D_{\mu} := \p_{\mu} - ig A^a_{\mu}(x)t^a,
    \ee 
}
\noindent where $g$ is the \textit{gauge-coupling}.\footnote{It plays the role $e$ played in QED.}

Ok now that we have our covariant derivative, we can ask the question about how the $A_{\mu}^a$s need to transform in order to give us a gauge invariant action. Well we have 
\bse 
    \begin{split}
        D_{\mu}\psi &\mapsto D_{\mu}' \psi'(x) \\
        & = \big(\p_{\mu} - ig A^{\prime a}_{\mu}(x) t^a\big)\big(U(x)\psi(x)\big) \\
        & = \big(\p_{\mu}U(x)\big)\psi(x) + U(x) \p_{\mu}\psi(x) - igA^{\prime a}_{\mu}(x) t^a U(x) \psi(x),
    \end{split}
\ese 
where we note we have \textit{not} commuted the $U(x)$ through the  $A^{\prime a}_{\mu}$ term; this is because we have a $t^a$ and $U(x)$ contains generators too, so we cannot simply commute them. If we then demand \Cref{eqn:NonAbelianCovariantDerivativeTransformation} holds, we can easily\footnote{Bonus exercise, show this.} show that we require
\mybox{
    \be 
    \label{eqn:NonAbelianGaugeFieldTransformation}
        A^{\prime a}_{\mu}(x)t^a = U(x) \bigg( A^a_{\mu}(x)t^a + \frac{i}{g}\p_{\mu}\bigg)U^{\dagger}(x).
    \ee 
}

\br 
    Firstly note that if we had defined \Cref{eqn:NonAbelianCovariantDerivative} with a plus sign instead of a minus sign everything would be the same expect that we would get a minus sign in \Cref{eqn:NonAbelianGaugeFieldTransformation}. Indeed you are completely free to put this minus sign in either place, and for the Abelian case (\Cref{eqn:GaugeInvariantDerivative,eqn:GaugeTransformationOnFields}) we did exactly this. This was not an inconsistency in these notes, but is one of the strange conventions that exist. That is it is standard to put the plus sign with $D_{\mu}$ for Abelian cases, while the plus sign goes with the $A_{\mu}'$ for non-Abelian cases.
\er 

\br 
    The $U^{\dagger}(x)$ appearing on the right-hand side of \Cref{eqn:NonAbelianGaugeFieldTransformation} comes strictly from the fact that we're using SU($N$). In general we put $U^{-1}(x)$ in its place. If you want to keep everything general, feel free to make this substitution, as very little of what follows depends on the SU($N$) specialisation.
\er 

\br 
    Recall (or look up) that the adjoing representation of a Lie group has the Lie algebra as its vector space and it acts as 
    \bse 
        U T^a U^{-1}.
    \ese 
    Keeping the above remark in mind, \Cref{eqn:NonAbelianGaugeFieldTransformation} looks \textit{almost} like the adjoint transformation. The pesky bit is that we also have the derivative term in there. As we will see in a moment, this derivative term essentially causes our gauge fields to rotate in the vector space.
\er 

Ok great, we have a formula for the transformation of our gauge fields and a gauge invariant action! The next instructive thing to do is to expand \Cref{eqn:NonAbelianGaugeFieldTransformation} and obtain the infinitesimal result. We shall leave this as an exercise. 

\bbox 
    Expanding to $\cO(\a)$ show that \Cref{eqn:NonAbelianGaugeFieldTransformation} is equivalent to 
    \bse
        A^{\prime a}_{\mu}(x) t^a = A^{\prime a}_{\mu}(x) t^a - \a^b(x) f^{bac} t^c A^a_{\mu}(x) + \frac{1}{g}\p_{\mu}\a^a(x) t^a.
    \ese
\ebox 

We can do some index gymnastics on the result of the above exercise to obtain 
\mybox{
    \be 
    \label{eqn:NonAbelianGaugeFieldInfitesimalTransformation}
        A^{\prime a}_{\mu}(x) = f^{abc} A^b_{\mu}(x) \a^c(x) + \frac{1}{g}\p_{\mu}\a^a. 
    \ee 
}

Why was this instructive? Well it allows us to make a few important points:
\ben[label=(\roman*)]
    \item The last term in \Cref{eqn:NonAbelianGaugeFieldInfitesimalTransformation} is identical to our Abeliean gauge field transformation, \Cref{eqn:GaugeTransformationOnFields}. 
    \item The first term is something new, though. We see that it comes with the structure constants (which necessarily vanish for the Abelian case), and so we see it stems fundamentally from the non-Abelian nature of our gauge group. In order to understand what this term means we need to recall that our index $a$ runs over the Lie algebra indices, which are the indices of a vector space. We can therefore think about our $A_{\mu}^a$s as being vectors in this space and pointing in some direction. What this first term is telling us is that non-Abelian-ness of the group is causing our vector to rotate. 
    \item Although we have used the fundamental representation in the above discussion, it is hopefully clear that essentially nothing relied on it specifically. Indeed it is likely that your action will contain several different fields, each of which might transform in different representations. For example in the standard model quarks transform in the fundamental representation of SU(3) while leptons transform in the trivial representation.\footnote{This is the theoretical statement that quarks carry colour charge while leptons do not. See a course on SM for more information.} 
    
    The important point is that the transformation of our gauge fields do \textit{not} depend on the representation of the fields they couple too, as otherwise we would have it transforming in multiple different ways, and this would change its charge structure which would be unreasonable physically. 
\een 

\bnn 
    Before moving on to constructing a full non-Abelian (or \textit{Yang-Mills} action), let's make a comment on some common notation. We will sometimes use 
    \be 
    \label{eqn:AmuMatrixValued}
        A_{\mu} := A_{\mu}^a t^a
    \ee 
    to denote the \textit{matrix} version of the gauge field.\footnote{More technically it is a Lie-algebra-valued one form.} This is just a short hand that saves us having to write a bunch of indices everywhere.  
\enn 

\section{The Yang-Mills Field Strength Action}

We are now in a position where we know how to introduce a covariant derivative in order to compensate the for the non-gauge invariant parts in derivative terms. We now want to use this in order to write down the most general gauge invariant action. 

As with the Abelian case, we want to introduce the field strength in order to make our gauge fields propagating degrees of freedom. Completely analogously to the Abelian case we obtain this by noting that 
\bse 
    [D_{\mu},D_{\nu}]\psi \mapsto U(x) [D_{\mu},D_{\nu}]\psi,
\ese 
and then we keep the definition \Cref{eqn:Fmunu}:
\be 
\label{eqn:FmunuNonAbeianOne}
    ig F_{\mu\nu} := [D_{\mu},D_{\nu}],
\ee 
where we have replaced $e\to g$ for consistency. We can now use \Cref{eqn:NonAbelianCovariantDerivative} to get an expression for the field strength in terms of the gauge fields $A_{\mu}$. Again we should remember that we really have $\dim\mathfrak{g}$ different gauge fields, labelled by the values of $a$, and so we should use \Cref{eqn:AmuMatrixValued} to write \Cref{eqn:FmunuNonAbeianOne} as
\bse 
    ig F_{\mu\nu}^a t^a = [D_{\mu}, D_{\nu}].
\ese 
From here we can show that 
\mybox{
    \be 
    \label{eqn:FmunuNonAbelian}
        F_{\mu\nu}^a = \big(\p_{\mu}A_{\nu}^a - \p_{\nu}A_{\mu}^a\big) + gf^{abc} A_{\mu}^b A_{\nu}^c,
    \ee 
}
\noindent where a sum over $b$ and $c$ is implied. 

\bbox 
    Prove \Cref{eqn:FmunuNonAbelian}. \textit{Hint: Note the first term is just the Abelian term, so it's just the second one you need to derive.}
\ebox 

Now there is an important difference between the Abelian field strength and the non-Abelian one: recall that $F_{\mu\nu}^{\text{Abelian}}$ was \textit{itself} gauge invariant, and so it was trivially true that the product term $F_{\mu\nu}F^{\mu\nu}$ was gauge invariant. The same isn't true for the non-Abelian case. We can show this by recalling that it transforms in the adjoint\footnote{Again we are assuming $SU(N)$ to get the dagger.}
\bse 
    F_{\mu\nu}^{\prime} = F_{\mu\nu}^{\prime a}t^a = U(x) F_{\mu\nu}^a t^a U^{\dagger}(x).
\ese 
Plugging through the algebra\footnote{Additional exercise, do this.} we can show that 
\bse 
    F_{\mu\nu}^{\prime a} = F_{\mu\nu}^a - f^{abc}\a^b F^c_{\mu\nu} + \cO(\a^2),
\ese 
where as always the sum over $b$ and $c$ is assumed. Note this result at least makes \textit{some} sort of sense as the non-gauge invariant bit is given by the structure constants $f^{abc}$, which vanish is we then consider the Abelian case. 

However not all hope is lost as we can still make a term that is gauge invariant in a product by taking a trace over our gauge indices $a$. Explicitly, 
\bse 
    \begin{split}
        \tr[F_{\mu\nu}^{\prime}F^{\prime \mu\nu}] & = \tr[U(x) F_{\mu\nu} U^{\dagger}(x) U(x) F^{\mu\nu}U^{\dagger}(x)] \\
        & = \tr[ U(x) F_{\mu\nu} F^{\mu\nu} U^{\dagger}(x)] \\
        & = \tr[F_{\mu\nu}F^{\mu\nu} U^{\dagger}(x)U(x)] \\
        & = \tr[F_{\mu\nu}F^{\mu\nu}],
    \end{split}
\ese 
where we have used the cyclicity of the trace. We can therefore write this part of the action as 
\mybox{
    \be 
    \label{eqn:SYM}
        S_{YM}[A] = \int d^4 x \bigg(-\frac{1}{2}\tr[F_{\mu\nu}F^{\mu\nu}]\bigg)
    \ee 
}

We have done this in terms of the matrix valued expressions, but we can just as easily do it in terms of the decomposition in terms of generators by introducing the following fact. 
\bcl 
    It is always possible to pick a basis for $\mathfrak{su}(N)$ such that the generators obey 
    \be 
    \label{eqn:TraceOfGenerators}
        \tr[t^at^b] = \frac{1}{2}\del^{ab}.
    \ee 
\ecl 

\bbox 
    Use \Cref{eqn:TraceOfGenerators} to show that \Cref{eqn:SYM} can be written in the, perhaps more familiar, form
    \bse
        S_{YM}[A] = \int d^4 x\bigg(-\frac{1}{4} F_{\mu\nu}^a F^{a\mu\nu}\bigg). 
    \ese  
    \textit{Hint: Note that the trace is over the SU($N$) space, so anything that doesn't have a $a$ index is just seen as a number from the trace's perspective.}
\ebox

This result is important and so we shall put it in one of our nice boxes. 
\mybox{
    \be
    \label{eqn:SYMWitha}
        S_{YM}[A] = \int d^4 x\bigg(-\frac{1}{4} F_{\mu\nu}^a F^{a\mu\nu}\bigg). 
    \ee 
}

\br 
    Perhaps this remark should have been sooner, but note that we have always put the $a$ index as a superscript, even when it came to doing sums. Although it might seem `more natural' (given our GR training) to write something like 
    \bse 
        F_{\mu\nu}^a F^{\mu\nu}_a.
    \ese 
    The problem with this is we don't actually have anything to `lower' indices in our SU($N$) space, whereas with Lorentz indices we can use the metric. To be a bit more clear we \textit{define}
    \bse 
        A_{\mu} := A^{\nu}g_{\mu\nu} 
    \ese  
    where $A$ is some abstract vector, however we can't make any such definition for the lowering of our $a$ indices.\footnote{To the differential geometry familiar, essentially what we're saying is that we can think of the metric as a linear mapping between vector fields and 1-forms, whereas no such map exists for our $a$ indices.}
\er 

\subsection{Some Classical Aspects}

Before moving on to study the quantum theory let's look at some of the \textit{classical} aspects of the theory. This subsection will hopefully help highlight the subtleties that often go on during the quantisation process in QFT, by which we mean that we often just look at the classical theory (i.e. the Lagrangian) and start doing quantum physics (i.e. drawing Feynman diagrams with loops etc). This clearly normally works pretty well, however let's now outline an example of what can go wrong in this na\"{i}ve extension. 

We start by finding the classical field equations that accompany our Yang-Mills action, which we obtain by varying the action 
\bse 
    \begin{split}
        \del S_{YM}[A] & = \int d^4 x \bigg(-\frac{1}{2}F^{a\mu\nu} \del F_{\mu\nu}^a \bigg) \\
        & = \int d^4x \bigg(-\frac{1}{2}F^{a\mu\nu}\big[ \p_{\mu}\del A^a_{\nu}  - \p_{\nu}\del A^a_{\mu} + g f^{abc} \big( \del A_{\mu}^b A_{\nu}^c + A_{\mu}^b \del A_{\nu}^c\big)\big]\bigg) \\
        %& = \int d^4 x \bigg(-\frac{1}{2}\bigg)\big[ -\p_{\mu}F^{a\mu\nu} \del A_{\nu}^a + \p_{\nu} F^{a\mu\nu} \del A_{\mu}^a - gf^{abc}F^{c\mu\nu} \big( \del A^b_{\mu} A^a_{\nu} + A_{\mu}^b \del A^a_{\nu} \big)\big] \\ 
        % Can't quite get the above line to work out correctly.
        & = \int d^4 x \big( \p_{\mu}F^{a\mu\nu} + gf^{abc}A^a_{\mu}F^{c\mu\nu}\big) \del A^b_{\nu},
    \end{split}
\ese
where several tricks have been played (e.g. integrate by parts and use the antisymmetric nature of $F^{\mu\nu}$). We set this to zero for all variations and so conclude 
\bse 
    \p_{\mu}F^{a\mu\nu} + g f^{abc}A^a_{\mu}F^{c\mu\nu} = 0.
\ese 
Then if we recall that $F_{\mu\nu}$ transforms in the adjoint we see this is nothing but\footnote{This comes from the general rule $D_{\mu}A^{\mu} = \p_{\mu}\Phi - g T^aA^a_{\mu}\Phi$ for some general field $\Phi$ which transforms in the the $T^a$ representation.}
\mybox{
    \be 
    \label{eqn:YangMillsEquation}
        D_{\mu}F^{a\mu\nu} = 0.
    \ee 
}
\noindent This is known as the classical \textit{Yang-Mills equation}, and, despite the simple looking notation, it is highly non-linear, as can easily be seen from the previous line ($A_{\mu}^aF^{c\mu\nu}$ contains products between the $A$s).

We now consider coupling our gauge fields $A_{\mu}^a$ to matter. In order to make connection with the standard model, and in particular the $SU(3)$ sector, we consider couplings to Fermions, $\psi$. However in order to stay somewhat general we do not specify the type of Fermion just yet. This last generality statement is the same as saying we leave the representation arbitrary and so our covariant derivative acts as 
\bse 
    D_{\mu}\psi = \p_{\mu}\psi - ig A^a_{\mu} T^a \psi,
\ese
where $T^a$ are our the relevant generators. We then have the action 
\bse 
    S[\psi,\overline{\psi},A] = \int d^4 x \, \bigg(-\frac{1}{2}\tr[F_{\mu\nu}F^{\mu\nu}] + \overline{\psi}\big(i\slashed{D} - m\big)\psi\bigg),
\ese 
which we can vary in order to obtain the following equations of motion 
\be 
\label{eqn:YangMillsWithField}
    D_{\mu}F^{a\mu\nu} = -g\big(\overline{\psi} \g^{\nu} T^a \psi\big) \qand \big(i\slashed{D}-m)\psi = 0.
\ee 
We can compare the first to \Cref{eqn:YangMillsEquation} and notice that the right-hand side would just be our classical Noether current if we had a \textit{global} symmetry.

We now specialise to quarks and replace $\psi$ with $q_i$, where the index $i$ labels the flavour of quark,\footnote{That is: up, down, charm, strange top or bottom.} and take $N=3$. Our action then becomes 
\bse 
    S[\psi,\overline{\psi},A] = \int d^4 x \, \bigg(-\frac{1}{2}\tr[F_{\mu\nu}F^{\mu\nu}] + \sum_i \overline{q}_i\big(i\slashed{D} - m_i\big)q_i\bigg),
\ese 
where we have allowed the masses for each quark to be different (as indeed they are). As we have mentioned a few times above, this is just the action for QCD, and here our $3^3-1=8$ gauge fields are the \textit{gluons}. It is an experimental fact that quarks carry one of three colours --- conventionally named red, green and blue --- from which it follows that they transform in the fundamental representation, i.e. $T^a=t^a$. 

Now comes the subtle points in the classical to quantum transition. We wont go into too much detail,\footnote{Namely because Dr. Iqbal doesn't and currently I want to get these notes finished so I can move on to other modules.} but the idea is as follows: imagine in the classical theory forcing a quark to sit still in at some point $x$. This is just the statement that we want to the spatial component of the current in \Cref{eqn:YangMillsWithField} to be a delta function. If we do this we find, to lowest order in $g$, that 
\bse 
    D_{\mu}F^{a\mu\nu} = g (t^a)_{11} \del^{(3)}(\vec{x}).
\ese
As we said before, the left-hand side of this equation is actually highly non-linear and messy, but apparently it turns out that you can show that these non-linear terms don't contribute in any way, and so we are left with a Maxwell-type expression 
\bse 
    \p_{\mu} F^{a\mu\nu} = g(t^a)_{11}\del^{(3)}(\vec{x}),
\ese 
which we know from our undergrad days has a inverse-square solution, namely
\bse 
    F^{art} \sim \frac{1}{r^2}.
\ese

This is where our problems arise. Physically this is telling us that if we increase the distance between two quarks the $SU(3)$-type electric field (or \textit{colour} field) dies off. This therefore means that we could essentially separate the two quarks and study them (essentially) completely independently. This is \textit{not} what we see in reality: confinement tells us that we cannot isolate a quark, never mind asking the question of "does it give off a $1/r^2$ colour field?" Another problem with the results we draw from here is that gluons are \textit{not} massless. 

So we see that, unlike for QED where the classical theory gave us almost the exact QFT, the classical QCD results are wildly wrong and so we should be careful in the future before na\"{i}vely powering through with classical results. 

\br 
    It is important to note that the massiveness of the gluons has \textit{nothing} to do with spontaneous symmetry breaking and the Higgs mechanism. It is for this reason we really shouldn't say that the Higgs gives everything mass. 
\er 

\section{Quantisation Non-Abelian Gauge Theories}

Now that we have seen the subtleties of the classical theory, let's look towards the quantisation of our non-Abelian theory. Of course if we keep our Fermion content we will also obtain the quantisation of the Fermions, but we have already seen this, so we can forget about it. Note we also drop the calculation of the gauge fields to the Fermions, but we are hopefully now trained enough to "see it" by considering the action. What we really want to focus on is the pure Yang Mills action \Cref{eqn:SYMWitha}. Before proceeding let's just comment on the things we expect to get from this term
\ben[label=(\roman*)]
    \item This term contains $\p A \p A$ terms, and so we expect to get the propagator for our gauge fields, 
    \item Recalling that $F_{\mu\nu}^a$ contains a nonlinear term of the form $gf^{abc} A^b_{\mu}A^c_{\nu}$, we expect to get couplings between the gauge fields themselves. In particular we expect to get both 3-point couplings (from $fAA\p A$) and 4-point couplings (from $fAA fAA$). Note as with the Abelian case, this interacting behaviour is completely fixed by gauge invariance.
\een 
We highlighted this now because we see that this expect additional coupling behaviour is completely new. The other thing we want to point out is that we only have one parameter in our theory, $g$, and so if our theory is going to renormalisable it better be the case that fixing $g$ fixes \textit{all} divergent behaviour at that level simultaneously. We do not discuss this in much detail here but more can be found via the QCD course. 

\subsection{Self Interactions}

First let's look at the gauge field self interactions. Our action is
\bse 
    S_{YM}[A] = \int d^4 x \bigg(-\frac{1}{4}\big[ \p_{\mu}A_{\nu}^a - \p_{\nu}A_{\mu}^a + g f^{abc}A_{\mu}^b A_{\nu}^c\big] \big[ \p^{\mu}A^{a\nu} - \p^{\nu}A^{a\mu} + g f^{ade}A^{d\mu} A^{e\nu}\big]\bigg),
\ese 
which we can expand out to see the form of the self-interaction terms. That is collect all the terms that have 3 $A$s and all the terms that have 4 $A$s and look at the prefactors. We do not do the expansion here but just not a couple things 
\ben[label=(\roman*)]
    \item The 3-point interactions will all come with derivatives. We are yet to see how to translate these into Feynman rules, but it is simpler then we might think: simply go to momentum space with the usual prescription $\p_{\mu} \to -ip_{\mu}$. We therefore expect the momenta of our gauge fields to enter the coupling strength. Note, we obviously have to be careful about the sign of $p_{\mu}$, and in the diagrams that follow we use the convention of having \textit{all} momenta flowing \textit{into} the vertex. 
    \item The 4-point interaction will not come with derivatives and so we do not expect any momentum dependence.
    \item The ordering of the gauge fields matters. That is, for example,
    \bse 
        A^b_{\mu}A^c_{\nu} A^{d\mu}A^{e\nu} \neq A^b_{\mu} A^{d\mu} A^c_{\nu} A^{e\nu}. 
    \ese 
    This will translate into us keeping track of the Lie algebra indices (the $a,b$ etc) on the Feynman diagrams. 
\een 

Very good, so now in order to get the vertices we have to do what we just outlined above (expand out and replace derivatives etc). We do not do this calculation here as it is quite tedious,\footnote{For the interested reader, a derivation was set on the first tutorial sheet for QCD by Dr. Huss, so if you take that course, see that sheet.} but simply give the result. Let's give the results first then see how we can remember them.

\bnn 
    In order to draw parallels with QCD we depict our non-Abelian gauge fields with a "springy" line, as this is standard for gluons. It is worth noting, though, that this is a gluon specific thing: we have already used a wavey line for the photon and it is standard to use either dashed or wavey lines for the $W^{\pm}$ and $Z$ gauge bosons. Of course none of this matters as it is just a convention, but this is just included for clarity.
\enn 

\mybox{
    \begin{center}
        \btik 
            \begin{scope}
                \gluon (0,1.5) -- (0,0);
                \node at (0,1.8) {$a,\mu$};
                \draw[thick, ->] (0.3,1.1) -- (0.3,0.4) node [midway, right] {$k$};
                \gluon[rotate around={120:(0,0)}] (0,1.5) -- (0,0);
                \node at (1.7,-0.8) {$b,\nu$};
                \draw[thick, ->, rotate around={120:(0,0)}] (0.3,1.1) -- (0.3,0.4);
                \node at (-0.85,0.05) {$q$};
                \gluon[rotate around={-120:(0,0)}] (0,1.5) -- (0,0);
                \node at (-1.7,-0.8) {$c,\rho$};
                \draw[thick, ->, rotate around={-120:(0,0)}] (-0.3,1.1) -- (-0.3,0.4);
                \node at (0.85,0.05) {$p$};
                \node[right] at (2,0.5) {$ = -gf^{abc}\big[\eta^{\mu\nu}(k-p)^{\rho} + \eta^{\nu\rho}(p-q)^{\mu} +\eta^{\rho\mu}(q-k)^{\nu}\big] $};
            \end{scope}
            \begin{scope}[yshift=-3cm]
                \gluon[rotate around={45:(0,0)}] (0,1.5) -- (0,0);
                \node at (-1.25,1.3) {$a,\mu$};
                \gluon[rotate around={-45:(0,0)}] (0,1.5) -- (0,0);
                \node at (1.25,1.3) {$b,\nu$};
                \gluon[rotate around={135:(0,0)}] (0,1.5) -- (0,0);
                \node at (1.25,-1.3) {$c,\rho$};
                \gluon[rotate around={-135:(0,0)}] (0,1.5) -- (0,0);
                \node at (-1.25,-1.3) {$d,\sig$};
                \node[right] at (2,0) {$=$};
                \node[right] at (2.5,1) {$-ig^2 \big[ f^{abe}f^{cde} (\eta^{\mu\rho}\eta^{\nu\sig} - \eta^{\mu\sig}\eta^{\nu\rho})$};
                \node[right] at (3.3,0.25) {$+ f^{ace}f^{bde} (\eta^{\mu\nu}\eta^{\rho\sig} - \eta^{\mu\sig}\eta^{\nu\rho})$};
                \node[right] at (3.3,-0.5) {$+ f^{ade}f^{bce} (\eta^{\mu\nu}\eta^{\rho\sig} - \eta^{\mu\rho}\eta^{\nu\sig})  \big]$};
            \end{scope}
        \etik 
    \end{center}
}

Ok so how do we go about remembering (at least the rough form of) these result? First let's look at the 3-point vertex. 

\subsubsection{3-Point Vertex}

Firstly we recall that the three gauge field interaction comes from the product on one derivative, $\p A$, and one structure constant, $g f AA$, term. The former doesn't have any factors of $g$ and so we expect our 3-point interaction to be of order $g$. We also include a factor of $-1$, because these always crop up at vertices.

Next, as we said in (iii) above, we have to keep track of the Lie algebra indices as well as the Lorentz ones. We have 3 gluons and so we expect each term in the expansion to have 3 Lie algebra indices and 3 Lorentz ones. Putting this together with the fact that we know the momenta enter the rule, as per (i) above, we see that the only reasonable combination are terms of the form 
\bse 
    -g f^{\cdot \cdot \cdot} \eta^{\bullet \bullet} (f(p,k,q))^{\bullet},
\ese 
where the $\cdot$s need to be filled by Lie algebra indices, the $\bullet$s by Lorentz ones and $f(p,q,k)$ is some \textit{linear} combination of the momenta.

Now we remember that ordering matters, and so before we can proceed we need to draw the vertex and distribute indices and momenta. We do just that, and for ease of explanation assume it is done as per the diagram above. Now comes the memory part; the idea is to start at the top and work clockwise when distributing indices and then take cyclic permutations for the Lorentz ones. That is each term comes with $f^{abc}$ but we get three different terms corresponding to the 3 cyclic permutations of $\mu,\nu$ and $\rho$. Finally the only other thing we have to remember is that our momentum dependence comes with a minus sign and again we go clockwise. Putting this all together, we expect to get, for example, the term $-gf^{abc}\eta^{\mu\nu}(k-p)^{\rho}$, which agrees exactly with the first term above.

\br 
    Note that we actually could go \textit{anti}clockwise with the same prescriptions and get the same result. This follows from the fact that we would need pick up a minus sign from swapping two $f^{abc}$ indices and another minus sign for swapping the momentum parts, i.e. $(k-p)^{\rho} = -(p-k)^{\rho}$. 
\er 

\bbox 
    Without looking at the answer, try and use the above memory steps to arrive at the 3-point interaction vertex. 
\ebox   

\subsubsection{4-Point Vertex}

Ok what about the $4$-point vertex. This is a lot uglier and I don't feel like I'm going to be able to give an explanation that makes memory any easier than anything any reader could come up with. 

\subsection{Propagator}

We now need to study the propagator for the non-Abelian gauge fields. As these are gauge fields, we still run into the same infinite contribution from the path integral over all, unphysical, gauge paths and so we will once again need to use the Fadeev-Poppov procedure. However, as we will shortly see, things become a bit more complicated when we have a non-Abelian gauge field. 

Again we work with the gauge fixing condition 
\bse 
	G^a_{\omega}(A) := \p_{\mu} A^{a \mu}(x) - \omega^a(x),
\ese 
where the $a$ reminds us that this is actually a gauge fixing condition for \textit{each} gluon. We then recall that our general gauge transformation (in the fundamental) is given by 
\bse 
	U(x) = \exp\big( i \a^a(x) t^a\big),
\ese
and that our gauge fields try to transform in the adjoint. That is, if we denote the transformed field via $ (A^{\a})^a_{\mu} $, we have 
\bse 
	(A^{\a})^a_{\mu} t^a = U(x) \bigg( A^a_{\mu} t^a + \frac{i}{g} \p_{\mu} \bigg) U^{\dagger}(x),
\ese 
which, by direct calculation, gives us 
\bse 
	(A^{\a})^a_{\mu} = A^a_{\mu} + \frac{1}{g}D_{\mu} \a^a.
\ese 

Now just as before we want to introduce a functional delta function in order to cleverly insert $1$ into our path integral. Again we use 
\bse 
	1 = \int [\pD \a ] \del\big[G_{\omega}(A^{\a}) \big] \det\bigg( \frac{\del G_{\omega}(A^{\a}) }{\del \a} \bigg)
\ese 
so that our path integral becomes 
\bse 
	Z = \int [\pD A] [\pD \a] \del\big[G_{\omega}(A^{\a}) \big] \det\bigg( \frac{\del G_{\omega}(A^{\a}) }{\del \a} \bigg) e^{iS[A]}.
\ese 

This is all \textit{exactly} the same as Abelian case, however we now notice a very important difference, which stems from the determinant term. Using the above expressions, we have (dropping the $a$ index for convenience)
\bse 
	\begin{split}
		\frac{\del G_{\omega}(A^{\a}) }{\del \a} & = \frac{\del}{\del \a} \big( \p_{\mu} (A^{\a})^{\mu} - \omega\big) \\
		& = \frac{\del}{\del \a}\bigg( \p_{\mu} \bigg[  A^{\mu} + \frac{1}{g}D^{\mu} \a \bigg] - \omega\bigg) \\
		& = \frac{1}{g} \p_{\mu} D^{\mu}.
	\end{split}
\ese 
The reason this is a problem is because our covariant derivative contains our gauge field $A^{\mu}$! This means that we \textit{cannot} simply strip the determinant factor outside the path integral as we did for the Abelian case. We shall return to this problem shortly, but first we note that this is going to have absolutely no effect on the propagator of the gauge field (which comes from the inverse of the quadratic term in the action, remember), and so following the same ideas as before we introduce a Gaussian weight etc to arrive at 
\bse 
	Z = \int [\pD A] \det\bigg( \frac{1}{g} \p_{\mu} D^{\mu} \bigg) \exp\bigg( iS[A] -i \int d^4 x \frac{1}{2\xi} \big(\p_{\mu} A^{\mu}\big) \bigg),
\ese 
from which we can read off the propagator. We have to remember, though, that our $A^{\mu}$ is a Lie-algebra valued expression and we should decompose it into the $\dim\mathfrak{g}$ different gauge fields $A_{\mu}^a$. In terms of the propagator this just corresponds to including a $\del^{ab}$ in the answer. This is just telling us that the colour (for the example of $SU(3)$) of the gluon doesn't change as it propagates. Explicitly we have the momentum space result
\mybox{
    \be 
        D_{\mu\nu}^{ab}(p) = - \frac{i}{p^2} \bigg(\eta_{\mu\nu} - (1-\xi) \frac{p_{\mu}p_{\nu}}{p^2}\bigg) \del^{ab}.
    \ee 
}
\noindent Here the gauge choice $\xi=1$ is referred to as \textit{Feynman-'t Hooft gauge}.

\subsection{Ghosts}

Now we need to work out what this pesky determinant term is doing. Clearly as it stands this determinant is ugly and ruins our ability to understand the problem. However, we now remember that a Gaussian integral over Grassman variables results in a expression with the determinant in the numerator and so we can \textit{interpret} this as such. That is we replace 
\be 
\label{eqn:DeterminantToGhosts}
	\det\big( \p_{\mu} D^{\mu} \big) = \int [\pD c][\pD \overline{c}] \exp \bigg( i  \int d^4 x \, \overline{c} \big(-\p_{\mu}D^{\mu} \big) c \bigg),
\ee
where $c$ and $\overline{c}$ are Grassman variables. They appear quadratically in our partition function, and so we can think of  the $c/\overline{c}$ as propagating degrees of freedom. That is we think of them as quantum fields that are Lorentz scalars (i.e. they transform in the trivial representation of the Lorentz group) but they transform in the adjoint of the gauge group. We call such fields \textit{Fadeev-Poppov ghosts}, and they play a huge role in non-Abelian gauge QFTs. Before moving on let’s make a few comments:
\ben[label=(\roman*)]
	\item Although these fields are anticommuting (they are Grassman variables), they are \textit{not} spinors. There are clearly many reasons for this but perhaps the easiest argument to make is they transform in the trivial representation of the Lorentz group, which spinors, by definition, do not. 
	\item It is very important to note that this is merely a trick we use to make our expression easier to understand and manipulate. This is important as it tells us that the Fadeev-Poppov ghosts are not physical fields and so we should never expect to get an external ghost field. However they can contribute to interior of diagrams, and as we shall briefly describe in a moment actually give nice intuitive results this way.
	\item As the ghost fields transform in the adjoint, they couple to gluons and so they have an actual qualitative effect. For example, the vacuum polarisation diagram for the gluon receives a contribution from a ghost loop. 
\een 

Ok let's find the propagator for the ghosts and their coupling to the gluons. We can rewrite the exponential in \Cref{eqn:DeterminantToGhosts} as a contribution to the action in the form
\bse 
    S_{\text{ghost}}[c,\overline{c},A] = \int d^4x \, \overline{c}^a \big(-\p^2 \del^{ac} - g \p^{\mu} f^{abc} A_{\mu}^b \big) c^c
\ese 
from which we can read off the propagator and interaction vertex. We summarise these in terms of diagrams now.
\mybox{
    \begin{center}
        \btik 
            \begin{scope}
                \midarrow[dashed] (-1,0) -- (1,0);
                \node at (-1.2,0) {$a$};
                \node at (1.2,0) {$b$};
                \node[right] at (2.5,0) {\Large{$= \frac{i}{p^2}\del^{ab}$}};
            \end{scope}
            \begin{scope}[yshift = -3cm]
                \gluon (0,1.5) -- (0,0);
                \node at (0,1.8) {$b,\mu$};
                \midarrow[dashed, rotate around={120:(0,0)}] (0,0) --  (0,1.5);
                \node at (1.5,-0.8) {$c$};
                \draw[thick, ->, rotate around={120:(0,0)}] (0.3,1.1) -- (0.3,0.4);
                \node at (-0.85,0.05) {$p$};
                \midarrow[dashed, rotate around={-120:(0,0)}] (0,0) -- (0,1.5);
                \node at (-1.6,-0.8) {$a$};
                \node[right] at (2.5,0.5) {\Large{$ = gf^{abc} p^{\mu} $}};
            \end{scope}
        \etik 
    \end{center}
}

\subsection{Wrapping Up}

We're now done! Well we need to put the coupling to Fermions back in, but we have already studied this and the ghosts don't interact with them (as they came purely from \Cref{eqn:DeterminantToGhosts}). Our full action is then 
\bse 
    S[\psi,\overline{\psi},A,c,\overline{c}] = S_{YM}[A] + S_{\text{quark}}[\psi,\overline{\psi},A] + S_{\text{ghost}}[c,\overline{c},A].
\ese 
From here we now have a full set of Feynman rules in order to start doing the quantum theory. For example if we make out gauge group SU($3$) we are now (theoretically) fully capable of calculating the contributions from QCD Feynman diagrams. We could then go on to look at the renormalisability of QCD and plug through that. We do not present that here (it is discussed in the QCD course), but just make a couple comments:
\ben[label=(\roman*)]
    \item QCD is in fact renormalisable. 
    \item From dimension counting arguments we can show\footnote{See QCD notes.} that a diagram corresponding to 2-gluons to 2-ghosts\footnote{Forgetting for the minute that we can't have external ghosts} processes would be superficially logarithmically divergent (i.e. they have SDOD $= 0$). However it turns out that these diagrams are all indeed convergent. This is very important for the renormalisability of the theory as our action doesn't contain any $2$-gluon-$2$-ghost terms, and so we couldn't absorb the infinite contribution. 
    \item As we mentioned before, the \textit{only} free parameter we have in our action is $g$, whereas we have multiple different divergent diagrams (e.g. we can have both gluon and Fermion loops). It turns out that it doesn't matter which diagram we choose to use to renormalise $g$, it will simultaneously make all the diagrams convergent. This is a really non-trivial result and will be discussed in more detail in the QCD notes. 
    \item Finally, we would be fair to ask the question "what are the Fadeev-Poppov ghosts of?" The answer to this is what we said we would return to in (ii) above. Recall that a general gauge field\footnote{In 4-dimensions, otherwise just use $d$.} has $4$ components but only 2 physical polarisations. Put another way, 2 of the polarisations are so-called \textit{spurious}. Physical things only make sense in external states (this is why stuff can be off shell as a propagator) and so there doesn't seem any reason why these two spurious polarisations couldn't propagate around a loop. However we do not want them to contribute to anything, as they are unphysical, and so we need some way to counteract them. This is where the ghosts come in: as ghosts are anticommuting fields the idea is that we compensate for our spurious loop with a ghost loop, with the latter picking up a minus sign for the same reason we get a minus sign for closed Fermion loops. These two contributions then exactly cancel and so we get the result we want. 
\een 

\br 
    Dr. Iqbal then has a discussion at the end of his notes entitled "Qualitative discussion: Non-Abelian gauge theory at long distances". It is mainly a discussion of confinement, but, as the title says, it is qualitative and all I could really do is just copy it across. I see little point in doing this, and so just end these notes here. If I get time later I shall read up on this stuff and try fill this in here more qualitatively. 
\er 

% -------------------------------------------------------------------
% Bibliography/Further Readings
% -------------------------------------------------------------------

\chapter*{Useful Texts \& Further Readings}

\section*{Textbooks}
\begin{itemize}
    \item M. E. Peskin and D. V. Schroeder, “\textit{An Introduction to QFT}” (Addison Wesley, 1995).
    \item M. Srednicki, "\textit{Quantum field theory}" (Cambridge University Press, 2007).
    \item A. Zee, “\textit{Quantum Field Theory in a Nutshell}” (Princeton University Press, 2010).
    \item R.J.Rivers “\textit{Path Integral methods in QFT}” (CUP 1987)
    \item R.P. Feynman and A.R. Hibbs, “\textit{Quantum Mechanics and Path Integrals}” (Dover Books on Physics, 2010)
\end{itemize}

\section*{Other Notes}
\begin{itemize}
    \item J. F. Donoghue, “\textit{Introduction to the effective field theory description of gravity}” in Advanced School on Effective Theories Almunecar, Spain, June 25-July 1, 1995. 1995. \href{https://arxiv.org/pdf/gr-qc/9512024.pdf}{Available on arXiv}.
    \item Prof. Timo Weigand "\href{https://www.thphys.uni-heidelberg.de/~weigand/QFT2-14/SkriptQFT2.pdf}{Quantum Field Theory I + II}", Institute for Theoretical Physics, Heidelberg University. 
\end{itemize}


%\bibliographystyle{agsm} 
%\bibliography{mybibliography} 
%\printbibliography[heading=bibintoc]


% -------------------------------------------------------------------
% Appendices
% -------------------------------------------------------------------

%\begin{appendices}
%\chapter{An Appendix of Some Kind}

\lipsum[1-3]  % Replace with your text

%\end{appendices}

\end{document}
